\documentclass[11pt]{article}

\usepackage[a4paper,margin=2.35cm]{geometry}
\usepackage{amsmath,amssymb,bm,mathtools,slashed}
\usepackage{booktabs,array,longtable}
\usepackage{graphicx}
\usepackage{enumitem}
\usepackage{xcolor}
\usepackage{microtype}
\usepackage[
  colorlinks=true,
  linkcolor=blue!55!black,
  citecolor=blue!55!black,
  urlcolor=blue!65!black
]{hyperref}

\numberwithin{equation}{section}
\allowdisplaybreaks[2]
\setlength{\jot}{5pt}
\setlist[itemize]{leftmargin=1.7em,itemsep=0.25em,topsep=0.35em}
\setlist[enumerate]{leftmargin=1.8em,itemsep=0.25em,topsep=0.35em}
\renewcommand{\arraystretch}{1.15}
\sloppy

\newcommand{\dd}{\mathrm{d}}
\newcommand{\ii}{\mathrm{i}}
\newcommand{\MSbar}{\overline{\mathrm{MS}}}
\newcommand{\cO}{\mathcal{O}}
\newcommand{\rflow}{r_{\mathrm{flow}}}
\newcommand{\tf}{t_{\mathrm f}}
\newcommand{\qtilde}{\widetilde q}
\newcommand{\hflow}{h_{\Gamma}^{(\tf)}}
\DeclareMathOperator{\Tr}{Tr}

\title{\textbf{Gradient-Flowed Quark and Gluon Quasi-PDFs in Lattice QCD}\\
\large Theory, Distillation, and the \(32^3\times96\) Implementations}
\author{Hongxin Dong}
\date{20 October 2026}

\begin{document}
\maketitle

\begin{abstract}
This note defines gradient-flowed quasi-PDFs and their implementation in
distillation. The production operator now starts from thin links, with no HYP
pre-smearing, so that its regulator scheme can be used with thin-link
gradient-flow quasi-PDF matching. Existing all-time light-quark perambulators
support the complete four-dimensional flowed connected calculation without
new Dirac inversions.
Time-diagonal strange perambulators support only a three-dimensional flowed
loop; complete four-dimensional flow needs off-diagonal time blocks. The
three-dimensional kernel suppresses spatial ultraviolet modes and regulates
the one-loop straight Wilson-line self-energy, but does not damp arbitrarily
large Euclidean energy and therefore does not inherit the standard all-order
gradient-flow finiteness or matching theorem. The
current \(32^3\times96\) outputs include finite-flow-time correlators and,
as of 20 September 2026, ringed matrix-element diagnostics from ten matched
configurations with 50 paired bootstrap replicas. The normalization uses a
published critical valence mass on a fixed massive sea. Its prescription
systematics, continuum and flow-time control, LaMET
matching, disconnected diagrams where required, and finite-\(z\) reconstruction
are still needed for a physical PDF.

For gluons, the present calculation instead flows the unsmeared gauge field,
constructs a four-leaf Clover field strength and straight spatial bilocals,
and forms vacuum-subtracted disconnected correlations with stored nucleon
two-point functions.  The raw terms are available at 14 flow times, and the
new covariance-preserving C3/C2 product uses a common 406-configuration sample
and resamples six operator choices to Nboot=1000.  The six choices are \(T\),
\(T+S\), and \(T-S\) in each channel, with even separation for unpolarized
and odd separation for helicity.  A convention audit finds that the stored
unpolarized combination is correct, whereas the stored helicity `combined'
array used a legacy plus sign.  Its raw terms are sufficient to construct
every requested choice without repeating gauge flow or the OPE measurement.
These raw products are finite-flow bare observables. A subsequent
unpolarized analysis propagates the 1000 paired replicas through flow
conversion, fixed-spacing flow-zero fits, finite-Ioffe-time reconstruction,
Fourier transformation, and a gluon-only LaMET inverse. A separate
20-configuration, 100-replica quark analysis covers three polarizations and
connected \(u,d,u-d\) flavors under the unvalidated \(Z_\chi=1\) assumption.
Section~\ref{sec:results-october} records these results, their source files,
and the frequent quark tail-fit boundary hits. Numerical completion does
not establish controlled physical PDFs.
Section~\ref{sec:blending-oct15} adds the Blending weight audit and the
completed ten-configuration L32x64 temporal-window tests. The subsequent
user-selected scheme B is isolated from the frozen historical scheme A;
implementation regression is not a proof of stochastic unbiasedness.
Section~\ref{sec:line-mass-oct20} records the L32x64 Wilson-loop analysis
with 1700 configurations and 3000 paired bootstrap replicas: direct
fixed-flow-pair time slopes, finite-time static energies, and candidate
line-mass coefficients. Flow-window/model dependence and adjoint signal
limitations currently prevent approval of a unique final coefficient.
No nonperturbative line-mass subtraction has been applied to the PDFs.
\end{abstract}

\tableofcontents

\section{Scope and central result}

Euclidean lattice QCD cannot evaluate a lightlike bilocal operator directly.
The quasi-PDF method instead uses an equal-time spatial bilocal in a nucleon
with large momentum \(P_z\), then matches it perturbatively to the light-front
PDF \cite{Ji2013}.

Gradient flow introduces a second ultraviolet resolution scale. At positive
flow time, fields are smooth over a radius
\begin{equation}
  \rflow=\sqrt{8\tf},
  \qquad [\tf]=\text{length}^2.
\end{equation}
Properly normalized flowed fields define a finite continuum operator at fixed
\(\tf>0\) \cite{Luscher2010,Luscher2013,MonahanOrginos2017}. This is a
short-distance scheme, not an \(\MSbar\) light-front PDF; flow-scheme and LaMET
matching remain necessary.

For the flowed-quark construction, the gauge flow, fermion flow, and Wilson
line all start from the same thin-link field.  For the flowed-gluon
construction, both Clover tensors and the Wilson line are built from that same
flowed thin-link gauge field.  A HYP-pre-smeared initial condition defines a
different finite-\(a\) scheme and is not used in either production
implementation described below.

The key identity for distillation is
\begin{equation}
  \chi(\tf,x)=\sum_y K(\tf;x,y)\psi(y),
  \qquad
  \bar\chi(\tf,x)=\sum_y\bar\psi(y)K^\dagger(\tf;y,x).
  \label{eq:kernel-solution-summary}
\end{equation}
Applying \(K\) to the all-time LapH basis moves the flow into the current
vertex and leaves the perambulators unchanged. Thus operator-only flow needs
no new inversion when the stored perambulators cover the temporal support of
\(K\). The light data satisfy this condition; the time-diagonal strange data
do not for complete four-dimensional flow.

\section{Euclidean quasi-PDF matrix elements}

\subsection{Spatially nonlocal operator}

Let the quark separation be \(z\hat{\bm z}\) on a fixed Euclidean time slice.
The conventional bare quasi-PDF operator is
\begin{equation}
  \cO_\Gamma(z,x)
  =
  \bar\psi(x)\,
  \Gamma\,
  W(x,x+z\hat{\bm z})\,
  \psi(x+z\hat{\bm z}),
  \label{eq:bare-quasi-operator}
\end{equation}
where
\begin{equation}
  W(x,x+z\hat{\bm z})
  =
  \mathcal{P}\exp\left[
    \ii g\int_0^z\dd\lambda\,
    A_z(x+\lambda\hat{\bm z})
  \right].
\end{equation}
Common unpolarized choices are \(\Gamma=\gamma_z\) and \(\gamma_4\). At finite
lattice spacing, Wilson-type fermions can mix some Dirac structures with
lower-symmetry operators. Matching must use the same \(\Gamma\), action, and
flowed scheme as the calculation.

The forward nucleon matrix element is
\begin{equation}
  h_\Gamma(z,P_z)
  =
  \langle N(P_z)|\cO_\Gamma(z,0)|N(P_z)\rangle.
\end{equation}
One Fourier convention for the quasi-PDF is
\begin{equation}
  \qtilde_\Gamma(x,P_z)
  =
  \int_{-\infty}^{\infty}\frac{\dd z}{4\pi}\,
  e^{-\ii xP_z z}\,
  h_\Gamma(z,P_z).
  \label{eq:quasi-fourier}
\end{equation}
The Fourier-sign convention is immaterial if used consistently.

\subsection{Two- and three-point functions}

For an interpolating operator \(N_\alpha\), define
\begin{align}
  C_2(\bm P,T;\Pi)
  &=
  \sum_{\bm x}e^{-\ii\bm P\cdot(\bm x-\bm x_s)}
  \Pi_{\alpha\beta}
  \langle N_\beta(\bm x,t_s+T)
          \overline N_\alpha(\bm x_s,t_s)\rangle,
  \\
  C_3^\Gamma(\bm P,T,\tau,z;\Pi)
  &=
  \sum_{\bm x,\bm y}e^{-\ii\bm P\cdot(\bm x-\bm x_s)}
  \Pi_{\alpha\beta}
  \langle N_\beta(\bm x,t_s+T)
          \cO_\Gamma(z,\bm y,t_s+\tau)
          \overline N_\alpha(\bm x_s,t_s)\rangle .
  \label{eq:c2-c3-definitions}
\end{align}
At \(0\ll\tau\ll T\), a spectral decomposition gives
\begin{align}
  C_2(\bm P,T;\Pi)
  &\longrightarrow
  |Z_N(\bm P)|^2 e^{-E_{\bm P}T}\,
  \mathcal A_2(\bm P,\Pi),
  \\
  C_3^\Gamma(\bm P,T,\tau,z;\Pi)
  &\longrightarrow
  |Z_N(\bm P)|^2 e^{-E_{\bm P}T}\,
  \mathcal A_3^\Gamma(\bm P,z,\Pi).
\end{align}
Factors of \(2E_{\bm P}\) depend on state normalization. A suitable \(C_3/C_2\)
ratio isolates the matrix element only after controlling excited states with
several source--sink separations and plateau, summation, or multi-state fits.
A raw three-point function is not a PDF.

\section{Four-dimensional gradient flow}

\subsection{Gauge flow}

Starting from the four-dimensional gauge field \(A_\mu(x)\), the continuum
flowed gauge field \(B_\mu(\tf,x)\) obeys
\begin{align}
  \partial_{\tf} B_\mu
  &=
  D_\nu G_{\nu\mu}
  +\alpha_0 D_\mu\partial_\nu B_\nu,
  &
  B_\mu(0,x)&=A_\mu(x),
  \label{eq:gauge-flow}
  \\
  G_{\mu\nu}
  &=
  \partial_\mu B_\nu-\partial_\nu B_\mu+[B_\mu,B_\nu],
  &
  D_\mu&=\partial_\mu+[B_\mu,\,\cdot\,].
\end{align}
The \(\alpha_0\) term fixes the flow-time gauge and drops out of gauge-invariant
observables. Lattice Wilson flow is the gauge-covariant steepest descent of the
Wilson action.

For the matching scheme adopted here, \(A_\mu\) is the original thin-link
field. No HYP transformation is applied before the initial condition in
Eq.~\eqref{eq:gauge-flow}.

\subsection{Fermion flow and its kernel}

For \(\alpha_0=0\), the quark flow equations are
\begin{align}
  \partial_{\tf}\chi(\tf,x)
  &=\Delta_B\chi(\tf,x),
  &\chi(0,x)&=\psi(x),
  \\
  \partial_{\tf}\bar\chi(\tf,x)
  &=\bar\chi(\tf,x)\overleftarrow{\Delta}_B,
  &\bar\chi(0,x)&=\bar\psi(x),
  \label{eq:fermion-flow}
\end{align}
where
\begin{equation}
  \Delta_B=D_\mu[B]D_\mu[B]
\end{equation}
is the four-dimensional covariant Laplacian. The kernel \(K\) is defined by
\begin{equation}
  \left[\partial_{\tf}-\Delta_{B(\tf)}^{(x)}\right]
  K(\tf;x,y)=0,
  \qquad
  K(0;x,y)=\delta_{xy}.
  \label{eq:kernel-equation}
\end{equation}
Because \(B_\mu\) evolves, \(K\) is a flow-time-ordered covariant heat kernel,
not simply \(\exp(\tf\Delta_{B(0)})\).

Under a gauge transformation \(\Omega\),
\begin{equation}
  K(\tf;x,y)\longrightarrow
  \Omega(x)K(\tf;x,y)\Omega^\dagger(y).
\end{equation}
This makes \(\chi\) covariant and the flowed bilocal gauge invariant.

In the free-field limit,
\begin{equation}
  K(\tf;p)=e^{-\tf p^2}.
\end{equation}
For the nearest-neighbor lattice Laplacian,
\begin{equation}
  \widehat p^{\,2}
  =
  \frac{4}{a^2}\sum_{\mu=1}^4
  \sin^2\left(\frac{ap_\mu}{2}\right),
  \qquad
  K_{\rm lat}(\tf;p)=e^{-\tf\widehat p^{\,2}}.
  \label{eq:lattice-free-kernel}
\end{equation}
The damping is four-dimensional, with radius \(\sqrt{8\tf}\).

\section{The gradient-flowed quasi-PDF operator}

\subsection{Definition and finite-flow scheme}

At positive flow time, define
\begin{equation}
  \cO_\Gamma^{(\tf)}(z,x)
  =
  \overline{\mathring\chi}(\tf,x)\,
  \Gamma\,
  W_{B(\tf)}(x,x+z\hat{\bm z})\,
  \mathring\chi(\tf,x+z\hat{\bm z}).
  \label{eq:flowed-quasi-operator}
\end{equation}
Here \(\mathring\chi\) is a ringed, multiplicatively normalized flowed field:
\begin{equation}
  \mathring\chi(\tf,x)
  =
  \zeta_\chi^{1/2}(\tf,\mu)\chi(\tf,x),
  \qquad
  \overline{\mathring\chi}(\tf,x)
  =
  \zeta_\chi^{1/2}(\tf,\mu)\bar\chi(\tf,x).
  \label{eq:ringed-field}
\end{equation}
At positive flow time, flowed fermions require the factor in
Eq.~\eqref{eq:ringed-field}, or an equivalent nonperturbative normalization
\cite{Luscher2013,MakinoSuzuki2014}.

The flowed matrix element and quasi-PDF are
\begin{align}
  \hflow(z,P_z)
  &=
  \langle N(P_z)|\cO_\Gamma^{(\tf)}(z,0)|N(P_z)\rangle,
  \\
  \qtilde_\Gamma^{(\tf)}(x,P_z)
  &=
  \int\frac{\dd z}{4\pi}\,
  e^{-\ii xP_z z}\,
  \hflow(z,P_z).
  \label{eq:flowed-quasi-pdf}
\end{align}
At fixed \(\tf>0\), flow regulates the Wilson-line power divergence. Its
short-distance behavior reappears in the matching coefficient as
\(\tf\to0\). Flow therefore defines a scheme; it does not eliminate matching.

\subsection{Scale hierarchy and matching}

A useful hierarchy is
\begin{equation}
  a \ll \rflow
  \ll P_z^{-1}
  \ll \Lambda_{\rm QCD}^{-1},
  \label{eq:scale-hierarchy}
\end{equation}
supplemented, where a fixed nonzero Wilson-line separation is used, by
\(\rflow\ll |z|\). Equivalently,
\begin{equation}
  \frac{a^2}{\tf}\ll1,
  \qquad
  \tf P_z^2\ll1,
  \qquad
  \tf\Lambda_{\rm QCD}^2\ll1,
  \qquad
  P_z^2\gg M_N^2,\Lambda_{\rm QCD}^2.
\end{equation}
This window can be narrow and must be tested numerically.

The factorization relation can be written schematically as
\begin{equation}
  \qtilde_\Gamma^{(\tf)}(x,P_z)
  =
  \int_{-1}^{1}\frac{\dd y}{|y|}\,
  C_\Gamma^{(\tf)}
  \left(
    \frac{x}{y},
    \frac{\mu}{|y|P_z},
    \sqrt{8\tf}\,|y|P_z,
    \sqrt{8\tf}\,\mu
  \right)
  q_\Gamma(y,\mu)
  +\delta_{\rm power},
  \label{eq:factorization}
\end{equation}
where \(q_\Gamma(y,\mu)\) is usually defined in \(\MSbar\). Schematically,
\begin{equation}
  \delta_{\rm power}
  =
  \cO\left(
    \frac{M_N^2}{P_z^2},
    \frac{\Lambda_{\rm QCD}^2}{P_z^2},
    \tf\Lambda_{\rm QCD}^2
  \right),
\end{equation}
plus discretization, finite-volume, and finite-\(z\) effects. The coefficient
\(C_\Gamma^{(\tf)}\) depends on the flowed operator and its normalization
\cite{MonahanOrginos2017,Monahan2018}; an unflowed kernel cannot be applied to
the raw flowed matrix element. Likewise, a thin-link flowed coefficient must
not be applied to a HYP-pre-smeared flow without an explicit scheme
conversion.

For nonsinglet moments, equivalently,
\begin{equation}
  \int\dd x\,x^{n-1}\qtilde_\Gamma^{(\tf)}(x,P_z)
  =
  C_n(\tf,\mu,P_z)
  \int_{-1}^{1}\dd x\,x^{n-1}q_\Gamma(x,\mu)
  +\text{power corrections}.
\end{equation}
The anomalous dimensions of twist-two operators determine the flow-time
dependence of \(C_n\).

\subsection{Order of limits}

The controlled order is:
\begin{enumerate}
  \item take \(a\to0\) at fixed physical \(\tf>0\), \(z\), and \(P_z\);
  \item normalize and match the flowed scheme where \(a^2/\tf\ll1\) and
        \(\tf P_z^2\ll1\);
  \item apply large-momentum matching and power corrections;
  \item reconstruct the \(x\) dependence from finite spatial separations.
\end{enumerate}
These dependences may be fitted jointly, but the scale hierarchy must remain
explicit.

\subsection{Operational matching workflow}
\label{sec:matching-workflow}

Matching is not applied to a raw three-point function.  The following order
keeps the lattice, flow, quasi-PDF, and light-cone schemes distinct.

\paragraph{Step 1: extract the finite-flow matrix element.}
For each \((z,P_z,\tf)\), fit \(C_2\) and \(C_3\) at several source--sink
separations.  After the chosen spin projection and kinematic normalization,
this gives
\begin{equation}
  h_{\Gamma}^{\rm lat}(z,P_z,\tf,a)
  =
  \langle P_z|
  \bar\chi(\tf,0)\Gamma W_{B(\tf)}(0,z)\chi(\tf,z)
  |P_z\rangle_a .
  \label{eq:matching-step-lattice-matrix-element}
\end{equation}
Excited-state, contact-term, finite-volume, and temporal-wrap effects are
removed or included in the fit before perturbative matching.

\paragraph{Step 2: normalize the flowed quark fields.}
Apply the ringed-field factor in the same discretization and convention used
to derive the coefficient:
\begin{equation}
  h_{\Gamma}^{\rm ring}(z,P_z,\tf,a)
  =\zeta_\chi(\tf,a,\mu_{\rm F})
   h_{\Gamma}^{\rm lat}(z,P_z,\tf,a).
  \label{eq:matching-step-ringed}
\end{equation}
There are two flowed fermion legs and hence one factor \(\zeta_\chi\) in the
bilinear.  A ratio may cancel part of this normalization, but cancellation
must be demonstrated for that ratio; it must not be assumed from the
cancellation of nucleon overlap factors.

\paragraph{Step 3: take the continuum limit at fixed physical flow time.}
Interpolate ensembles to common physical \(z,P_z,\tf\) and form
\begin{equation}
  h_{\Gamma}^{\rm GF}(z,P_z,\tf)
  =\lim_{a\to0}
  h_{\Gamma}^{\rm ring}(z,P_z,\tf,a),
  \qquad \frac{a^2}{\tf}\to0.
  \label{eq:matching-step-continuum}
\end{equation}
The flow time must not be held fixed only in lattice units.  Several \(\tf\)
values are needed to identify \(a^2/\tf\) artifacts and
\(\tf P_z^2\) power corrections.

\paragraph{Step 4: determine the flow-scheme conversion.}
For a two-step treatment, write the small-flow-time expansion of the nonlocal
operator schematically as
\begin{equation}
  \cO_{\Gamma}^{\rm GF}(\tf,z)
  =
  \int\!\dd z'\,
  \mathcal Z_{\Gamma}^{\rm GF\to q}
    (z,z';\tf,\mu)
  \cO_{\Gamma}^{\rm quasi,\MSbar}(z',\mu)
  +\cO(\tf\Lambda_{\rm QCD}^2,\tf/z^2).
  \label{eq:flow-to-quasi-conversion}
\end{equation}
In a multiplicatively renormalizable coordinate-space channel,
\(\mathcal Z\) is diagonal in \(z\); the more general notation allows Dirac,
singlet--gluon, and lattice-symmetry mixing.  Its perturbative terms contain
the finite line contribution proportional to \(|z|/\sqrt{\tf}\), logarithms
such as \(\ln(\mu^2\tf)\), and endpoint constants.  All are scheme dependent.

The coefficient is obtained from renormalized partonic matrix elements.  With
the same infrared regulator on both sides,
\begin{equation}
  \mathcal Z_{\Gamma}^{(1)}
  =
  \left[h_{\Gamma,\mathrm{parton}}^{\mathrm{GF},(1)}\right]_{\mathrm{ren}}
  -
  \left[h_{\Gamma,\mathrm{parton}}^{\mathrm{quasi},\MSbar,(1)}\right]_{\mathrm{ren}},
  \label{eq:matching-coefficient-difference}
\end{equation}
because infrared poles and logarithms cancel in the difference.  The flowed
partonic calculation must use the same 4D kernel, ringed normalization,
\(\Gamma\), Wilson-line orientation, and thin-link initial condition as the
lattice operator.

\paragraph{Step 5: perform LaMET matching.}
After Fourier transformation, or directly in coordinate space, the
renormalized quasi-PDF obeys
\begin{equation}
  \widetilde q_{\Gamma}^{\rm quasi}(x,P_z,\mu)
  =
  \int_{-1}^{1}\frac{\dd y}{|y|}
  C_{\Gamma}^{\rm LaMET}
  \left(\frac{x}{y},\frac{\mu}{|y|P_z}\right)
  q_{\Gamma}(y,\mu)
  +\cO\!\left(
    \frac{M_N^2}{P_z^2},
    \frac{\Lambda_{\rm QCD}^2}{P_z^2}
  \right).
  \label{eq:standard-lamet-step}
\end{equation}
Equations~\eqref{eq:flow-to-quasi-conversion} and
\eqref{eq:standard-lamet-step} may be combined into the one-step flowed
coefficient in Eq.~\eqref{eq:factorization}.  One must use either a consistent
two-step factorization or a consistent one-step coefficient, not multiply
coefficients that were defined with overlapping subtractions.

\paragraph{Step 6: reconstruct \(x\) and test the matching window.}
The lattice supplies finitely many noisy separations.  Fourier truncation,
Bayesian reconstruction, derivative methods, or coordinate-space fits must
propagate their own uncertainties.  The final result must be stable under
variations of
\begin{equation}
  P_z,\quad \tf,\quad a,\quad z_{\max},\quad
  \mu,\quad \text{fit model},\quad \text{matching order}.
  \label{eq:matching-systematic-variations}
\end{equation}
For flavor singlets, quark and gluon distributions form a matching matrix;
connected \(u\) and \(d\) contractions alone are not a complete singlet PDF.

\subsection{Order of $\lambda$ treatment, Fourier transformation, and quasi matching}
\label{sec:lambda-fourier-order}

Three operations must be distinguished.  The small-flow-time limit is
$\tf\to0$.  The coordinate variable is $\lambda=zP_z$ (with $\zeta^2=z^2$),
and a coordinate-space analysis may interpolate in $\lambda$, model the tail
beyond the largest measured $|\lambda|$, or extrapolate the $\zeta^2$
dependence at fixed $\lambda$.  None of these operations means taking a
physical $\lambda\to\infty$ limit.  A tail model is only a reconstruction
device for the Fourier integral; a $\zeta^2\to0$ extrapolation belongs to a
pseudo-ITD analysis.

For the $x$-space gluon quasi-to-light-cone kernels used here, the recommended
two-step chain is
\begin{equation}
\begin{gathered}
 h_{\tau}^{\rm GF}(z,P_z)
 \xrightarrow{\rm GF\ conversion} h_R(z,P_z;\tau)
 \xrightarrow{\tau\to0} h_R(z,P_z) \\
 \xrightarrow{\lambda\ \rm interpolation/tail} h_R(\lambda,P_z)
 \xrightarrow{\rm Fourier} \widetilde g_{\rm quasi}(x,P_z)
 \xrightarrow{C_{gg}^{-1}\ \rm (and\ mixing)} g(x,\mu).
\end{gathered}
 \label{eq:lambda-fourier-chain}
\end{equation}
For an even unpolarized matrix element, the convention used in the present
diagnostic is
\begin{equation}
 x\widetilde g(x,P_z)
 =\frac{P_z}{\pi}\int_0^{\infty}\!\dd z\,
 \cos(xP_z z)\,h_R(z,P_z),
 \label{eq:gluon-quasi-fourier-order}
\end{equation}
with the corresponding odd sine transform for the helicity channel.  A finite
$z_{\max}$, interpolation prescription, and tail model must be varied as
reconstruction systematics.  Missing short-distance points are invalid data
and must remain masked rather than zero-filled.

The kernel used in arXiv:2510.26425v2 is an $x$-space convolution, so the
Fourier transform (including any finite-$\lambda$ reconstruction) precedes
that quasi-to-light-cone matching.  A coordinate-space kernel can instead be
used before Fourier transformation only when it is derived in the same scheme
and the transform is performed consistently.  With finite $z_{\max}$,
interpolation, ratio normalization, or bootstrap-dependent nonlinear steps,
the two numerical procedures are not interchangeable.  The GF-to-quasi
conversion and the subsequent quasi-to-light-cone LaMET matching are separate
coefficients; applying both twice would double count a subtraction.

The current N406 gluon product follows the first part of this chain: it uses a
$\tau/a^2\ge1$ flow-time extrapolated coordinate-space matrix element, then a
finite-$z$ cosine transform, and only afterwards the one-loop gluon kernel.
This remains a matched quasi diagnostic rather than a physical PDF because
the nonperturbative ratio/hybrid denominator, singlet quark--gluon mixing,
continuum limit, finite-$\lambda$ tail uncertainty, and target-mass and
higher-twist corrections have not yet been controlled.

\paragraph{What changes for 3D flow.}
Every occurrence of \(\mathcal Z_{\Gamma}^{\rm GF\to q}\) or the direct
flowed coefficient must be recalculated with
\(F_{\tf}^{\rm 3D}=e^{-2\tf\bm k^2}\).  Substituting a 3D lattice matrix
element into a 4D coefficient leaves a finite, generally
\(z\)- and \(P_z\)-dependent scheme error even if both matrix elements have a
well-defined continuum limit.  Agreement at \(\tf=0\) does not determine this
finite conversion at \(\tf>0\).

\subsection{Explicit small-flow-time coefficient for the ringed bilocal}
\label{sec:explicit-sfte}

The general conversion in Eq.~\eqref{eq:flow-to-quasi-conversion} becomes
multiplicative in the asymptotic small-flow region for the straight
off-lightcone bilocal.  The corrected one-loop result of Brambilla and Wang
\cite{BrambillaWang2024} is
\begin{equation}
  \cO_{\Gamma}^{\mathrm{GF,R}}(z,\tf)
  =
  \mathcal C_\psi(\tf,\mu)
  \exp\!\left[\delta m(\tf)|z|\right]
  \cO_{\Gamma}^{\mathrm{quasi},\MSbar}(z,\mu)
  +\cO(\tf).
  \label{eq:explicit-sfte-operator}
\end{equation}
Here ``R'' denotes ringed flowed fermions.  For one flavor a convenient
nonperturbative definition is
\begin{equation}
  \mathring\chi(\tf,x)
  =\mathring Z_\chi^{1/2}(\tf)\chi(\tf,x),
  \qquad
  \mathring Z_\chi(\tf)
  =
  \frac{-2N_c}
  {(4\pi)^2\tf^2
   \left\langle
     \bar\chi(\tf,x)
     \overleftrightarrow{\slashed{\mathcal D}}
     \chi(\tf,x)
   \right\rangle_{m=0}}.
  \label{eq:ringed-zchi-nonperturbative}
\end{equation}
The derivative discretization and the massless or chiral-extrapolation
prescription are part of the scheme.  Since the bilocal has two flowed
endpoints, it contains one overall factor \(\mathring Z_\chi\), not two:
\begin{equation}
  h_\Gamma^{\mathrm{GF,R}}
  =\mathring Z_\chi h_\Gamma^{\mathrm{GF}}.
  \label{eq:ringed-bilocal-factor}
\end{equation}

In the convention of Eq.~\eqref{eq:explicit-sfte-operator},
\begin{align}
  \mathcal C_\psi(\tf,\mu)
  &={}
  1-\frac{\alpha_s(\mu)C_F}{4\pi}
  \left[
    3\ln\!\left(2\mu^2\tf e^{\gamma_E}\right)
    +2+\ln 432
  \right]
  +\cO(\alpha_s^2),
  \label{eq:cpsi-one-loop}
  \\
  \delta m(\tf)
  &={}
  -\frac{\alpha_s(\mu)C_F}{4\pi}
  \frac{\sqrt{2\pi}}{\sqrt{\tf}}
  +\cO(\alpha_s^2).
  \label{eq:delta-m-one-loop}
\end{align}
The inverse conversion applied to a lattice matrix element is therefore
\begin{equation}
  h_\Gamma^{\mathrm{quasi},\MSbar}(z,P_z,\mu)
  =
  \frac{\exp[-\delta m(\tf)|z|]}
       {\mathcal C_\psi(\tf,\mu)}
  \mathring Z_\chi(\tf)
  h_\Gamma^{\mathrm{GF}}(z,P_z,\tf)
  +\cO\!\left(
    \frac{\tf}{z^2},\tf P_z^2,
    \tf\Lambda_{\rm QCD}^2
  \right).
  \label{eq:explicit-sfte-matrix-element}
\end{equation}
The scale
\begin{equation}
  \mu_{\tf}
  =\frac{e^{-\gamma_E/2}}{\sqrt{2\tf}}
  \label{eq:sfte-natural-scale}
\end{equation}
removes the explicit logarithm in Eq.~\eqref{eq:cpsi-one-loop}; one may match
near \(\mu_{\tf}\) and evolve to a common scale.  The coefficient displayed
above is independent of \(\Gamma\) at this order in the stated limit.  This
does not make the later LaMET kernel universal: vector, axial, and tensor
quasi-PDFs still require their own kernels.

Equation~\eqref{eq:explicit-sfte-operator} assumes four-dimensional gauge and
fermion flow, two endpoints and a straight Wilson line flowed to the same
\(\tf\), ringed fermions, massless quarks, and no unaccounted operator mixing.
It uses the corrected 2024 result.  The affected expression in the original
2018 calculation \cite{Monahan2018} should not be used without its erratum and
the later correction.

\subsection{Independent ringed-fermion measurement: implementation update}
\label{sec:ringed-kinetic-implementation}

This subsection records the derivation, implementation, and measurements
through 20 September 2026. Numerical tests and the ten-configuration paired
normalization are complete in a published-critical-valence-mass prescription.
The original finite-mass pilot is retained below as historical evidence;
it is not the normalization used for the new plots. Prescription systematics
and formal physical acceptance remain open. The existing nucleon
correlators are retained unchanged.  The new calculation measures the vacuum
quantity in Eq.~\eqref{eq:ringed-zchi-nonperturbative}, not another nucleon
three-point function and not the difference of the degenerate light-flavor
vacuum loops \cite{MakinoSuzuki2014,Luscher2013}.

\subsubsection{Why this normalization is needed, and where its numerator comes from}
\label{sec:ringed-derivation-detail}

The ratio \(C_3/C_2\) removes nucleon overlap factors. Its denominator does
not contain the two flowed fermion fields at the current insertion, so their
normalization does not cancel. Flow damps high momenta but ordinary flowed
fermions still require multiplicative field renormalization. Ringed fields
fix that scale through a vacuum observable \cite{MakinoSuzuki2014}.
Here \(Z_\chi^{\rm code}\equiv\mathring Z_\chi\): the field
named \texttt{Zchi} in NPZ files is the ringed factor, not the unringed
minimal-subtraction wave-function constant also denoted \(Z_\chi\) in some
papers. We use the one-flavor, massless prescription of
Ref.~\cite{BrambillaWang2024}.

The continuum normalization can be derived explicitly. With Hermitian
Euclidean gamma matrices, \(\{\gamma_\mu,\gamma_\nu\}=2\delta_{\mu\nu}\),
the free massless propagator and heat kernel are
\begin{equation}
 S_0(p)=\frac{-\ii\slashed p}{p^2},\qquad
 K_{\tf}(p)=e^{-\tf p^2},\qquad
 \overleftrightarrow{\slashed\partial}\longrightarrow 2\ii\slashed p.
\end{equation}
Two flowed endpoints give \(K K^\dagger=e^{-2\tf p^2}\). The closed
Grassmann contraction has a minus sign, and the spin trace is four:
\begin{align}
 E_0(\tf)
 &=-2N_c\int\frac{\dd^4p}{(2\pi)^4}e^{-2\tf p^2}
       \operatorname{tr}_{s}\!\left[(\ii\slashed p)
                                  \frac{-\ii\slashed p}{p^2}\right]
 \nonumber\\
 &=-8N_c\int\frac{\dd^4p}{(2\pi)^4}e^{-2\tf p^2}
  =-\frac{2N_c}{(4\pi)^2\tf^2},
 \label{eq:ringed-free-numerator}
\end{align}
where \(\int\dd^4p\,e^{-\alpha p^2}/(2\pi)^4=1/(16\pi^2\alpha^2)\).
Thus \(\mathring Z_\chi=E_0/E\), with
\(E=\langle\bar\chi\overleftrightarrow{\slashed D}\chi\rangle_{m_R=0}\).
The dimension of \(E\) is mass to the fourth power, so \(\tf^2E\) and
\(\mathring Z_\chi\) are dimensionless. No absolute value should be
inserted to hide the sign of the measured kinetic expectation.

A constant field rescaling \(\chi\mapsto c^{1/2}\chi\),
\(\bar\chi\mapsto c^{1/2}\bar\chi\) gives \(E\mapsto cE\) and
\(\mathring Z_\chi\mapsto\mathring Z_\chi/c\). The ringed fields
are invariant under this common normalization change. The two endpoints of
the bilocal give exactly one \(\mathring Z_\chi\), as in
Eq.~\eqref{eq:ringed-bilocal-factor}; neither its square nor its square root
is the bilocal factor. The same scalar can be applied to the already fitted
matrix element. This reuses the declared fit and its replicas; it does not
assert equivalence to a new fit with a changed covariance or new window selection.

For degenerate light valence quarks, \(h_{u-d}^{R}=\mathring Z_\chi
(h_u^{\rm GF}-h_d^{\rm GF})\). The denominator is a \emph{one-flavor}
vacuum kinetic trace. The vacuum \(u-d\) difference would vanish. If one
instead sums over \(N_f\) degenerate flavors, its numerator must also
contain \(N_f\) \cite{MakinoSuzuki2014}.

\subsubsection{Vacuum trace, lattice units, and the implemented estimator}

\paragraph{From a vacuum contraction to a stochastic estimator.}
Let \(\theta=\tf/a^2\), \(V_4=N_xN_yN_zN_t\), and
\(S=D_{\rm val}^{-1}\) in the mass-normalized lattice-field convention.
The flow kernel \(K_\theta\) is spin diagonal but is fully four dimensional.
With a volume average and one flavor, the dimensionless kinetic expectation is
\begin{align}
 E_{\rm lat}(\theta)
 &\equiv a^4\left\langle\bar\chi
              \overleftrightarrow{\slashed{\mathcal D}}\chi\right\rangle
 =-\frac{2}{V_4}\left\langle
       \Tr\!\left[\slashed\nabla_{B(\theta)}
                   K_\theta S K_\theta^\dagger\right]\right\rangle_U,
 \label{eq:ringed-kinetic-full-trace}\\
 (\nabla_\mu f)(x)
 &=\frac{1}{2}\left[
 U_\mu(\theta,x)f(x+\hat\mu)
 -U_\mu^\dagger(\theta,x-\hat\mu)f(x-\hat\mu)\right].
 \label{eq:ringed-kinetic-lattice-derivative}
\end{align}
Here \(\nabla_\mu=a\mathcal D_\mu+\cO(a^3)\).  The minus sign is the
Grassmann contraction; integration by parts in the volume sum gives the
factor two.  The trace includes all sites, four spins, and three colors:
there is no extra division by spin or color dimensions.
Writing \(\tf=a^2\theta\) and \(E_{\rm lat}=a^4E\) gives
\begin{equation}
 \mathring Z_\chi(\theta)
 =\frac{-6}{(4\pi)^2\theta^2\langle\operatorname{Re}E_{{\rm lat},U}\rangle_U}.
 \label{eq:ringed-lattice-z}
\end{equation}
The powers of \(a\) cancel: the code must not multiply by another \(a^4\)
or convert the kinetic expectation to \(\mathrm{fm}^{-4}\) first.

For full-volume complex \(Z_4\) noise with
\(\mathbb E_\eta[\eta\eta^\dagger]=\mathbf1\), solve once and flow twice:
\begin{equation}
 D_{\rm val}q_r=\eta_r,\qquad
 \rho_r(\theta)=K_\theta\eta_r,\qquad
 \phi_r(\theta)=K_\theta q_r .
\end{equation}
Both fields must follow the same thin-link gauge history.  Cyclicity of the
trace then gives the unbiased noise estimator
\begin{equation}
 \widehat E_{{\rm lat},U}(\theta)
 =-\frac{2}{V_4N_\eta}\sum_{r=1}^{N_\eta}
       \rho_r^\dagger(\theta)
       \slashed\nabla_{B(\theta)}\phi_r(\theta).
 \label{eq:ringed-noise-estimator}
\end{equation}
The complex samples are saved before the real projection.  The noise is not
rescaled to unit norm.  A single inverse per noise serves every flow
checkpoint; neither a dense \(K\) nor an explicit adjoint-flow integration is
needed.  All temporal slices are retained in this measurement.
Explicitly, the identity used is
\begin{equation}
 \mathbb E_\eta[\eta^\dagger K^\dagger\slashed\nabla_BKS\eta]
 =\Tr[K^\dagger\slashed\nabla_BKS]
 =\Tr[\slashed\nabla_BKSK^\dagger].
\end{equation}
The estimate of \(E_U\) is unbiased in noise; its finite-sample reciprocal
need not be. In \texttt{core.py}, \texttt{kinetic\_sample} evaluates
\(-2\,\texttt{vdot}(\rho,\slashed\nabla\phi)/V_4\).
The fermion axes are \((t,z,y,x,s,c)\); gauge axes are
\((\mu,t,z,y,x,c,c')\). The index \(\mu=0,1,2,3\) means
\(x,y,z,t\), and its array-axis index is \(3-\mu\). The gamma basis is
the DeGrand--Rossi basis of \texttt{gamma\_dr()}, shared with the flowed
current code. Noise components are independent elements of
\(\{1,\ii,-1,-\ii\}\), with no unit-vector normalization.

\paragraph{A finite lattice is not the continuum tree-level limit.}
The exact free Wilson reference used by \texttt{free\_wilson\_kinetic} is
\begin{align}
 E_{\rm lat}^{\rm free}(\theta,m)
 &=-\frac{8N_c}{V_4}\sum_p e^{-2\theta\hat p^2}
 \frac{\sum_\mu\sin^2p_\mu}
 {\left[m+\sum_\mu(1-\cos p_\mu)\right]^2+\sum_\mu\sin^2p_\mu},
 \label{eq:ringed-free-lattice}\\
 \hat p^2&=4\sum_\mu\sin^2(p_\mu/2),\qquad
 p_t=2\pi(n_t+\tfrac12)/N_t.
\end{align}
At fixed lattice spacing the trace remains ultraviolet-cutoff regulated as
\(\theta\to0\), while the continuum numerator in
Eq.~\eqref{eq:ringed-lattice-z} grows as \(\theta^{-2}\).
Consequently \(\mathring Z_\chi\) need not be near one at small
\(\theta\), even in the free theory. This is a cutoff diagnostic, not a
proof that every observed deviation is a cutoff effect. Replacing the
continuum numerator by a lattice tree-level value would define an additional
finite-\(a\) normalization prescription; the present analysis does not do so.

\paragraph{Three distinct uses of gauge links.}
The historical light-perambulator source script
\nolinkurl{/public/home/sush/EasyDistillation/work/perambulator.3.py}
specifies \(am=-0.2295\), \(c_{\rm SW}=1.141151096\), isotropy, and one
four-dimensional stout step with \(\rho=0.125\) for the valence inverse.
That is distinct both from thin-link operator flow and from the 20-step
stout links used for the distillation basis.  The new pilot reproduces the
stated valence action while starting \(K\) and \(B\) from the thin links.
The association of every historical perambulator artifact with this source
script remains an acceptance check, not a conclusion from its filename.

For the full, unpreconditioned MAT solution, QUDA's asymmetric mass
normalization gives \(D_{\rm mass}=D_\kappa/(2\kappa)\), with
\(\kappa=[2(am+4)]^{-1}\).  The implementation checks the full residual
\(\|D_{\rm mass}q-\eta\|/\|\eta\|\) and inserts no additional
\(2\kappa\).  Temporal boundary conditions are antiperiodic.  The boundary
phase is embedded during flow and removed from the saved gauge before the
derivative applies it once.  The integrator is coupled gauge/fermion RK3,
not Euler; the six checkpoints use
\(\theta=0.05,0.1,0.2,0.5,0.8,1.0\) with step size \(0.01\).

\paragraph{Paired statistics and the massless gate.}
For bootstrap draw \(b\), use exactly the original bare-fit configuration
indices \(I_b\), average the per-configuration noise means, and only then
take the inverse:
\begin{equation}
 \overline E^{(b)}_{\rm lat}(\theta)
 =\frac{1}{|I_b|}\sum_{U\in I_b}\operatorname{Re}
                  \widehat E_{{\rm lat},U}(\theta),\qquad
 \mathring Z_\chi^{(b)}(\theta)
 =\frac{-2N_c}{(4\pi)^2\theta^2\overline E^{(b)}_{\rm lat}(\theta)}.
 \label{eq:ringed-paired-bootstrap}
\end{equation}
Multiply the same replica of \(h^{\rm GF}\) by this factor once.
``Matched configurations'' means that the kinetic means use the same
configurations and ordering as the bare fit. It does not replace the
ensemble scalar by a random per-configuration factor inside the operator.
Averaging \(1/E_U\), or multiplying each configuration's ratio by such an
inverse, defines a different estimator. For \(N_b=50\), report
\begin{align}
 h^R_{j\ell}&=Z_j h^{\rm bare}_{j\ell},\qquad
 h^{R,(b)}_{j\ell}=Z_j^{(b)}h^{{\rm bare},(b)}_{j\ell},\\
 \sigma^2_{j\ell}&=\frac1{N_b-1}\sum_b
 (h^{R,(b)}_{j\ell}-\overline h^R_{j\ell,\rm boot})^2,
 \label{eq:ringed-product-error}
\end{align}
where \(j\) labels flow time and \(\ell=(P_z,z)\); the current plots show
the real unpolarized component. The center retains the original bare-fit
central estimator, rather than replacing it by the replica mean.
There is no further division by \(\sqrt{N_b}\). Cross-flow and cross-point
covariances use these same replicas. At linear order,
\begin{equation}
 \operatorname{Var}(Zh)\simeq Z^2\operatorname{Var}(h)
 +h^2\operatorname{Var}(Z)+2Zh\operatorname{Cov}(Z,h).
\end{equation}
Independent bootstrap draws for \(Z\) and \(h\), or independent quadrature
of their errors, discard the last term. An additional naive
nested-noise bootstrap would double-count noise already present in the
configuration means; noise convergence must nevertheless be checked using
the retained individual measurements.

Neither \(am=-0.2295\) nor an arbitrary choice \(am=0\) establishes the
interacting chiral limit. The 18 September pilot was explicitly finite mass.
The original formal exporter accepts a mass-independent ringed factor only after an
independently tuned critical valence mass and matching field normalization
have been documented.  A chiral extrapolation, including PCAC-mass
uncertainties, is a separate alternative and is not implemented by simply
relabelling that pilot. Its diagnostic exports keep formal \(Z_\chi\) as
NaN and are rejected by the formal SFTE consumer. The new paired NPZ products
below contain finite diagnostic values with explicit non-acceptance flags;
they do not promote that older interface to a validated product.

The \(c_{\rm fl}\) contact subtraction in Eq.~(6.8) of
Ref.~\cite{Luscher2013} replaces \(KSK^\dagger\) by
\(K(S-c_{\rm fl}\mathbf1)K^\dagger\).  Its contribution to the full kinetic
trace vanishes for the spin-independent flow used here, since
\(\Tr[\slashed\nabla K K^\dagger]=0\) by
\(\operatorname{tr}_{\rm spin}\gamma_\mu=0\). A finite-noise estimate of
that zero trace need not vanish for each noise. This observation does not
establish complete \(\cO(a)\) improvement or remove mass-dependent and
cutoff effects.

\paragraph{Code and validation boundary.}
The independent module resides in
\begin{center}
\nolinkurl{/public/group/lqcd/donghx/Gradient_Flow_GluonPDF/}\par
\nolinkurl{quark_ringed_kinetic_v1/}.
\end{center}
Its entry points are \texttt{measure.py}, \texttt{collect.py}, and
\texttt{compare.py}; the directory README gives reproducible commands.
The full-volume stencil currently requires one MPI rank and one GPU.
The implementation has CPU tests for gamma conventions, APBC, gauge
covariance, exact free-field traces, same-noise Fourier comparisons, RK3
step halving, immutable restart records, and paired bootstrap export.
The following GPU measurements now supplement those CPU tests. No existing
nucleon matrix element was overwritten.

\paragraph{Completed GPU validation and finite-mass pilot.}
Jobs 3643531, 3643532, and 3643534 completed successfully on single A800 GPUs.
On a \(4^3\times8\) free lattice, four noise vectors were compared to the
exact lattice Fourier estimator with the same noise. The maximum relative
error decreased from \(1.4938\times10^{-5}\) at step size \(0.01\) to
\(1.8037\times10^{-6}\) at \(0.005\), consistent with the expected RK3
order. For configuration 10000, paired two-noise comparisons gave maximum
relative differences \(1.8625\times10^{-15}\) for checkpoint versus direct
flow, \(1.4944\times10^{-5}\) for step halving, and
\(3.3440\times10^{-12}\) for the corrected versus original QUDA backend.
The exact-unit-gauge test exposed a zero-generator matrix-exponential
singularity. An isolated library now returns the original link for an
exactly zero generator; the nonzero branch and shared installation are
unchanged. Actual loaded-library hashes are recorded.

At \(am=-0.2295\), the 32-noise pilot took 11 minutes 48 seconds. Its
maximum full-MAT residual was \(7.4581\times10^{-11}\), and the maximum
gauge unitarity/determinant deviation was \(4.3993\times10^{-14}\).
The measured real kinetic expectations and their noise standard errors are
\begin{center}
\begin{tabular}{rrr}
\toprule
\(\theta\) & \(\operatorname{Re}\widehat E_{{\rm lat},U}\) & Noise standard error \\
\midrule
0.05 & -2.23775135 & \(8.2594\times10^{-5}\) \\
0.10 & -1.36202226 & \(6.2976\times10^{-5}\) \\
0.20 & -0.59540806 & \(4.3218\times10^{-5}\) \\
0.50 & -0.12081402 & \(2.2398\times10^{-5}\) \\
0.80 & -0.04722902 & \(1.5581\times10^{-5}\) \\
1.00 & -0.02984961 & \(1.2973\times10^{-5}\) \\
\bottomrule
\end{tabular}
\end{center}
These are stochastic errors on one gauge configuration, not ensemble
uncertainties. The imaginary means are each within 1.64 noise standard
errors of zero. This validates the tested numerical chain, not a massless
ringed factor. Full results, input hashes, and immutable run directories
are recorded under the module's
\nolinkurl{validation/gpu_summary_20260918.json} and
\nolinkurl{validation/gpu_report_20260918.md}.

\subsubsection{Critical mass: literature input and prescription uncertainty}
\label{sec:ringed-critical-mass}
The ensemble HMC input independently confirms the stated valence-action
parameters. The later literature check identifies
\(am_{\rm crit}=-0.23539(3)\) for F32P30 in Table 3 of
Ref.~\cite{DuEtAl2024CriticalMass}. Its \(\beta=6.41\),
\(u_0=0.863437\), \(v_0=0.956942\),
\(c_{\rm SW}=v_0^{-3}=1.141151096\), and one-step stout action agree
with the checked local action. The vacuum calculation uses this central
value, while the sea masses remain \(am_l=-0.2295\), \(am_s=-0.2050\).
No local PCAC remeasurement or full sea-quark chiral extrapolation has been
performed. Pion and \(\eta_s\) masses identify the ensemble; they are not
Dirac bare-mass inputs. The older \(-0.23545(3)\) in
Ref.~\cite{HuEtAl2023CriticalMass} is a separate sensitivity choice, not an
independent measurement to average with the newer value.

If independent local tuning is required, measure nonsinglet
pseudoscalar--pseudoscalar and axial--pseudoscalar correlators with declared
Euclidean phases and source conjugation. In lattice units, a chosen axial
improvement convention gives
\begin{equation}
 am_{\rm PCAC}(n)=
 \frac{\tfrac12(\partial_n+\partial_n^*)C_{AP}(n)
       +c_A\partial_n^*\partial_n C_{PP}(n)}{2C_{PP}(n)}.
\end{equation}
Use common configurations and sources at several nearby valence masses,
check the time plateaus, and fit the zero of
\(am_{\rm PCAC}=A(am_0-am_{\rm crit})\), including fit-window and curvature
checks. An unimproved axial-current prescription must explicitly set
\(c_A=0\) and retain its cutoff uncertainty. Existing nucleon two-point
functions do not supply these meson correlators. Tuning only the valence
PCAC mass to zero on this fixed-sea-mass ensemble defines a valence chiral
prescription, not a full sea-quark chiral limit. Critical-mass uncertainty
must propagate into the kinetic expectation and \(\mathring Z_\chi\).
The ten configurations and their paired bootstrap indices have now been
completed. Action-to-artifact provenance, prescription uncertainty, and the
massless/improvement review remain separate issues. The finite-mass pilot
itself retains \texttt{measured\_unvalidated} status.

A four-noise, single-configuration study used identical noise IDs at
\(am_c=-0.23539\) and \(am_c\pm0.00003\). Define
\begin{equation}
 \Delta_m Z_j=\frac{Z_j(am_c+\sigma_m)-Z_j(am_c-\sigma_m)}2,
 \qquad (C_m)_{jk}\simeq\Delta_mZ_j\Delta_mZ_k.
 \label{eq:ringed-mass-response}
\end{equation}
The six signed responses \(100\Delta_mZ/Z\) were
\(0.00235,0.00286,0.00410,0.00663,0.00486,0.00111\) percent,
with paired-noise uncertainties
\(0.00090,0.00041,0.00072,0.00267,0.00467,0.00607\) percent.
The last two responses are not significantly resolved. This four-noise
single-configuration study is not an ensemble mass-systematic estimate.
Its files are under
\nolinkurl{quark_ringed_mass_pilot_v1/runs/critical_mass_20260918_cuonexuz/analysis/}.
Some 32-noise mass-scan products exist, but a complete four-mass reviewed
report is still to be confirmed. The ten-configuration error bars below
are conditional on the central mass and do not include this mass systematic.

\subsubsection{Completed ten-configuration products and reproducible operations}
\label{sec:ringed-tenconf-results}

The gauge configurations are \(10000,10100,\ldots,10900\), with
\(N_\eta=32\) on each. Noise generation uses
\texttt{SeedSequence([9182026,conf,noise\_id])}; different configurations
do not reuse identical noise fields. The six checkpoints have
\(\theta=(0.05,0.1,0.2,0.5,0.8,1.0)\), step size \(0.01\), and
accumulated steps \((5,10,20,50,80,100)\). All flows share the original
\((50,10)\) bare-fit bootstrap-index array, with seed 1115. The two
existing bare-fit methods are retained separately: six-flow Sumratio
differences with AIC, and four-flow direct ratio AIC marked as an
unweighted diagnostic. No fit window or fallback flag is changed.

\begin{table}[htbp]
 \centering
 \begin{tabular}{rrr}
 \toprule
 \(\theta\) & \(\mathring Z_\chi\) & Paired bootstrap error \\
 \midrule
 0.05 & 6.7679838065 & 0.0003478494 \\
 0.10 & 2.7791305139 & 0.0001365254 \\
 0.20 & 1.5885047901 & 0.0000585582 \\
 0.50 & 1.2514982449 & 0.0000741719 \\
 0.80 & 1.2505405670 & 0.0001376585 \\
 1.00 & 1.2665887094 & 0.0002108594 \\
 \bottomrule
 \end{tabular}
 \caption{Ten-configuration ringed factors at the published critical valence
 mass, on the fixed massive sea. Errors use the original 50 gauge-bootstrap
 draws of the 32-noise means; mass-prescription and other systematics are
 not included.}
 \label{tab:ringed-tenconf}
\end{table}

The complete sample passes the saved noise/aggregate identity and hash
checks, flowed-link SU(3) checks, and the true full-MAT residual gate
\(10^{-9}\). Its maximum residual is \(7.52886\times10^{-13}\).
All six imaginary kinetic means pass the declared five-standard-error
consistency check. These numerical gates do not establish complete physics
acceptance.

The initial solve for configuration 10800 stalled. Its isolated recovery
used single-precision MG preconditioning with half-precision null-vector
storage, GCR Krylov dimension 16, and reliable-update parameter 0.1.
The physics, noise IDs, tolerance, and true-residual gate were unchanged.
All 32 noises were recomputed; all passed on the initial solve without
defect correction. The old and new accepted noise-0 kinetic values differ
by at most \(1.70589\times10^{-10}\) relatively, and the recovery's
same-noise free-field comparison differs by \(1.44009\times10^{-5}\),
below the prescribed \(5\times10^{-4}\) threshold. GPU job 3707765 and
CPU analysis/plotting job 3709104 both completed with exit code zero.

\paragraph{Locations and code-to-formula map.}
All relative paths in this paragraph start at
\nolinkurl{/public/group/lqcd/donghx/Gradient_Flow_GluonPDF/}.
The completed output root is
\begin{center}
 \nolinkurl{quark_ringed_ensemble_v1/}\par
 \nolinkurl{results_10conf_after_Zchi_20260920/}.
\end{center}
The original campaign is
\nolinkurl{quark_ringed_ensemble_v1/runs/critical_10conf_20260920_rz65f_im/}.
Its \texttt{reused/conf10000} and eight complete new configuration
directories supply nine inputs. Configuration 10800 comes from
\nolinkurl{quark_ringed_recovery_v7/measurement/kinetic_conf.npz};
the corresponding \texttt{solver\_recovery\_review.json} binds the
replacement path, unchanged physics, and numerical-equivalence evidence.
It does not overwrite the old failed data or serve as a formal physics review.

The estimator is \texttt{quark\_ringed\_kinetic\_v1/core.py}.
The actual recovery measurement code is the immutable copy under
\nolinkurl{quark_ringed_recovery_v7/code/quark_ringed_kinetic_v1/}.
In \nolinkurl{quark_ringed_ensemble_v1/pipeline.py}, \texttt{audit\_raw}
checks the measurements, \texttt{paired} implements
Eqs.~\eqref{eq:ringed-paired-bootstrap}--\eqref{eq:ringed-product-error},
and \texttt{analyze} preserves the original fit products and masks.
The executed CPU snapshot is
\nolinkurl{quark_ringed_recovery_v7/analysis_code/}.

Under the output root, each of
\nolinkurl{analysis/sixflow_sumdiff/} and
\nolinkurl{analysis/fourflow_ratio/} contains
\texttt{paired\_matrix\_elements.npz} and \texttt{manifest.json}.
The former uses schema \texttt{quark\_ringed\_paired\_10conf\_v1}.
For six flows, kinetic noise has axes \((6,10,32)\) for
\((\theta,\mathrm{conf},\mathrm{noise})\);
\texttt{Zchi\_bootstrap} has shape \((6,50)\);
matrix centers/errors have shape \((6,3,21)\), and matrix replicas have
shape \((6,50,3,21)\), ordered as \((\theta,b,P_z,z)\).
The momenta are \(n_z=3,4,5\) and separations are \(z/a=0,\ldots,20\).
The four-flow method replaces the first dimension by four.
The original conf list, indices, bare samples, complex kinetic data, quality
masks, and source hashes are retained.

To reproduce the plots from existing products, run the following in the
project root, replacing \texttt{ANALYSIS} by the output root's
\texttt{analysis} directory and choosing a new \texttt{NEW\_FIGURES}
directory. No additional GPU inversion is needed:
\begin{quote}\small\ttfamily
/public/home/donghx/AnaConda/bin/python3 -B -m\par
quark\_ringed\_ensemble\_v1.plot\_after\_zchi\par
ANALYSIS --output NEW\_FIGURES
\end{quote}
These are wrapped lines of one shell command. Set
\texttt{OPENBLAS\_NUM\_THREADS=1} and a writable
\texttt{MPLCONFIGDIR} when running on a CPU node.
Reanalysis uses \texttt{pipeline analyze ORIGINAL\_RUN} with
\texttt{--recovery-review REVIEW --output NEW\_ANALYSIS} from the frozen
CPU snapshot. Existing output directories must not be overwritten.
The complete absolute paths, original bare-input filenames, and exact
commands are recorded in the output README and Section 14 of
\nolinkurl{/public/home/donghx/.codex/Markdown/gradient_flow.md}.

\begin{figure}[htbp]
 \centering
 \includegraphics[width=\linewidth,page=1]{figures/ringed_10conf_after_zchi.pdf}
 \caption{After ringed normalization for the six-flow Sumratio-difference
 AIC input. Each panel has a different momentum; flow times are displaced
 horizontally and points are not joined. Open points retain the original
 fit-fallback status. The small-flow values are shown with their full error
 range, including the large normalization at \(\theta=0.05\). No SFTE
 factor is applied.}
 \label{fig:ringed-tenconf-all}
\end{figure}
\begin{figure}[htbp]
 \centering
 \includegraphics[width=\linewidth,page=1]{figures/ringed_10conf_subset.pdf}
 \caption{The same paired products, displaying only
 \(\theta=0.2,0.5,0.8,1.0\) for readability. This is a display selection,
 not a refit or an accepted SFTE window. The three momenta are approximately
 1.50, 2.00, and 2.50 GeV at \(a=0.0775\) fm.}
 \label{fig:ringed-tenconf-subset}
\end{figure}

\paragraph{Interface to matching and outstanding checks.}
The after-ringed plots use \(\mathring Z_\chi h^{\rm GF}\) only.
The extra saved \texttt{sfte\_*} diagnostics multiply the paired ringed
samples by \(\mathcal C_\psi^{-1}e^{-\delta m|z|}\) at
\(\mu=2\) GeV, \(\alpha_s=0.30\); they do not multiply by
\(\mathring Z_\chi\) again. Their \(z=0\) entries are masked, since the
nonlocal small-flow expression is not a local-current matching formula.
Ringed \(z=0\) entries remain finite. The negative \(\delta m\) convention
of Eq.~\eqref{eq:delta-m-one-loop} is unchanged.

For these six flow times, \(\rflow/a\) is approximately
\(0.632,0.894,1.265,2,2.530,2.828\). The first two are below one lattice
spacing. Neither normalization nor numerical convergence establishes a
window \(a\ll\rflow\ll|z|\) or \(\sqrt{\tf}P_z\ll1\).
Next steps are to review action/artifact and boundary-improvement provenance,
complete the paired mass-response study, assess noise and gauge-autocorrelation
stability, and retain the original fit and six-direction-estimator caveats.
Flow-window, continuum, scale, and perturbative-order systematics remain
separate. Metadata record complete statistical pairing but
\texttt{formal\_physics\_acceptance=False}, and all
\texttt{formal\_result\_valid} entries remain false. Fifty replicas do not
increase ten configurations into fifty independent configurations.

\subsection{Status and derivation of the available matching kernels}
\label{sec:matching-status}

Two perturbative conversions must be distinguished.  Gradient-flow matching
converts a finite-flow operator into a renormalized equal-time quasi or pseudo
operator.  LaMET (or pseudo-ITD) matching then converts that operator into a
light-cone PDF.  The two kernels have different schemes and must not be
multiplied twice when one combined coefficient has already been used.

For quarks, one-loop kernels exist for vector (unpolarized), axial
(helicity), and tensor (transversity) quasi-PDFs in the nonsinglet sector;
coordinate-space ratio/hybrid versions are also available.  Transversity has
no leading-twist gluon mixing.  Singlet unpolarized and helicity matching are
matrix-valued in the quark--gluon sector, so connected $u$ or $d$ contractions
alone do not constitute the complete singlet result.

For gluons, one-loop ratio, hybrid, and pseudo-distribution kernels exist for
the unpolarized channel \cite{KhanEtAl2021UnpolarizedGluon,BalitskyMorrisRadyushkin2022,ChenEtAl2025GluonPDF}.
The helicity gluon pseudo-ITD and its invariant-amplitude matching are given
in Refs.~\cite{EgererEtAl2022GluonHelicity,BalitskyMorrisRadyushkin2022}.
These are quasi/pseudo-to-light-cone kernels and do not replace the
flowed-operator conversion.

The 4D gradient-flow-to-$\MSbar$ conversion needed before those kernels was
derived at one loop by Brambilla and Wang \cite{BrambillaWang2024}.  Introduce
an auxiliary adjoint field $h_v$ for the straight Wilson line and decompose
the field strength relative to its direction $v$:
\begin{equation}
 J_{\parallel\perp}^{\mu\nu}=gF_{\parallel\perp}^{\mu\nu}h_v,
 \qquad J_{\perp\perp}^{\mu\nu}=gF_{\perp\perp}^{\mu\nu}h_v .
\end{equation}
The small-flow expansion of each local current is
\begin{equation}
J_i^{\rm GF}(\tf)=c_i(\tf,\mu)J_i^{\MSbar}(\mu)+O(\tf),
 \qquad i=\parallel\perp,\perp\perp .
\end{equation}
This should not be read as a complete channel-by-channel catalogue of
flow-to-quasi kernels.  The published result is a one-loop small-flow-time
conversion for the four-dimensional flowed Wilson-line basis (and the quark
coefficient is Dirac-structure independent at this order).  The vector,
axial, tensor, and gluon LaMET kernels quoted above are a separate, subsequent
matching step.  A complete all-order GF-to-quasi result for every channel,
including the singlet quark--gluon mixing matrix and a nonperturbative line
mass, is not currently available.  Three-dimensional flow and finite-flow
bare correlators therefore cannot be assigned these coefficients without a
separate matching calculation.
At one loop, in the convention of Ref.~\cite{BrambillaWang2024},
\begin{equation}
 c_{\parallel\perp}=1+O(\alpha_s^2),\qquad
 c_{\perp\perp}=1+\frac{\alpha_s C_A}{4\pi}
 \ln(2\mu^2\tf e^{\gamma_E})+O(\alpha_s^2).
\end{equation}
The Wilson-line bilocal receives the corresponding products
\begin{equation}
 \mathcal C_{\parallel\perp}=c_{\parallel\perp}^{2},\qquad
 \mathcal C_{\perp\perp}=c_{\perp\perp}^{2},
\end{equation}
and a line-mass factor
\begin{equation}
 \exp[\delta m(\tf)|z|],\qquad
 \delta m(\tf)=-\frac{\alpha_s C_A}{4\pi}
 \frac{\sqrt{2\pi}}{\sqrt{\tf}}+O(\alpha_s^2).
\end{equation}
The exact sign follows the Wilson-line orientation and the Euclidean
convention.  The derivation computes flowed and unflowed partonic diagrams
with the same IR regulator; the unflowed integrals are scaleless in dimensional
regularization, so subtracting their IR poles from the renormalized flowed
amplitudes leaves the matching coefficient.  Multiplicative renormalization
of the auxiliary-field currents then promotes the local result to the
nonlocal Wilson line.

For the present gluon operators, this conversion must be performed on the
separately stored temporal and spatial pieces before forming the channel.  In
the code convention (transverse sums included) the relevant combinations are
\begin{align}
 \mathcal O_U^{\rm GF}&=T_U^{\rm GF}-S_U^{\rm GF},
 &T_U^{\rm GF}&=M_{tx;tx}^{\rm GF}+M_{ty;ty}^{\rm GF},
 &S_U^{\rm GF}&=2M_{xy;xy}^{\rm GF},\\
 \mathcal O_{\Delta g}^{\rm GF}&=T_H^{\rm GF}-S_H^{\rm GF},
 &T_H^{\rm GF}&=M_{tx;\widetilde{tx}}^{\rm GF}+M_{ty;\widetilde{ty}}^{\rm GF},
 &S_H^{\rm GF}&=2M_{xy;\widetilde{xy}}^{\rm GF}.
\end{align}
The tensor-basis assignment is fixed by the Wilson-line direction and the
Minkowski--Euclidean dictionary; it must not be inferred from the letters
`t', `x', and `y' alone.  Here the Wilson line is along the spatial $z$
direction.  Consequently $F_{tx}$, $F_{ty}$, and $F_{xy}$ are all
$\perp\perp$ components relative to that line.  The unpolarized conversion
is therefore applied to the two separately stored terms as
\begin{equation}
 \mathcal O_U^{\overline{\rm MS}}
 =e^{-\delta m|z|}
 \left[
 \frac{T_U^{\rm GF}}{\mathcal C_{\perp\perp}}
 -\frac{S_U^{\rm GF}}{\mathcal C_{\perp\perp}}
 \right]+O(\tf).
 \label{eq:gluon-unpolarized-flow-matching}
\end{equation}
For helicity, dualization exchanges the two endpoint tensor sectors.  Both
pieces therefore have the same coefficient
\begin{equation}
 \mathcal C_{\Delta g}
 =c_{\parallel\perp}c_{\perp\perp}
 =\sqrt{\mathcal C_{\parallel\perp}\mathcal C_{\perp\perp}}
 =1+\frac{\alpha_s C_A}{4\pi}L_t+O(\alpha_s^2),
 \label{eq:gluon-helicity-flow-coefficient}
\end{equation}
and
\begin{equation}
 \mathcal O_{\Delta g}^{E,\overline{\rm MS},\,\mathrm{odd}}
 =\frac{e^{-\delta m|z|}}{\mathcal C_{\Delta g}}
 \left\{[T_H-S_H](+z)-[T_H-S_H](-z)\right\}+O(\tf).
 \label{eq:gluon-helicity-flow-matching}
\end{equation}
Thus a single common scalar factor for `Mtiti' and `Mijij' is the
unpolarized prescription for this spatial Wilson-line geometry.  It is also
appropriate for helicity only after the endpoint-factorized
$\mathcal C_{\Delta g}$ derivation.  The latter is a
natural one-loop consequence of the published local current coefficients,
not yet an independently published complete helicity partonic matching
calculation.  The current code has produced finite-flow matched operator
diagnostics and an exploratory unpolarized quasi matching; these outputs are
not yet a complete renormalized physical PDF.

The operational chain is
\begin{center}
\(\text{flowed Clover}
 \xrightarrow{\;\mathcal C_i(\tf,\mu),\,e^{-\delta m|z|}\;}
 \text{renormalized quasi/pseudo matrix element}\)
\par\vspace{2pt}
\(\xrightarrow{\;C^{\rm LaMET}_{q,g}\;}\quad
 \text{light-cone PDF}.\)
\end{center}
The same logic applies to quark vector, axial, and tensor channels, with the
ringed-fermion normalization included for flowed quarks.  A physical PDF also
requires the appropriate singlet mixing, target-mass and power corrections,
and controlled continuum and flow-time limits.

\subsection{Small-flow window and implementation in the present pipeline}
\label{sec:sfte-implementation}

The asymptotic coefficient is useful only where the smoothing radius is both
resolved by the lattice and shorter than the physical scales entering the
bilocal:
\begin{equation}
  a\ll\sqrt{8\tf}\ll|z|,
  \qquad
  \sqrt{\tf}P_z\ll1,
  \qquad
  \tf\Lambda_{\rm QCD}^2\ll1.
  \label{eq:sfte-window}
\end{equation}
Each product should therefore store the dimensionless diagnostics
\begin{equation}
  \frac{a}{\sqrt{8\tf}},\qquad
  \frac{\sqrt{8\tf}}{|z|},\qquad
  \sqrt{\tf}P_z,\qquad
  \tf\Lambda_{\rm QCD}^2,
  \label{eq:sfte-diagnostics}
\end{equation}
together with a declared validity mask.  Window thresholds are global
systematic choices; selecting them separately at every \((z,P_z)\) biases the
analysis.

The present generalized vertex already uses the required common gauge history,
\begin{equation}
  J(\tf;t_J,\tau_a,\tau_b)
  =[K(\tf)V_{\tau_a}]^\dagger
  O[B(\tf),t_J]
  [K(\tf)V_{\tau_b}],
\end{equation}
and the Wilson line is constructed from flowed thin links without HYP
pre-smearing.  The source and sink nucleon interpolators and \(C_2\) remain
unflowed, which is consistent because the matched insertion is the bilocal
current. The required factor is \(\mathring Z_\chi(\tf)\), now measured in
the ten-configuration diagnostic of Section~\ref{sec:ringed-tenconf-results}.
If this quantity is
an ensemble/flow-time scalar, it can be multiplied into the completed C3 or
the fitted C3/C2 matrix element.  Existing flowed vertices and C3 therefore
need not be regenerated, and no new operator-flow Dirac inversion is required.
The vacuum expectation value in Eq.~\eqref{eq:ringed-zchi-nonperturbative} may
nevertheless require new stochastic estimators or inversions unless the
accuracy of its reconstruction from the available all-to-all basis has been
demonstrated.

The independent post-processing implementation is
\begin{center}
  \nolinkurl{/public/home/donghx/Gradient_Flow/Blending_L32x96/small_flow_time/}.
\end{center}
It refuses placeholder values of \(\mathring Z_\chi\), propagates the existing
bootstrap samples, records the one-loop coefficient and sign convention, and
saves the masks in Eq.~\eqref{eq:sfte-diagnostics}.  Its output is an
\(\MSbar\) coordinate-space bilocal diagnostic, not a matched PDF.

At \(\tf/a^2=0.5\), \(\sqrt{8\tf}/a=2\).  The points \(|z|/a\le2\) therefore
do not satisfy the simple asymptotic separation condition.  Moreover,
\(\sqrt{\tf}P_z\simeq0.42,0.56,0.69\) for momentum integers 3, 4, and 5 on
the length-32 lattice.  The largest momenta are marginal at this flow time.
The multi-flow checkpoints provide a scan, but the three ensembles must be
interpolated to common physical \(\tf\) or \(\sqrt{8\tf}\); a common value of
\(\tf/a^2\) is not a common physical flow time.

Finally, the exponential in Eq.~\eqref{eq:explicit-sfte-matrix-element}
contains the one-loop flowed line-mass subtraction.  It must not be multiplied
on top of an overlapping nonperturbative hybrid or ratio subtraction.  If a
nonperturbative line mass is preferred, the residual finite conversion must be
derived in that same scheme.  After SFTE and the controlled limit analysis,
one applies the channel-specific LaMET kernel.  The first clean target is the
nonsinglet \(u-d\) combination; individual connected \(u,d\) results do not
include disconnected and singlet--gluon mixing effects.

\section{Lattice evolution algorithm and its PyQUDA realization}
\label{sec:pyquda-gflow}

\subsection{Covariant lattice Laplacian}

For a flowed link field \(V_\mu(\tf,x)\), the nearest-neighbor
four-dimensional Laplacian is
\begin{equation}
  \Delta_V\phi(x)
  =
  \frac{1}{a^2}\sum_{\mu=1}^{4}
  \left[
    V_\mu(x)\phi(x+\hat\mu)
    +V_\mu^\dagger(x-\hat\mu)\phi(x-\hat\mu)
    -2\phi(x)
  \right].
  \label{eq:lattice-laplacian}
\end{equation}
The production QUDA call works in lattice units.  Its
\texttt{ApplyLaplace} arguments are \texttt{nDim=4}, hopping coefficient
\(a_{\rm hop}=1\), and diagonal coefficient \(b=-8\).  Thus its dimensionless
operator is
\begin{equation}
  (a^2\Delta_V\phi)(x)
  =\sum_{\mu=1}^{4}
  \left[
    V_\mu(x)\phi(x+\hat\mu)
    +V_\mu^\dagger(x-\hat\mu)\phi(x-\hat\mu)
  \right]-8\phi(x).
  \label{eq:quda-lattice-laplacian}
\end{equation}
The four directions are communicated across MPI boundaries before each
application.  The operator is diagonal in Dirac spin and acts nontrivially
only on position and color.  On unit links its eigenvalue is
\(-a^2\widehat p^{\,2}\), reproducing
Eq.~\eqref{eq:lattice-free-kernel}.

\subsection{The coupled third-order Runge--Kutta step}

The current QUDA source implements the forward gauge--fermion flow algorithm
of Appendix~D of Ref.~\cite{Luscher2013}.  Let \(V_0\) and \(\chi_0\) denote
the gauge and fermion fields at the beginning of one step. Here
\(\tau=\tf/a^2\), \(\epsilon=\Delta\tau\) and
\(L[V]=a^2\Delta_V\) are dimensionless; the following stages use
\(\partial_\tau\chi=L[V]\chi\). Define
\begin{equation}
  k_0=L[V_0]\chi_0.
\end{equation}
The first fermion and gauge stages are
\begin{align}
  \chi_1
  &=\chi_0+\frac{\epsilon}{4}k_0,
  \\
  V_1
  &=\exp\!\left\{
      \frac{\epsilon}{4}\mathcal Z[V_0]
    \right\}V_0,
  \label{eq:quda-rk-stage1}
\end{align}
where \(\mathcal Z[V]\) is the anti-Hermitian Wilson gauge-force generator
constructed from the staples.  Next set
\begin{equation}
  k_1=L[V_1]\chi_1
\end{equation}
and update
\begin{align}
  \chi_2
  &=\chi_0+\epsilon\left(
      \frac{8}{9}k_1-\frac{2}{9}k_0
    \right),
  \\
  V_2
  &=\exp\!\left\{
      \epsilon\left(
        \frac{8}{9}\mathcal Z[V_1]
        -\frac{17}{36}\mathcal Z[V_0]
      \right)
    \right\}V_1.
  \label{eq:quda-rk-stage2}
\end{align}
The final stage uses
\begin{equation}
  k_2=L[V_2]\chi_2
\end{equation}
and gives
\begin{align}
  \chi_{\rm new}
  &=\chi_1+\frac{3\epsilon}{4}k_2,
  \\
  V_{\rm new}
  &=\exp\!\left\{
      \epsilon\left(
        \frac{3}{4}\mathcal Z[V_2]
        -\frac{8}{9}\mathcal Z[V_1]
        +\frac{17}{36}\mathcal Z[V_0]
      \right)
    \right\}V_2.
  \label{eq:quda-rk-stage3}
\end{align}
The definition of \(\mathcal Z\) absorbs the sign and Lie-algebra convention
used in the link exponential.  Equations~\eqref{eq:quda-rk-stage1}--
\eqref{eq:quda-rk-stage3} expose the essential point: every fermion Laplacian
is evaluated on the matching intermediate gauge field.  Flowing the gauge
first and then evolving the fermion through a fixed final field would be a
different discretization.

The QUDA temporary arrays correspond directly to these equations:
\begin{center}
\begin{tabular}{p{0.23\textwidth}p{0.67\textwidth}}
  \toprule
  QUDA object & Mathematical role during one step \\
  \midrule
  \texttt{f\_temp3}
  & Current solution \(\chi_0\), then \(\chi_1\), and finally
    \(\chi_{\rm new}\) \\
  \texttt{f\_temp0}
  & First derivative \(k_0=L[V_0]\chi_0\) \\
  \texttt{f\_temp1}
  & First stage \(\chi_1\), then reused for
    \(k_1=L[V_1]\chi_1\) \\
  \texttt{f\_temp2}
  & Second stage \(\chi_2\), then reused for
    \(k_2=L[V_2]\chi_2\) \\
  \texttt{f\_temp4}
  & Scratch output of \texttt{ApplyLaplace} \\
  \texttt{gin}, \texttt{gout}, \texttt{gaugeTemp}
  & Input, output, and stored-force gauge fields for the three link stages \\
  \bottomrule
\end{tabular}
\end{center}

\subsection{The numerical kernel is a product of linear step maps}

For a fixed gauge trajectory, one RK step is linear in the input fermion
field.  Denote this linear map by
\begin{equation}
  \chi_{n+1}=\mathcal R_\epsilon[V_n]\chi_n.
\end{equation}
After \(N=\tau/\epsilon=\tf/(a^2\epsilon)\) steps, QUDA has applied
\begin{equation}
  K_{\epsilon,N}[V_0]
  =
  \mathcal R_\epsilon[V_{N-1}]\cdots
  \mathcal R_\epsilon[V_1]\mathcal R_\epsilon[V_0],
  \qquad
  \chi_N=K_{\epsilon,N}[V_0]\chi_0.
  \label{eq:quda-discrete-kernel-product}
\end{equation}
This converges to the flow-time-ordered heat kernel as
\(\epsilon\to0\) at fixed \(\tf\).  QUDA never forms the matrix on the
right-hand side of Eq.~\eqref{eq:quda-discrete-kernel-product}; it evaluates
its action on supplied right-hand sides.  Constructing every column would
require one initial delta field for every space--time and color index.

The LapH eigenvectors therefore do not enter the differential equation and do
not diagonalize \(\Delta[V(\tf)]\).  They enter only as the right-hand sides
needed by distillation.  For an eigenvector on source time \(\tau\),
\begin{equation}
  \eta_{\tau n}(t,\bm x,c)
  =\delta_{t\tau}V_{\tau n}(\bm x,c),
  \qquad
  \Phi_{\tau n}^{(\tf)}
  =K_{\epsilon,N}\eta_{\tau n}.
  \label{eq:quda-k-times-v}
\end{equation}
Thus the implemented calculation obtains columns of \(KV\), not columns of
the full coordinate-space \(K\).

\subsection{PyQUDA binding and parameter flow}

The Python binding accepts two arrays of spinor pointers and forwards them to
the QUDA C interface:
\begin{equation}
  \texttt{performGFlowQuda}
  (h_{\rm out},h_{\rm in},p_{\rm inv},p_{\rm smear},p_{\rm obs},N_{\rm rhs}).
  \label{eq:pyquda-gflow-signature}
\end{equation}
The binding checks that the leading pointer-array dimensions of
\(h_{\rm in}\) and \(h_{\rm out}\) both equal \(N_{\rm rhs}\).  In this API,
\texttt{QudaInvertParam} supplies the color-spinor layout, precision, field
location, and related metadata required to construct QUDA spinor objects.
No linear solver and no Dirac inversion is called by
\texttt{performGFlowQuda}.

The relevant smear parameters are
\begin{center}
\begin{tabular}{p{0.30\textwidth}p{0.60\textwidth}}
  \toprule
  Parameter & Meaning \\
  \midrule
  \texttt{smear\_type}
  & \texttt{QUDA\_GAUGE\_SMEAR\_WILSON\_FLOW} \\
  \texttt{n\_steps}
  & Number of coupled gauge--fermion RK steps in this call \\
  \texttt{epsilon}
  & Step size in lattice flow-time units \\
  \texttt{t0}
  & Flow time attached to the beginning of this segment \\
  \texttt{restart}
  & Reuse the resident flowed gauge rather than restart from the thin links \\
  \texttt{meas\_interval}
  & Interval for plaquette and spinor-norm diagnostics \\
  \bottomrule
\end{tabular}
\end{center}
When \texttt{restart=false}, QUDA copies the resident thin-link field into its
flowed-gauge workspace.  When \texttt{restart=true}, it requires the previous
flowed gauge to remain resident.  In both cases the supplied input spinor is
copied to the device, evolved, and copied to the supplied output pointer.

\subsection{Tensor layout and mode packing in the present code}

For a batch of \(N_b\) LapH modes, the host lexicographic source has shape
\begin{equation}
  (N_{\rm spinor},L_t^{\rm local},L_z,L_y,L_x,N_s,N_c),
  \qquad
  N_{\rm spinor}=\left\lceil\frac{N_b}{4}\right\rceil,
  \label{eq:gflow-host-layout}
\end{equation}
with \(N_s=4\) and \(N_c=3\).  Mode \(n\) is inserted as
\begin{equation}
  \texttt{source}
  \left[\left\lfloor n/4\right\rfloor,
        \tau,\bm x,n\bmod4,c\right]
  =V_{\tau n}(\bm x,c).
  \label{eq:gflow-mode-pack}
\end{equation}
PyQUDA converts this array to its even--odd local layout before passing the
pointer array to QUDA.  Because the covariant Laplacian is spin diagonal, the
four spin components evolve independently.  They are used here as four
storage lanes for four LapH modes; this does not introduce physical spin
mixing.  After the call, the code converts back to lexicographic order,
unpacks the four lanes, and saves
\begin{equation}
  \Phi_{\tau n}^{(\tf)}(t_J,\bm x,c)
  \quad\text{with shape}\quad
  (N_v,L_z,L_y,L_x,N_c)=(400,32,32,32,3)
  \label{eq:gflow-saved-basis-shape}
\end{equation}
for each retained pair \((\tau,t_J)\).

The temporal anti-periodic boundary condition is implemented by applying a
center phase to the temporal links before loading the gauge into QUDA.  This
phase affects temporal hopping in Eq.~\eqref{eq:quda-lattice-laplacian}; it is
removed before the flowed gauge links are written for Wilson-line use.

\subsection{Kernel dimension and the number of flow evolutions}

There is one mathematical kernel for a fixed gauge trajectory and flow time,
but it must be applied to every independent initial vector required by the
estimator.  Suppressing spin, define
\begin{equation}
  N_{4c}=N_tN_xN_yN_zN_c
  =96\times32^3\times3
  =9\,437\,184.
\end{equation}
The full four-dimensional color-coordinate kernel would be the dense matrix
\begin{equation}
  K(\tf)\in\mathbb C^{N_{4c}\times N_{4c}},
  \qquad
  K_{\alpha a x;\beta b y}(\tf)
  =\delta_{\alpha\beta}K_{a x;b y}(\tf).
  \label{eq:full-kernel-dimension}
\end{equation}
It contains approximately \(8.91\times10^{13}\) complex entries.  Dense
storage would require about \(0.63\,\mathrm{PiB}\) in complex64 or
\(1.27\,\mathrm{PiB}\) in complex128, even after exploiting its spin-diagonal
form.  The implementation must therefore remain matrix-free.

For one kernel source time \(\tau\), the embedded stochastic-distillation
basis and its flowed image have dimensions
\begin{equation}
  \mathcal V_\tau\in\mathbb C^{N_{4c}\times400},
  \qquad
  \Phi_\tau^{(\tf)}=K(\tf)\mathcal V_\tau
  \in\mathbb C^{N_{4c}\times400}.
  \label{eq:k-v-dimension}
\end{equation}
Equivalently, the complete flowed array has shape
\((400,96,32,32,32,3)\).  It would occupy about \(28.1\,\mathrm{GiB}\) in
complex64 or \(56.3\,\mathrm{GiB}\) in complex128 for one \(\tau\).  The code
therefore evolves mode batches and retains only the requested output-time
slices.

With the default \(N_b=32\), four modes are packed into each Wilson spinor,
so a full batch passes eight spinors to \texttt{performGFlowQuda}.  For every
source time \(\tau\), the number of calls at one flow time is
\begin{equation}
  N_{\rm call}(\tau)
  =\left\lceil\frac{400}{32}\right\rceil
  =13.
  \label{eq:gflow-call-count}
\end{equation}
The first twelve calls contain 32 modes and the last contains 16 modes, or
four Wilson spinors.  The four GPUs implement one time-decomposed lattice
evolution collectively; they do not turn the 400 modes into four independent
100-mode solves.

Thus ``one kernel'' means one linear map shared by all right-hand sides, not
one scalar-vector evolution from which the other 399 outputs can be inferred.
A sufficiently large multi-right-hand-side allocation could set
\texttt{FLOW\_MODE\_BATCH=400} and make one QUDA call per \(\tau\), but that
call would still evolve 400 independent initial modes and would require 100
Wilson spinor fields.  In the present single-flow path,
\texttt{restart=false} is used for every batch.  Consequently each batch
repeats the same thin-link gauge-flow trajectory; only the right-hand sides
inside that batch share it.  The multi-flow-time path reuses successive
flow-time segments with \texttt{restart=true} within a batch, but starts again
from the thin links for the next mode batch.  This is mathematically
consistent and is also an identifiable performance limitation.

\subsection{Restarted flow-time checkpoints}

Suppose \(0<\tf^{(1)}<\tf^{(2)}\).  The first segment uses the thin-link gauge,
the unflowed source \(\eta\), and \texttt{restart=false}:
\begin{equation}
  \Phi^{(1)}=K(\tf^{(1)},0)\eta.
\end{equation}
The next segment passes \(\Phi^{(1)}\) as the input spinor, keeps the flowed
gauge resident, sets \texttt{restart=true}, and advances only
\(\tf^{(2)}-\tf^{(1)}\):
\begin{equation}
  \Phi^{(2)}
  =K(\tf^{(2)},\tf^{(1)})\Phi^{(1)}
  =K(\tf^{(2)},0)\eta.
  \label{eq:gflow-restart-composition}
\end{equation}
The production code alternates two output work fields and publishes each
checkpoint under its own flow-time tag.  This avoids repeating the trajectory
from zero for every requested \(\tf\).  Equality with independent direct
evolution is a required numerical regression; it follows analytically from
Eq.~\eqref{eq:gflow-restart-composition} only if the same resident gauge,
input spinor, step size, and boundary convention are retained.

\subsection{What is solved and what is not}

The operation above is an explicit initial-value evolution.  It is not a
Krylov solve, does not invert \(\Delta\), and does not invert the Wilson or
Clover Dirac matrix.  Its output for each initial vector is
\(K\eta\).  The expensive Dirac inversions are already contained in the
perambulators.  The flowed current is assembled later as
\begin{equation}
  J^{(\tf)}
  =
  (KV_{\tau_a})^\dagger
  O[B(\tf)]
  (KV_{\tau_b}),
\end{equation}
which is then contracted with the ordinary all-time perambulators.

The minimum implementation checks are: the unit-link damping
\(e^{-\tf\widehat p^2}\), gauge covariance, zero-flow recovery, equality of
direct and restarted trajectories, stability under
\(\epsilon\to\epsilon/2\), and the discarded temporal norm for a finite
storage window.

\section{Distillation representation}

\subsection{LapH eigenvectors and perambulators}

On each time slice \(\tau\), let \(V(\tau)\) contain \(N_v\) orthonormal
eigenvectors of a gauge-covariant three-dimensional Laplacian:
\begin{equation}
  -\nabla_{\rm 3D}^2(\tau)V_n(\tau)
  =
  \lambda_n(\tau)V_n(\tau),
  \qquad
  V^\dagger(\tau)V(\tau)=\mathbf 1.
  \label{eq:laph}
\end{equation}
The distillation projector is
\begin{equation}
  \Box(\tau)=V(\tau)V^\dagger(\tau).
\end{equation}
For flavor \(f\), the conventional perambulator is
\begin{equation}
  \tau_f(t,\tau)
  =
  V^\dagger(t)\,D_f^{-1}(t,\tau)\,V(\tau),
  \label{eq:perambulator}
\end{equation}
with spin indices suppressed. This object contains the costly Dirac
inversions. The light perambulators use anti-periodic temporal boundary
conditions and cover all source times.

The stored compressed basis must satisfy, for example,
\begin{equation}
  \max_{m,n}\left|
  (V^\dagger V)_{mn}-\delta_{mn}
  \right|.
\end{equation}
Compression and its weight \(\omega_{mn}\) affect the estimator, not the
continuum flow equation.

The 400 vectors used here are not the lowest 400 Laplacian eigenvectors.  The
stored ordering and the perambulator source range select
\begin{equation}
  400=100_{\rm exact}+100_{\rm compressed}+200_{\rm high\ noise}.
  \label{eq:stochastic-basis-content}
\end{equation}
The middle part consists of orthonormal random combinations from the five
spectral groups
\([100,200,400,600,800]\), compressed to 20 vectors per group.  The high-mode
vectors are generated after orthogonalization against the explicitly
constructed modes.  Although the file tag records that 400 high-noise vectors
were generated, the current \(N_v=400\) source range takes only the first 200
of them after the 100 exact and 100 compressed vectors.  Accordingly the
two-index weight used by the correlator has \(N_{\rm noise}=200\).

Schmidt orthogonalization targets
\begin{equation}
  V^\dagger V=\mathbf 1_{400},
\end{equation}
which must still be checked on the actual input files; this note does not
claim a complete Gram audit of all files. It does not establish the sample-wise completeness relation
\(VV^\dagger=\mathbf 1_{32^3N_c}\).  If a spectral sector has full dimension
\(D_r\) and contains \(n_r\) sampled orthonormal vectors, its first-moment
completeness estimator is schematically
\begin{equation}
  \widehat{\mathbf 1}_r
  =\frac{D_r}{n_r}\sum_{a=1}^{n_r}v_{ra}v_{ra}^\dagger,
  \qquad
  \mathbb E_\eta[\widehat{\mathbf 1}_r]=\mathbf 1_r,
  \label{eq:stochastic-completeness}
\end{equation}
under the assumed noise ensemble.  Repeated indices require joint inclusion
weights rather than a product of independent one-index factors.  The current
\(\omega_{mn}\) implements the corresponding two-index sector-size and
same-index corrections and is checked against the reference connected-diagram
code. This describes the historical unified-table scheme A. The subsequent
scheme B distinguishes equal and unequal original source times; the
definitions and validation limits are in Section~\ref{sec:blending-oct15}.

It is therefore appropriate to call the coordinate-space object a
\emph{stochastic reconstruction of the conventional propagator}, rather than
an exact equality on each configuration and noise realization.  For the
two-index convention used by the current vertex, it may be represented as
\begin{equation}
  \widehat S_\omega(x,y)
  =\sum_{m,n}V_m(x)
    \bigl[\omega_{mn}\tau_{mn}\bigr]V_n^\dagger(y)
  =V\bigl(\omega\odot\tau\bigr)V^\dagger,
  \label{eq:stochastic-coordinate-propagator}
\end{equation}
where \(\odot\) denotes elementwise multiplication and all spin and time
indices are suppressed.  Since \(\omega\) is symmetric, a trace contraction
may equivalently attach this factor to the two-index current vertex, as done
in the production code.  This algebraic transfer does not justify multiplying
three separately reconstructed propagator estimators made from the same noise
set: a baryon contraction can require higher joint-noise weights or
independent noise sets.  The factorized contraction preserves the original
index structure and is therefore the production definition.

\subsection{Flowed all-time basis}

In the older equations below, \(\tau\) labels an original lattice source
time, not flow time. In Section~\ref{sec:blending-oct15} this label is
renamed \(\sigma\), reserving \(\tau=\tf/a^2\) for flow time. The cache
directory \path{basis/tauNNN/} still uses the older source-time spelling.

Embed a time-slice eigenvector in four dimensions,
\begin{equation}
  \eta_{\tau n}(y)
  =
  \delta_{y_4,\tau}\,V_n(\tau,\bm y),
\end{equation}
and evolve it with the four-dimensional kernel:
\begin{equation}
  \Phi_{\tau n}^{(\tf)}(x)
  =
  \sum_y K(\tf;x,y)\eta_{\tau n}(y).
  \label{eq:flowed-basis}
\end{equation}
For \(\tf>0\), \(\Phi_{\tau n}^{(\tf)}(\bm x,t_J)\) has support at
\(\tau\ne t_J\); three-dimensional smearing has no such mixing.

For a displacement \(z\) in direction \(i\), define the spin-independent
generalized distillation vertex
\begin{align}
  J_{mn}^{(\tf)}
  (t_J;\tau_a,\tau_b;z,i)
  &=
  \sum_{\bm x}
  \left[
    \Phi_{\tau_a m}^{(\tf)}(\bm x,t_J)
  \right]^\dagger
  W_{B(\tf),i}(\bm x,\bm x+z\hat{\bm i};t_J)
  \Phi_{\tau_b n}^{(\tf)}
  (\bm x+z\hat{\bm i},t_J).
  \label{eq:generalized-vertex}
\end{align}
The vertex is spin independent.  The subsequent C3 files retain both
insertion-spin indices, so \(\Gamma\) can be applied later without repeating
the flowed vertex calculation.

At zero flow time,
\begin{equation}
  \Phi_{\tau n}^{(0)}(\bm x,t_J)
  =
  \delta_{\tau,t_J}V_n(t_J,\bm x),
\end{equation}
and therefore
\begin{equation}
  J^{(0)}(t_J;\tau_a,\tau_b;z,i)
  =
  \delta_{\tau_a,t_J}\delta_{\tau_b,t_J}
  V^\dagger(t_J)W_i(z;t_J)V(t_J).
  \label{eq:zero-flow-identity}
\end{equation}
This identity provides the exact zero-flow regression against the original
\(V^\dagger D(z)V\) vertex.

\subsection{From the equations to the production code}
\label{sec:basis-vertex-code}

This subsection fixes the implementation meaning of ``basis'' and
``vertex''.  They are not two independent physical smearings.  The flowed
basis is the set of right-hand sides obtained by applying the four-dimensional
fermion-flow kernel to the time-slice LapH vectors.  The vertex is the
matrix-valued representation of the flowed bilocal between two such flowed
bases.  The nucleon perambulators are not flowed again at this stage.

For the L32x96 and L32x64 production branches, the shared configuration object
stores the flow time, step size, temporal window, number of modes, separation
range, input gauge scheme, and all input/output roots.  Its cache tag records
the flow time, step size, temporal window, mode count and separation count.
For example,
\begin{center}
\path{thin_tf0p50000000_eps0p01000000_tw4_Nv400_dz21}
\end{center}
identifies \(\tf/a^2=0.5\), \(\epsilon=\Delta(\tf/a^2)=0.01\), \(w=4\), 400 current
modes and 21 stored separation lengths.
The ``thin'' prefix is a scheme identifier: the LIME gauge field is read
without HYP pre-smearing.  A cache made from HYP links must not be silently
reused for the thin-link matching scheme.

\paragraph{Basis stage.}
For every source time \(\tau\), the code reads the stored eigenvectors
\(V_n(\tau,\bm x,c)\), embeds them as
\(\eta_{\tau n}(t,\bm x,c)=\delta_{t\tau}V_n(\tau,\bm x,c)\), and packs
four modes into the four spin lanes of a QUDA multi-right-hand-side field.
The spin lanes are storage lanes: the covariant Laplacian is spin diagonal and
does not mix the physical spin components.  A batch of 32 modes therefore
contains eight Wilson spinors, while the final batch contains four.  With
\(N_v=400\), one source time requires
\begin{equation}
  \left\lceil\frac{400}{32}\right\rceil=13
\end{equation}
QUDA calls.  Each call applies the coupled gauge--fermion Runge--Kutta flow
to the supplied right-hand sides; it does not invert a Dirac matrix and does
not construct the dense kernel \(K\).

The output is unpacked from QUDA's even--odd layout, and the retained current
time slices are written as
\begin{equation}
  \Phi_{\tau n}^{(\tf)}(t_J,\bm x,c),
  \qquad
  \text{shape }(N_v,L_z,L_y,L_x,N_c)=(400,32,32,32,3).
  \label{eq:code-flowed-basis-shape}
\end{equation}
For a finite temporal window \(w\), the saved slices satisfy
\( |t_J-\tau|_{\rm periodic}\leq w\).  The basis metadata records the flow
schema, \(\tau\), all stored \(t_J\), \(N_v\), data type, discarded norm,
and gauge provenance.  On restart, a ``done.json'' marker is accepted only
when all corresponding \texttt{.npy} slices exist and the metadata passes
validation.

\paragraph{Vertex stage.}
For each current time \(t_J\), each ordered source-time pair
\((\tau_a,\tau_b)\) in the retained window, and each spatial direction
\(i\in\{x,y,z\}\), the code loads the two flowed basis slices and the flowed
gauge links.  It builds positive and negative straight Wilson products
recursively.  In matrix notation the two stored signs are
\begin{align}
  J_{-,mn}^{(\tf)}(t_J;\tau_a,\tau_b;d,i)
  &=\sum_{\bm x}\Phi_{\tau_a m}^{(\tf)\dagger}(t_J,\bm x)
    W_i(\bm x,\bm x-d\hat i)\Phi_{\tau_b n}^{(\tf)}(t_J,\bm x-d\hat i),\\
  J_{+,mn}^{(\tf)}(t_J;\tau_a,\tau_b;d,i)
  &=\sum_{\bm x}\Phi_{\tau_a m}^{(\tf)\dagger}(t_J,\bm x)
    W_i(\bm x,\bm x+d\hat i)\Phi_{\tau_b n}^{(\tf)}(t_J,\bm x+d\hat i).
  \label{eq:code-vertex-signs}
\end{align}
The implementation stores both signs and \(d=0,\ldots,N_z-1\) in an array
with axes
\begin{equation}
  (\texttt{vdv\_sign},\texttt{delta\_z},
   \texttt{eigen\_left},\texttt{eigen\_right})
  =(2,N_z,N_v,N_v).
  \label{eq:code-vertex-shape}
\end{equation}
In historical scheme A, the stochastic two-index factor \(\omega_{mn}\) from
Eq.~\eqref{eq:stochastic-coordinate-propagator} is multiplied into this
vertex before it is saved. Scheme B instead uses the source-time-dependent
table in Eq.~\eqref{eq:scheme-B-oct15}, once at the same location.
The spin matrix \(\Gamma\) is deliberately not
contracted here.  Open current-spin indices appear in the C3 product after
combining this spin-independent vertex with the spinful perambulators; the
same vertex then feeds unpolarized, helicity, and transversity projections.

With \texttt{FLOW\_WORK\_ROOT} as the cache root, the durable files follow
\begin{center}
\begin{tabular}{>{\raggedright\arraybackslash}p{0.20\textwidth}
                >{\raggedright\arraybackslash}p{0.69\textwidth}}
\toprule
Object & Path below \texttt{FLOW\_WORK\_ROOT/<conf>/<tag>/} \\
\midrule
Basis slice & \path{basis/tauNNN/tJMMM.npy}; one
\path{basis/tauNNN/done.json} per kernel source time \(\tau\). \\
Flowed gauge & \path{flowed_gauge/index.json} plus time-sharded
\texttt{.npy} links. \\
Vertex block & \path{vertex/tJMMM/taNNN_tbKKK_z.npy} and a same-name
\texttt{.json} manifest; \(z\) is the production direction. \\
\bottomrule
\end{tabular}
\end{center}

The corresponding code path is
\begin{center}
\begin{tabular}{>{\raggedright\arraybackslash}p{0.32\textwidth}
                >{\raggedright\arraybackslash}p{0.57\textwidth}}
\toprule
Physical object & Implementation and contract \\
\midrule
Flow parameters and cache paths
 & \texttt{flow\_common\_l32x64.py}, \texttt{FlowConfig};
 \texttt{n\_steps=round(FLOW\_TIME/FLOW\_EPSILON)}, thin-link guard,
 \texttt{basis/}, \texttt{vertex/}, and \texttt{flowed\_gauge/} roots. \\
Flowed basis
 & \texttt{generate\_flowed\_vertex\_l32x64.py},
 \texttt{generate\_flowed\_basis}; mode packing,
 \texttt{performGFlowQuda}, unpacking, norm and metadata. \\
Flowed gauge
 & \texttt{generate\_flowed\_vertex\_l32x64.py},
 \texttt{save\_flowed\_gauge\_once};
 four-dimensional links and an \texttt{index.json} shard manifest. \\
Generalized vertex
 & \texttt{flow\_common\_l32x64.py}, \texttt{generalized\_vertex};
 recursive \(\pm d\) Wilson products and the \((2,N_z,N_v,N_v)\) result. \\
Vertex production
 & \texttt{generate\_flowed\_vertex\_l32x64.py},
 \texttt{generate\_vertices};
 ordered \((t_J,\tau_a,\tau_b,i)\) blocks, metadata validation, and atomic
 \texttt{.npy}/\texttt{.json} publication. \\
End-to-end orchestration
 & \texttt{full\_chain.py}; basis, vertex,
 input validation, C2/C3 contraction, output validation, and vertex-only
 cache cleanup. \\
\bottomrule
\end{tabular}
\end{center}

All filenames in this table are relative to
\path{/public/group/lqcd/donghx/Gradient_Flow_GluonPDF/quark_l32x64_flow_20261007/};
the basis and vertex implementation files are in its \texttt{reference/}
subdirectory.  The L32x64 orchestrator runs four single-GPU basis workers on
disjoint source times, then four MPI ranks for vertex generation and C2/C3
contraction.  This scheduling is distinct from the four-rank domain-decomposed
flow description used for the L32x96 branch.

The L32x96 branch has the corresponding files
\nolinkurl{/public/group/lqcd/donghx/Gradient_Flow/Blending_L32x96/flow_common_l32x96.py},
\nolinkurl{generate_flowed_vertex_l32x96.py}, and the factorized C2/C3 driver.
The tensor contract is the same, while \(N_t\), flow-time list, temporal
windows, and ensemble paths differ.  Since the vertex is reconstructible from
the retained basis and flowed gauge, production may delete only
\texttt{vertex/} after a successful output guard; deleting \texttt{basis/} or
\texttt{flowed\_gauge/} removes the inexpensive restart path and requires
repeating the flow stage.

\paragraph{Contraction boundary.}
The subsequent C2/C3 driver reads the vertex blocks, contracts them with the
ordinary all-time perambulators, applies the open-spin matrices and proton
projectors, and writes the full time-axis correlators.  Thus the logical chain
is
\begin{equation}
  V,\;D_f^{-1}
  \;\longrightarrow\;
  \Phi=KV,\;B(\tf)
  \;\longrightarrow\;
  J=\Phi^\dagger W_{B(\tf)}\Phi
  \;\longrightarrow\;
  C_2,C_3.
  \label{eq:basis-vertex-c3-chain}
\end{equation}
The gluon OPE branch does not use this quark distillation basis or generalized
vertex: it flows the gauge field, constructs Clover field strengths and the
gluon Wilson line, and evaluates a vacuum-subtracted disconnected nucleon
correlator.  Reusing the words ``basis'' or ``vertex'' for the gluon OPE
would therefore mix two different computational layers.

\subsection{Why all-time perambulators avoid new Dirac inversions}

In a Wick contraction, the active lines contain
\begin{equation}
  D_f^{-1}K^\dagger
  \quad\text{and}\quad
  KD_f^{-1},
\end{equation}
After inserting the all-time distillation basis, both kernels enter
Eq.~\eqref{eq:generalized-vertex}; propagation remains in ordinary
perambulators:
\begin{equation}
  C_{3,f}^{(\tf)}
  \sim
  \sum_{\tau_a,\tau_b}
  \mathcal W_f
  \left[
    \tau_f(t_{\rm sink},\tau_a)\,
    J^{(\tf)}(t_J;\tau_a,\tau_b)\,
    \tau_f(\tau_b,t_{\rm src})
  \right],
  \label{eq:c3-perambulator-schematic}
\end{equation}
where \(\mathcal W_f\) contains spectator lines, proton tensors, spin matrices,
and Wick signs.

No new system involving \(D_f\) is needed if all blocks in the double sum are
stored. The present light data meet this condition. New inversions are needed
if the required times or distillation subspace are absent, as for the
time-diagonal strange data in Section~\ref{sec:3d-vs-4d}, or if flow changes
the fermion action rather than only the operator.

For completeness, a propagator with both endpoints flowed can be written
\begin{equation}
  S_f^{(\tf,\tf')}(x,y)
  =
  \sum_{u,v}
  K(\tf;x,u)\,
  D_f^{-1}(u,v)\,
  K^\dagger(\tf';v,y).
  \label{eq:flowed-propagator}
\end{equation}
This is a composition of \(D_f^{-1}\) and flow kernels, not a new Dirac
inverse.

In the present stochastic representation, the corresponding coordinate-space
estimator is
\begin{equation}
  \widehat S_{f,\omega}^{(\tf,\tf')}
  =K(\tf)V\bigl(\omega\odot\tau_f\bigr)V^\dagger K^\dagger(\tf').
  \label{eq:flowed-stochastic-propagator}
\end{equation}
Linearity of \(K\) preserves the intended stochastic completion: the flow
changes the sampled vectors but not their source-sector probabilities.  Since
the heat kernel is not unitary, \((KV)^\dagger(KV)\neq\mathbf 1\) is expected.
Reorthogonalizing \(KV\) and discarding the associated transformation would
change the flowed operator and invalidate the original weights.

\section{Three-dimensional and four-dimensional flow with the available
perambulators}
\label{sec:3d-vs-4d}

\subsection{Precise definitions}

Here, three-dimensional flow evolves each Euclidean time slice independently:
only spatial links enter the gauge flow, and fermions use the spatial
covariant Laplacian:
\begin{align}
  \partial_{\tf}B_i^{\rm 3D}
  &=
  D_jG_{ji}^{\rm 3D},
  &
  B_i^{\rm 3D}(0,\bm x,x_4)
  &=
  A_i(\bm x,x_4),
  \qquad i,j=1,2,3,
  \label{eq:3d-gauge-flow}
  \\
  \partial_{\tf}\chi_{\rm 3D}
  &=
  \Delta_{\rm 3D}\chi_{\rm 3D},
  &
  \Delta_{\rm 3D}
  &=
  \sum_{i=1}^{3}D_iD_i,
  \qquad
  \chi_{\rm 3D}(0)=\psi .
  \label{eq:3d-fermion-flow}
\end{align}
Whether temporal links are fixed or omitted must be stated. Standard
four-dimensional flow evolves all directions and uses
\begin{equation}
  \Delta_{\rm 4D}
  =
  \sum_{\mu=1}^{4}D_\mu D_\mu
  =
  \Delta_{\rm 3D}+D_4D_4.
  \label{eq:3d-4d-laplacians}
\end{equation}

The kernels obey
\begin{align}
  K_{\rm 3D}
  (\tf;\bm x,x_4;\bm y,y_4)
  &=
  \delta_{x_4y_4}\,
  K_{\rm 3D}(\tf;\bm x,\bm y;x_4),
  \label{eq:3d-time-diagonal-kernel}
  \\
  K_{\rm 4D}
  (\tf;\bm x,x_4;\bm y,y_4)
  &\neq 0
  \quad\text{in general for }x_4\neq y_4 .
  \label{eq:4d-time-mixing-kernel}
\end{align}
For a free lattice field,
\begin{align}
  K_{\rm 3D}(\tf;\bm p)
  &=
  \exp\!\left[-\tf\widehat{\bm p}^{\,2}\right],
  &
  \widehat{\bm p}^{\,2}
  &=
  \frac{4}{a^2}
  \sum_{i=1}^{3}\sin^2\!\left(\frac{ap_i}{2}\right),
  \label{eq:free-3d-kernel}
  \\
  K_{\rm 4D}(\tf;p)
  &=
  \exp\!\left[
    -\tf\left(\widehat{\bm p}^{\,2}+\widehat p_4^{\,2}\right)
  \right].
  \label{eq:free-4d-kernel}
\end{align}
The root-mean-square diffusion distances are
\begin{equation}
  r_{\rm rms}^{\rm 3D}=\sqrt{6\tf},
  \qquad
  r_{\rm rms}^{\rm 4D}=\sqrt{8\tf}.
  \label{eq:3d-4d-radii}
\end{equation}
Thus \(\sqrt{8\tf}\) is the four-dimensional smoothing radius.

Flowing only links or only quarks defines other schemes. A calculation must
state which fields are flowed; ``3D'' or ``4D'' alone is insufficient.

\subsection{Ultraviolet damping: what can and cannot be inferred}

The difference between Eqs.~\eqref{eq:free-3d-kernel} and
\eqref{eq:free-4d-kernel} is decisive in ultraviolet power counting.  A
four-dimensional flowed leg supplies
\begin{equation}
  e^{-\tf p^2}
  =e^{-\tf(\bm p^2+p_4^2)},
\end{equation}
so every component of a large Euclidean loop momentum is exponentially
damped.  This property is used in the renormalization proof for standard
gradient-flow observables: after the parameters and boundary fields of QCD
are renormalized, positive-flow-time gauge observables are finite, while
flowed quark fields require only a multiplicative normalization
\cite{LuscherWeisz2011,Luscher2013}.

Three-dimensional flow instead supplies only
\begin{equation}
  e^{-\tf\bm p^2}.
\end{equation}
It suppresses the corner \(|\bm p|\to\infty\), but it does not suppress the
independent limit \(|p_4|\to\infty\) at fixed \(\bm p\).  This can be made
explicit with the generic integral
\begin{equation}
  I_{r,n}^{\rm 3D}(\tf)
  =
  \int\!\frac{\dd^3\bm p\,\dd p_4}{(2\pi)^4}
  \frac{(p_4^2)^r e^{-2\tf\bm p^2}}
       {(p_4^2+\bm p^2)^n}.
  \label{eq:3d-generic-uv-integral}
\end{equation}
For fixed \(\bm p\), its large-energy behavior is
\begin{equation}
  I_{r,n}^{\rm 3D}
  \underset{|p_4|\to\infty}{\sim}
  \int^{\Lambda_4}\!\dd p_4\,
  |p_4|^{2r-2n}.
  \label{eq:3d-energy-power-counting}
\end{equation}
The energy integral converges only if \(n-r>1/2\).  There is no universal
exponential bound.  By contrast, standard 4D flow multiplies the integrand by
\(e^{-2\tf p_4^2}\), making every polynomial numerator integrable.

A concrete counterexample is a composite operator containing temporal
derivatives of a spatially flowed gauge field.  At tree level it contains
\begin{equation}
  \left\langle
    \Tr\bigl[\partial_4B_i^{\rm 3D}(\tf,x)
             \partial_4B_i^{\rm 3D}(\tf,x)\bigr]
  \right\rangle
  \propto
  \int\!\dd^3\bm p\,\dd p_4\,
  e^{-2\tf\bm p^2}
  \frac{p_4^2}{p_4^2+\bm p^2}.
  \label{eq:3d-temporal-derivative-counterexample}
\end{equation}
The integrand approaches \(e^{-2\tf\bm p^2}\) as
\(|p_4|\to\infty\), so the remaining \(p_4\) integral is linearly divergent.
Replacing 3D by 4D flow adds \(e^{-2\tf p_4^2}\) and removes this divergence.
Equation~\eqref{eq:3d-temporal-derivative-counterexample} does not say that
every equal-time 3D-flowed bilocal diverges.  It proves the narrower and
important statement that spatial flow alone does not imply the general
positive-flow-time finiteness theorem.

Consequently, three separate claims must not be conflated:
\begin{enumerate}
  \item 3D flow suppresses large \emph{spatial} momentum;
  \item it may render a particular equal-time diagram finite;
  \item it renders the complete composite operator finite to all orders with
        only the standard QCD and ringed-field counterterms.
\end{enumerate}
The first statement is immediate, the second must be checked diagram by
diagram, and the third requires a new anisotropic renormalization proof.

\subsection{One-loop Wilson-line self-energy and the linear divergence}

The straight spatial Wilson line has a power-divergent self-energy in a cutoff
scheme.  Its leading one-loop structure, suppressing gauge-parameter and
inessential numerical factors, is
\begin{equation}
  \mathcal I(\tf,z)
  =g^2C_F
  \int\!\frac{\dd^4k}{(2\pi)^4}
  \frac{4\sin^2(k_z z/2)}{k_z^2}
  \frac{F_{\tf}(k)}{k^2},
  \label{eq:wilson-line-self-energy}
\end{equation}
where the two gauge fields attached to the flowed line give
\begin{align}
  F_{\tf}^{\rm 4D}(k)
  &=e^{-2\tf(k_4^2+\bm k^2)},
  &
  F_{\tf}^{\rm 3D}(k)
  &=e^{-2\tf\bm k^2}.
  \label{eq:3d-4d-line-flow-factors}
\end{align}
For a line long compared with the regulator radius,
\begin{equation}
  \frac{4\sin^2(k_z z/2)}{k_z^2}
  \longrightarrow
  2\pi|z|\,\delta(k_z).
  \label{eq:long-line-delta}
\end{equation}
Standard 4D flow therefore gives
\begin{equation}
  \mathcal I_{\rm lin}^{\rm 4D}
  \propto
  g^2C_F|z|
  \int\!\frac{\dd^3k_\perp}{(2\pi)^3}
  \frac{e^{-2\tf k_\perp^2}}{k_\perp^2}
  \propto
  g^2C_F\frac{|z|}{\sqrt{\tf}}.
  \label{eq:4d-flow-linear-term}
\end{equation}
The line is finite at every fixed \(\tf>0\); the original linear divergence
reappears as the calculable short-flow-time behavior \(|z|/\sqrt{\tf}\).

For 3D flow, write \(\bm q=(k_x,k_y)\) after setting \(k_z=0\).  The undamped
energy integral in this particular graph is nevertheless convergent:
\begin{equation}
  \int\!\frac{\dd k_4}{2\pi}
  \frac{1}{k_4^2+\bm q^2}
  =\frac{1}{2|\bm q|}.
  \label{eq:3d-line-energy-integral}
\end{equation}
Hence
\begin{equation}
  \mathcal I_{\rm lin}^{\rm 3D}
  \propto
  g^2C_F|z|
  \int\!\frac{\dd^2\bm q}{(2\pi)^2}
  \frac{e^{-2\tf\bm q^2}}{2|\bm q|}
  \propto
  g^2C_F\frac{|z|}{\sqrt{\tf}}.
  \label{eq:3d-flow-linear-term}
\end{equation}
This establishes a useful but limited result: spatial flow can regulate the
one-loop straight-line self-energy and therefore can suppress the lattice
linear divergence associated with that subgraph.  The coefficient in
Eq.~\eqref{eq:3d-flow-linear-term} is not the 4D coefficient, because the
regulator functions in Eq.~\eqref{eq:3d-4d-line-flow-factors} differ.
Endpoint, vertex, quark-normalization, and higher-loop subgraphs still require
a dedicated 3D analysis.  The multiplicative renormalization structure of the
unflowed straight operator is discussed, for example, in
Ref.~\cite{Ishikawa2017}; replacing its cutoff by a flow radius is a scheme
choice, not the removal of renormalization physics.

The order of limits is also essential.  At fixed physical \(\tf>0\), the
quantity \(|z|/\sqrt{\tf}\) is finite as \(a\to0\).  Holding instead
\(\tf/a^2\) fixed makes \(\sqrt{\tf}\propto a\), so
Eq.~\eqref{eq:3d-flow-linear-term} again scales as \(|z|/a\).  A fixed number
of lattice flow steps is therefore not by itself a continuum
renormalization prescription.

\subsection{Why the established matching calculation uses 4D flow}

The equal-time kinematics of a quasi-PDF do not mathematically forbid a 3D
smearing scheme.  Four-dimensional flow is required here for a more specific
reason: the available flowed quasi-PDF finiteness arguments and perturbative
matching coefficients define the operator with standard 4D gauge and fermion
flow \cite{MonahanOrginos2017,Monahan2018}.  They rely on
\begin{equation}
  F_{\tf}(k)=e^{-2\tf k^2},
  \qquad
  \Delta=D_\mu D_\mu,
  \qquad
  \rflow=\sqrt{8\tf}.
  \label{eq:standard-flow-matching-contract}
\end{equation}
Changing any of these ingredients changes the regulator and the finite part
of the matching kernel.

Moreover, 3D flow introduces a preferred unit vector
\(n_\mu=(0,0,0,1)\) before the external hadron momentum and Wilson-line
direction are specified.  A 3D matching coefficient may depend separately on
\begin{equation}
  \tf\bm p^2,
  \qquad p_4,
  \qquad n\!\cdot p,
  \qquad n\!\cdot z,
\end{equation}
and the reduced symmetry permits counterterms and operator mixing that are
not organized by 4D Euclidean representations.  A consistent 3D program is
possible in principle, but it must
\begin{enumerate}
  \item prove renormalizability of the anisotropically flowed bilocal;
  \item define a 3D analogue of ringed-fermion normalization;
  \item calculate its Wilson-line mass term, endpoint terms, and mixing;
  \item derive a new 3D-flow-to-\(\MSbar\) or direct light-cone matching
        coefficient;
  \item validate the result nonperturbatively and take the continuum limit at
        fixed physical spatial flow radius.
\end{enumerate}
Until these steps are completed, a 3D-flowed matrix element is a legitimate
smeared observable but not an input to the existing 4D-flow matching formula.

\subsection{Physical and renormalization differences}

An equal-time quasi-PDF operator does not make 3D and 4D regulators
equivalent. Four-dimensional flow preserves continuum Euclidean covariance
and underlies the established finiteness and matching framework
\cite{Luscher2013,MonahanOrginos2017,Monahan2018}. Timeslice flow selects a
preferred direction and requires separate matching.

\begin{longtable}{
  >{\raggedright\arraybackslash}p{0.22\textwidth}
  >{\raggedright\arraybackslash}p{0.35\textwidth}
  >{\raggedright\arraybackslash}p{0.35\textwidth}}
  \caption{Comparison of the three- and four-dimensional flowed operator
  definitions.}\label{tab:3d-4d-comparison}\\
  \toprule
  Property & Three-dimensional flow & Four-dimensional flow \\
  \midrule
  \endfirsthead
  \toprule
  Property & Three-dimensional flow & Four-dimensional flow \\
  \midrule
  \endhead
  Covariant Laplacian
  & \(\sum_{i=1}^{3}D_iD_i\)
  & \(\sum_{\mu=1}^{4}D_\mu D_\mu\) \\
  Temporal support
  & Exactly time diagonal:
    \(K_{\rm 3D}\propto\delta_{x_4y_4}\)
  & Nonzero temporal tails for every \(\tf>0\) \\
  Free damping
  & \(e^{-\tf\widehat{\bm p}^{\,2}}\)
  & \(e^{-\tf(\widehat{\bm p}^{\,2}+\widehat p_4^{\,2})}\) \\
  Large-\(p_4\) behavior
  & Undamped; convergence depends on the numerator and propagator powers of
    each subgraph
  & Exponentially damped by \(e^{-\tf p_4^2}\) \\
  Straight-line self-energy
  & Finite at one loop for fixed \(\tf>0\), with a scheme-specific
    \(|z|/\sqrt{\tf}\) term
  & Finite at fixed \(\tf>0\), with the standard flowed
    \(|z|/\sqrt{\tf}\) term \\
  General finiteness theorem
  & Not implied by spatial damping; requires a new anisotropic proof and
    counterterm analysis
  & Established after ordinary QCD renormalization and flowed-quark
    normalization \\
  Fermion temporal boundary
  & Does not enter the kernel evolution
  & Must be implemented; the present light calculation uses anti-periodic
    temporal boundary conditions \\
  Continuum symmetry
  & Spatial rotations with a preferred time direction
  & Four-dimensional Euclidean symmetry before the external quasi-PDF
    kinematics are imposed \\
  Standard flowed matching
  & Requires a separate 3D derivation or nonperturbative scheme conversion
  & Existing gradient-flow quasi-PDF framework applies when all other scheme
    details agree \\
  Perambulator time blocks
  & Diagonal block \(\tau_f(t_J,t_J)\) is sufficient
  & Off-diagonal blocks \(\tau_f(\tau_b,\tau_a)\) are required \\
  Main numerical advantage
  & Independent time slices and substantially smaller storage
  & Standard flowed-field scheme, at the cost of temporal mixing \\
  \bottomrule
\end{longtable}

A 3D-flowed strange observable is well defined at finite lattice spacing as a
distinct spatially smeared scheme.  Equations~\eqref{eq:3d-line-energy-integral}
and \eqref{eq:3d-flow-linear-term} show that it can suppress the straight
Wilson-line linear divergence, but
Eq.~\eqref{eq:3d-temporal-derivative-counterexample} shows why this is not a
general UV-finiteness proof.  It cannot use 4D matching solely because
\(\tf/a^2\) is the same.

\subsection{Time-block content of the light and strange data}

The time labels of a projected perambulator are
\begin{equation}
  \tau_f(t_2,t_1)
  =
  V^\dagger(t_2)\,
  D_f^{-1}(t_2,t_1)\,
  V(t_1),
  \label{eq:perambulator-time-block}
\end{equation}
Here \(t_1\) and \(t_2\) are source and sink times. Availability of arbitrary
\(\tau_f(t_2,t_1)\), rather than only \(\tau_f(t_1,t_1)\), is decisive.

Each light-flavor file contains all 96 sink times for one source time, so the
contraction can read
\begin{equation}
  \tau_\ell(t_{\rm sink},\tau_a),
  \qquad
  \tau_\ell(\tau_b,t_{\rm src}),
  \label{eq:light-required-blocks}
\end{equation}
for every retained \(\tau_a,\tau_b\). Hence the 4D light calculation needs no
new inversion.

The strange reader instead uses files of the form
\begin{equation}
  \texttt{tXXX\_tsinkXXX\_gen\ldots},
\end{equation}
with equal source and sink labels. Its generator saves only the source-time
slice. Thus the available data are
\begin{equation}
  \bigl\{\tau_s(t,t)\,;\ t=0,\ldots,N_t-1\bigr\},
  \label{eq:strange-diagonal-data}
\end{equation}
rather than the full \(\tau_s(t_2,t_1)\).

\subsection{Why the diagonal strange data are sufficient for 3D flow}

Equation~\eqref{eq:3d-time-diagonal-kernel} keeps both quark fields at
\(t_J\). The flowed vertex is
\begin{equation}
  J_{\rm 3D}^{(\tf)}(t_J;z)
  =
  \bigl[K_{\rm 3D}(\tf)V(t_J)\bigr]^\dagger
  W_{B^{\rm 3D}(\tf)}(z;t_J)
  \bigl[K_{\rm 3D}(\tf)V(t_J)\bigr].
  \label{eq:3d-vertex}
\end{equation}
The disconnected strange loop has the schematic form
\begin{equation}
  L_{s,\Gamma}^{\rm 3D}(t_J,z)
  =
  -\Tr_{\rm spin,V}
  \left[
    \Gamma\,
    J_{\rm 3D}^{(\tf)}(t_J;z)\,
    \tau_s(t_J,t_J)
  \right].
  \label{eq:3d-strange-loop}
\end{equation}
No off-diagonal block appears, so 3D flow can reuse the strange data without
new inversions.

The eigenvector, spin, color, stochastic-weight, smearing, and boundary
conventions must still agree; orthonormality alone is insufficient.

\subsection{Why complete 4D flow needs off-diagonal strange blocks}

Four-dimensional flow has two independent kernel source times:
\begin{equation}
  L_{s,\Gamma}^{\rm 4D}(t_J,z)
  =
  -\sum_{\tau_a,\tau_b}
  \Tr_{\rm spin,V}
  \left[
    \Gamma\,
    J_{\rm 4D}^{(\tf)}
      (t_J;\tau_a,\tau_b;z)\,
    \tau_s(\tau_b,\tau_a)
  \right].
  \label{eq:4d-strange-loop}
\end{equation}
Stored axes may transpose \(\tau_s(\tau_b,\tau_a)\), but both time labels must
be available.

It is useful to separate Eq.~\eqref{eq:4d-strange-loop} into
\begin{align}
  L_{s,\Gamma}^{\rm 4D}
  &=
  L_{s,\Gamma}^{\rm diag}
  +
  L_{s,\Gamma}^{\rm off},
  \\
  L_{s,\Gamma}^{\rm diag}
  &=
  -\sum_{\tau}
  \Tr_{\rm spin,V}
  \left[
    \Gamma J_{\rm 4D}^{(\tf)}(t_J;\tau,\tau;z)
    \tau_s(\tau,\tau)
  \right],
  \\
  L_{s,\Gamma}^{\rm off}
  &=
  -\sum_{\tau_a\neq\tau_b}
  \Tr_{\rm spin,V}
  \left[
    \Gamma J_{\rm 4D}^{(\tf)}(t_J;\tau_a,\tau_b;z)
    \tau_s(\tau_b,\tau_a)
  \right].
  \label{eq:4d-strange-diag-offdiag}
\end{align}
The existing data determine only the diagonal term. At \(\tf>0\), the
off-diagonal term cannot be reconstructed from \(\tau_s(t,t)\),
orthonormality, \(\gamma_5\)-hermiticity, or \(K\). Restricting to equal times
is an approximation, not complete 4D flow.

\subsection{Required recomputation and economical storage}

Without saved full solutions, exact 4D strange flow requires new inversions.
For each source time, spin, and distillation vector, solve
\begin{equation}
  D_s X_{\tau_a n\beta}
  =
  \delta_{t,\tau_a}\,V_n(\tau_a)\,e_\beta ,
  \label{eq:strange-required-inversion}
\end{equation}
and projects the solution at the required sink times:
\begin{equation}
  \tau_s(\tau_b,\tau_a)_{\alpha\beta;mn}
  =
  V_m^\dagger(\tau_b)\,
  X_{\tau_a n\beta}(\tau_b)_\alpha .
  \label{eq:strange-required-projection}
\end{equation}
A solve contains all sink times before output truncation, so the generator
should project and save each required \(\tau_b\) block.

With a temporal kernel window \(w\), both \(\tau_a\) and \(\tau_b\) obey
\begin{equation}
  |\tau_a-t_J|_{\rm periodic}\leq w,
  \qquad
  |\tau_b-t_J|_{\rm periodic}\leq w.
  \label{eq:window-two-times}
\end{equation}
Only the band
\begin{equation}
  |\tau_b-\tau_a|_{\rm periodic}\leq2w
  \label{eq:banded-strange-storage}
\end{equation}
is needed. For \(w=4\), storing temporal distances through eight reduces
storage. All 96 current times still require all source times.

Before new inversions, audit any legacy full-time strange files. Mass, actions,
smearing, boundary conditions, eigenvectors, weights, axis order, and
configuration coverage must match. Compatible files can be reused.

\paragraph{Practical conclusion.}
The light data support complete 4D connected flow without new inversions. The
strange data support 3D flow only. Complete 4D strange flow needs validated
off-diagonal blocks or new inversions; any diagonal approximation must be
labelled and tested separately.

\section{Gradient-flowed gluon quasi-PDF observables}
\label{sec:flowed-gluon}

This section specifies the theory conventions and the actual estimator used
by the present gluon calculation.  The unpolarized operator follows the
structure used in Ref.~\cite{ChenEtAl2025GluonPDF}; the helicity tensor
combination and its odd-separation projection follow
Ref.~\cite{EgererEtAl2022GluonHelicity}.  Those references do not, by
themselves, define the renormalization of the present thin-link flowed
operator.  In particular, the HYP-smearing and hybrid-renormalization factors
of Ref.~\cite{ChenEtAl2025GluonPDF} cannot be copied into this flow scheme.

\subsection{From light-cone gluon PDFs to a Euclidean bilocal}

The light-cone unpolarized and helicity gluon distributions are encoded in
forward nucleon matrix elements of two field strengths separated along a
lightlike direction.  Such real-time separations cannot be implemented in a
Euclidean Monte Carlo calculation.  As in LaMET or the pseudo-PDF approach,
one instead measures an equal-time bilocal with a straight spatial separation

\begin{equation}
  z_\mu=(0,0,z,0)_{\mathrm E},
  \qquad
  \nu=P_z z,
  \label{eq:gluon-ioffe-time}
\end{equation}
where the implementation uses spatial direction \(z\equiv2\) and Euclidean
time direction \(t\equiv3\).  The Ioffe time \(\nu\) is dimensionless.
Large \(P_z\) or short spacelike \(z^2\), together with a perturbative
factorization formula, is needed to connect the Euclidean matrix element to a
light-cone PDF.  The bare lattice correlator is only the first layer of this
construction.

\subsection{Thin-link four-dimensional Wilson flow}

Let \(B_\mu(\tau,x)\) be the flowed gauge field and
\(G_{\mu\nu}(\tau,x)\) its field strength.  In the continuum, Wilson flow is
defined by

\begin{align}
  \partial_\tau B_\mu(\tau,x)
  &=D_\nu G_{\nu\mu}(\tau,x),
  &B_\mu(0,x)&=A_\mu(x),
  \label{eq:gluon-wilson-flow-continuum}\\
  G_{\mu\nu}
  &=\partial_\mu B_\nu-\partial_\nu B_\mu+[B_\mu,B_\nu].
\end{align}
The flow time has dimensions of length squared and introduces the smoothing
radius
\begin{equation}
  r_{\rm flow}=\sqrt{8\tau}.
  \label{eq:gluon-flow-radius}
\end{equation}
The lattice code evolves the original, unsmeared links with a third-order
Runge--Kutta Wilson-flow integrator using
\(\epsilon=0.01\).  Thus \(\tau/a^2=N_{\rm step}\epsilon\).  No HYP-smearing
is applied before the flow.  Positive flow time suppresses ultraviolet gauge
modes, but it also changes the operator: varying \(\tau\) is a scheme and
resolution change, not merely a noise-tuning choice.

\subsection{Clover field strength and dual tensor}

At each flowed link field the code constructs the oriented four-leaf Clover
sum \(Q_{\mu\nu}(\tau,x)\) and its anti-Hermitian part,
\begin{equation}
  F^{\rm clov}_{\mu\nu}(\tau,x)
  =-\frac{\ii}{8}\left[Q_{\mu\nu}(\tau,x)
       -Q_{\mu\nu}^{\dagger}(\tau,x)\right],
  \label{eq:flowed-clover-strength}
\end{equation}
in the normalization of the production code.  The physical production field
is projected onto the Lie algebra,
\begin{equation}
  F_{\mu\nu}\longrightarrow
  F_{\mu\nu}-\frac{1}{3}\Tr(F_{\mu\nu})\mathbf 1.
  \label{eq:color-traceless}
\end{equation}
The legacy untraced projection is saved for diagnostics but is not the
production choice.  The Euclidean dual is
\begin{equation}
  \widetilde F_{\mu\nu}
  =\frac12\epsilon_{\mu\nu\rho\sigma}F_{\rho\sigma},
  \qquad \epsilon_{\mathrm E}^{xyzt}=+1,
  \label{eq:euclidean-dual-strength}
\end{equation}
where the letters \(x,y,z,t\), rather than the integers printed in an older
note, define the orientation.  The sign is fixed by the source routine
\texttt{levi4}; it is not a freely adjustable convention after the OPE files
have been generated.  Any comparison with Minkowski formulae must explicitly
carry the Wick-rotation phase rather than infer a sign from notation alone.

\subsection{One convention for Minkowski and Euclidean tensors}
\label{sec:gluon-wick-convention}

The unpolarized and helicity channels use one common convention throughout
this note.  Minkowski space has
\begin{equation}
 \eta_{\mu\nu}=\operatorname{diag}(1,-1,-1,-1),
 \qquad \epsilon_{\mathrm M}^{0123}=+1,
 \qquad
 \widetilde F^{\mathrm M}_{\mu\nu}
 =\frac12\epsilon^{\mathrm M}_{\mu\nu\rho\sigma}
 F_{\mathrm M}^{\rho\sigma}.
 \label{eq:minkowski-tensor-convention}
\end{equation}
The implementation orders Euclidean coordinates as
\begin{equation}
 (x,y,z,t)=(0,1,2,3),
 \qquad \epsilon_{\mathrm E}^{xyzt}=+1,
 \qquad
 \widetilde F^{\mathrm E}_{\mu\nu}
 =\frac12\epsilon^{\mathrm E}_{\mu\nu\rho\sigma}
 F^{\mathrm E}_{\rho\sigma}.
 \label{eq:euclidean-code-convention}
\end{equation}
Choose the Wick rotation
\begin{equation}
 x^0_{\mathrm M}=-\ii x^t_{\mathrm E},
 \qquad A^E_t=-\ii A^M_0.
 \label{eq:wick-rotation-choice}
\end{equation}
It gives, for spatial indices \(k,l\),
\begin{align}
 F^E_{tk}&=-\ii F^M_{0k},
 &F^E_{kl}&=F^M_{kl},
 \label{eq:wick-field-strength}\\
 \widetilde F^M_{0k}&=\widetilde F^E_{tk},
 &\widetilde F^M_{kl}&=-\ii\widetilde F^E_{kl}.
 \label{eq:wick-dual-strength}
\end{align}
Thus an electric--electric product and an electric--dual-electric product do
not acquire the same Wick phase.  This is the essential point that is lost if
one merely replaces the Minkowski label \(0\) by the Euclidean label \(t\).

The source routine \texttt{dual\_field} evaluates
Eq.~\eqref{eq:euclidean-code-convention} directly.  In particular, its signs
are
\begin{equation}
 \widetilde F^E_{tx}=-F^E_{yz},
 \qquad
 \widetilde F^E_{ty}=+F^E_{xz},
 \qquad
 \widetilde F^E_{xy}=+F^E_{zt}.
 \label{eq:code-dual-components}
\end{equation}
Equations~\eqref{eq:wick-field-strength}--\eqref{eq:code-dual-components}
provide an executable sign dictionary: the code stores the Euclidean Clover
tensor and does not attach an additional temporal-link phase.  Any comparison
with a Minkowski operator must therefore apply these equations to every
factor.

\subsubsection{Dictionary for the old HYP note}

The internal file \texttt{DeltaG\_Note.pdf} wrote analytic formulae in the
order
\begin{equation}
 (0,1,2,3)_{\rm old}=(t,x,y,z),
 \qquad \epsilon_{\rm old}^{txyz}=+1,
 \label{eq:old-note-index-order}
\end{equation}
whereas both its Python implementation and the present C++ implementation use
Eq.~\eqref{eq:euclidean-code-convention}.  The component relabelling is
\begin{center}
\small
\begin{tabular}{c@{\(\;\longrightarrow\;\)}c c@{\qquad}c@{\(\;\longrightarrow\;\)}c}
\toprule
old note & current code & & old note & current code\\
\midrule
\(F_{01}\) & \(F_{tx}=F_{30}\) && \(F_{23}\) & \(F_{yz}=F_{12}\)\\
\(F_{02}\) & \(F_{ty}=F_{31}\) && \(F_{13}\) & \(F_{xz}=F_{02}\)\\
\(F_{12}\) & \(F_{xy}=F_{01}\) && \(F_{03}\) & \(F_{tz}=F_{32}\)\\
\bottomrule
\end{tabular}
\end{center}
Moving \(t\) from the first to the last position is an odd permutation;
therefore Eq.~\eqref{eq:old-note-index-order}, when merely relabelled as
\(xyzt\), has the opposite orientation to the actual code tensor.  A dual
field acquires this orientation sign, while an ordinary \(F_{\mu\nu}\) does
not.  Consequently the component expansions in Eqs.~(35)--(36) of the old
note cannot be copied term by term into the current basis.  In what follows,
``old plus'' and ``old temporal'' mean diagnostic channels inspired by those
formulae, but they are redefined unambiguously from the current \(H^E_{\mu\nu}\).

\subsection{Gauge-invariant flowed bilocal}

For a separation in the third spatial direction, define
\begin{align}
  M_{\mu\lambda;\nu\rho}^{(\tau)}(z,t)
  ={}&\sum_{\bm x}\Tr\bigl[
  F_{\mu\lambda}(\tau,\bm x+z\hat{\bm z},t)
  W_\tau(\bm x+z\hat{\bm z},\bm x;t)\nonumber\\
  &\hspace{2.6cm}\times
  F_{\nu\rho}(\tau,\bm x,t)
  W_\tau(\bm x,\bm x+z\hat{\bm z};t)
  \bigr].
  \label{eq:flowed-gluon-bilocal}
\end{align}
Both field strengths and every link of the straight Wilson line are made from
the same flowed gauge field at the same \(\tau\).  The trace is in the
fundamental representation.  Gauge transformations at the two endpoints
cancel inside the trace, so Eq.~\eqref{eq:flowed-gluon-bilocal} is gauge
invariant without gauge fixing.

The code stores both line orientations and forms
\begin{align}
  M^{\rm even}(z)&=M(+z)+M(-z),\\
  M^{\rm odd}(z)&=M(+z)-M(-z).
  \label{eq:gluon-z-projections}
\end{align}
These are a sum and difference, not averages: there is no factor \(1/2\) in
the stored arrays.  Keeping both raw orientations is important for testing
Hermiticity, parity, and the exact linear identities.

To put the two physics channels on identical footing, define the Euclidean
primitive correlators
\begin{align}
 U^E_{\mu\nu}(z)
 &=\sum_{\bm x}\Tr\!\left[
 F^E_{\mu\nu}(x+z)W^\dagger
 F^E_{\mu\nu}(x)W\right],
 \label{eq:euclidean-unpolarized-primitive}\\
 H^E_{\mu\nu}(z)
 &=\sum_{\bm x}\Tr\!\left[
 F^E_{\mu\nu}(x+z)W^\dagger
 \widetilde F^E_{\mu\nu}(x)W\right].
 \label{eq:euclidean-helicity-primitive}
\end{align}
The flow-time and time-slice labels are suppressed.  In both equations the
first tensor is evaluated at the displaced endpoint and the second at the
origin, exactly as in the implementation.  The only channel difference is
\(F\) versus \(\widetilde F\) in the second factor.

\subsection{Unpolarized operator}

For a boost and separation along \(z\), the transverse sum means
\(\sum_{\perp=x,y}\); a repeated \(i\) below never includes the longitudinal
direction.  Start from the Minkowski pseudo-PDF current
\cite{KhanEtAl2021UnpolarizedGluon}
\begin{align}
 \mathcal M_U^M
 &=\sum_{i=x,y}M^M_{0i;i0}+2M^M_{xy;yx},
 \label{eq:unpolarized-minkowski-start}\\
 M^M_{\mu\nu;\rho\sigma}
 &\equiv
 \Tr\!\left[F^M_{\mu\nu}(z)W^\dagger
               F^M_{\rho\sigma}(0)W\right].
\end{align}
Antisymmetry gives \(F_{i0}=-F_{0i}\) and \(F_{yx}=-F_{xy}\).
The Wick dictionary gives
\(F^M_{0i}=+\ii F^E_{ti}\) and leaves spatial components unchanged.
Consequently, term by term,
\begin{align}
 M^M_{0i;i0}
 &=\Tr\!\left[(+\ii F^E_{ti})W^\dagger
                 (-\ii F^E_{ti})W\right]
   =U^E_{ti},\\
 2M^M_{xy;yx}
 &=2\Tr\!\left[F^E_{xy}W^\dagger(-F^E_{xy})W\right]
   =-2U^E_{xy}.
 \label{eq:unpolarized-termwise-wick}
\end{align}
Therefore
\begin{equation}
 \mathcal M_U^M
 =U^E_{tx}+U^E_{ty}-2U^E_{xy}.
 \label{eq:unpolarized-wick-result}
\end{equation}
This explicit reversed-index derivation prevents the overall sign ambiguity
in Eq.~(25) of the old note.  After translating that equation to the current
basis, its printed Euclidean expression is
\(-T_U+S_U=-(T_U-S_U)\): it differs from the present convention only by an
overall operator sign, not by a relative-sign ambiguity.  We fix the
production convention to the right-hand side of
Eq.~\eqref{eq:unpolarized-wick-result}.  With code indices
\((x,y,z,t)=(0,1,2,3)\), define
\begin{align}
  \cO_U^{(\tau)}(z)
  &=U^E_{tx}(z)+U^E_{ty}(z)-2U^E_{xy}(z)\nonumber\\
  &\equiv T_U(z)-S_U(z),
  \label{eq:flowed-unpolarized-operator}\\
  T_U&\equiv U^E_{tx}+U^E_{ty},
  &S_U&\equiv2U^E_{xy}.
\end{align}
This is exactly the index structure
\(M_{tx;tx}+M_{ty;ty}-2M_{xy;xy}\) used in
Ref.~\cite{ChenEtAl2025GluonPDF}.  The production physics selector uses its
even-separation part.  The code saves \texttt{combined}, \(M_{ti;ti}\), and
\(M_{ij;ij}\) separately; in the unified notation these are
\(T_U-S_U\), \(T_U\), and \(S_U\), so the identity can be verified configuration by
configuration.

The statement that an analogous thin-link operator belongs to a convenient
multiplicative-renormalization class does not prove that a finite-flow
matrix element already has the paper's hybrid-renormalized normalization.
The flowed operator needs its own flow-scheme conversion and mixing analysis.

\subsection{Helicity operator, odd separation, and Euclidean phase}

The polarized continuum bilocal is schematically
\begin{equation}
  \widetilde m_{\mu\alpha;\lambda\beta}(z,P,S)
  =\langle P,S|G_{\mu\alpha}(z)W(z,0)
      \widetilde G_{\lambda\beta}(0)|P,S\rangle,
  \label{eq:gluon-helicity-bilocal}
\end{equation}
and its spin-dependent contribution is odd under reversal of the separation,
\begin{equation}
  \widetilde M(z,P,S)=\widetilde m(z,P,S)
                     -\widetilde m(-z,P,S).
  \label{eq:helicity-odd-z}
\end{equation}
The Minkowski target used for the helicity pseudo-distribution is
\cite{BalitskyMorrisRadyushkin2022,EgererEtAl2022GluonHelicity}
\begin{equation}
 \mathcal M_{\Delta g}^M
 =\sum_{i=x,y}\widetilde M^M_{0i;0i}
  +2\widetilde M^M_{xy;xy}.
 \label{eq:minkowski-helicity-target}
\end{equation}
Here the second field strength in every term is dualized.  It is useful first
to expose the electric--magnetic content.  With
\begin{equation}
 E_i\equiv F^M_{0i},\qquad
 F^M_{ij}=-\epsilon_{ijk}B_k,\qquad
 \widetilde F^M_{0i}=B_i,\qquad
 \widetilde F^M_{ij}=\epsilon_{ijk}E_k,
 \label{eq:minkowski-electric-magnetic}
\end{equation}
the target is
\begin{equation}
 \mathcal M_{\Delta g}^M
 \sim E_xB_x+E_yB_y-2B_zE_z.
 \label{eq:helicity-eb-check}
\end{equation}
The longitudinal term already has the relative minus sign; it arises from
\(F^M_{xy}=-B_z\), not from a numerical preference of the data.
The source-code dual identities in
Eq.~\eqref{eq:code-dual-components} provide an immediate independent check:
\begin{align}
 T_H
 &=-\ii\left(E_xB_x+E_yB_y\right),&
 S_H&=-2\ii B_zE_z,\nonumber\\
 \ii(T_H-S_H)
 &=E_xB_x+E_yB_y-2B_zE_z.
 \label{eq:code-helicity-eb-map}
\end{align}
Thus the target relative sign is already visible without manipulating any
four-index epsilon tensor.

The same conclusion follows directly from Wick rotation.  From
Eq.~\eqref{eq:wick-field-strength},
\(F^M_{0i}=+\ii F^E_{ti}\), while
Eq.~\eqref{eq:wick-dual-strength} gives
\(\widetilde F^M_{0i}=\widetilde F^E_{ti}\),
\(F^M_{xy}=F^E_{xy}\), and
\(\widetilde F^M_{xy}=-\ii\widetilde F^E_{xy}\).  Hence
\begin{align}
 \widetilde M^M_{0i;0i}
 &=\Tr\!\left[(+\ii F^E_{ti})W^\dagger
                  \widetilde F^E_{ti}W\right]
   =+\ii H^E_{ti},\\
 2\widetilde M^M_{xy;xy}
 &=2\Tr\!\left[F^E_{xy}W^\dagger
                 (-\ii\widetilde F^E_{xy})W\right]
   =-2\ii H^E_{xy},
 \label{eq:helicity-termwise-wick}\\
 \mathcal M_{\Delta g}^M
 &=+\ii\left[H^E_{tx}+H^E_{ty}-2H^E_{xy}\right].
 \label{eq:helicity-wick-result}
\end{align}
The two derivations agree: the temporal and spatial factors acquire opposite
Wick phases.  Removing the common factor \(\ii\) yields the code-basis target
\begin{align}
 \cO_{\Delta g}^{E,(\tau)}(z)
 &=H^E_{tx}(z)+H^E_{ty}(z)-2H^E_{xy}(z)\nonumber\\
 &\equiv T_H(z)-S_H(z),
 \label{eq:flowed-helicity-operator}\\
 T_H&\equiv H^E_{tx}+H^E_{ty},
 &S_H&\equiv2H^E_{xy}.
\end{align}
Its physical orientation is the unaveraged odd difference
\begin{equation}
 \cO_{\Delta g}^{E,\mathrm{odd}}(z)
 =\cO_{\Delta g}^{E}(+z)-\cO_{\Delta g}^{E}(-z).
 \label{eq:helicity-correct-odd-combination}
\end{equation}
Thus unpolarized and helicity use the same Euclidean coefficient vector
\((+1,+1,-2)\), or \(T_X-S_X\).  Their differences are the
second tensor (\(F\) or \(\widetilde F\)), even versus odd separation, and
the nucleon spin projector.  A different Wick direction or an overall dual
orientation flips the entire helicity operator, but cannot change the
relative sign between \(T_H\) and \(S_H\).

The old note's Eq.~(35), after its \(txyz\to xyzt\) orientation is treated
consistently, maps to the diagnostic \(T_H+S_H\), while its Eq.~(36) maps to
\(T_H\).  They are useful legacy cross-checks but do not replace the target
derived in Eqs.~\eqref{eq:helicity-eb-check} and
\eqref{eq:flowed-helicity-operator}.
For complete transparency, the actual code dual
Eq.~\eqref{eq:code-dual-components} expands the three choices as
\begin{align}
 T_H-S_H
 &=-K_{tx;yz}+K_{ty;xz}+2K_{xy;tz},\nonumber\\
 T_H+S_H
 &=-K_{tx;yz}+K_{ty;xz}-2K_{xy;tz},\nonumber\\
 T_H
 &=-K_{tx;yz}+K_{ty;xz},
 \label{eq:helicity-three-component-expansion}
\end{align}
where, only in this equation,
\[
 K_{\mu\nu;\rho\sigma}(z)
 \equiv\sum_{\bm x}\Tr\!\left[
 F^E_{\mu\nu}(x+z)W^\dagger F^E_{\rho\sigma}(x)W\right].
\]
These signs reproduce the individual contractions in the old Eq.~(35) after
the index dictionary is applied and make the sole disputed longitudinal term
explicit.

For a boost along \(z\), Ref.~\cite{EgererEtAl2022GluonHelicity} relates this
combination to invariant amplitudes by
\begin{equation}
  \widetilde M_{0i;0i}+\widetilde M_{ij;ij}
  =-2P_zP_0\widetilde{\mathcal M}_{sp}^{(+)}(\nu,z^2)
   +2P_0^3z\widetilde{\mathcal M}_{pp}(\nu,z^2).
  \label{eq:helicity-invariant-amplitudes}
\end{equation}
The second term produces a target-mass contamination of order
\(m_N^2/P_z^2\) in the reduced pseudo-ITD construction.  It must be fitted or
subtracted; it is not removed by gradient flow.

The Euclidean helicity correlator is naturally complex.  In the present gamma
basis, the numerator uses the complete complex `pol35' two-point function;
it does \emph{not} use only its imaginary part.  No real/imaginary projection
is allowed before vacuum subtraction, ratio formation, and resampling.  The
operator-level Wick factor in Eq.~\eqref{eq:helicity-wick-result} and the
phase of the `pol35' nucleon projector are logically separate conventions.
After correcting the operator combination, the physical ratio component must
be calibrated from the full projector and gamma-matrix definitions (and
checked under spin and momentum reversal), rather than inherited from the
legacy plus-sign plot.

\subsection{Three helicity selectors as a controlled diagnostic}

The independently saved raw groups permit the following comparison without
repeating Wilson flow or the OPE contraction:
\begin{center}
\begin{tabular}{
  >{\raggedright\arraybackslash}p{0.24\textwidth}
  >{\raggedright\arraybackslash}p{0.27\textwidth}
  >{\raggedright\arraybackslash}p{0.38\textwidth}}
\toprule
Label & Current-basis definition & Interpretation\\
\midrule
\texttt{TH\_minus\_SH}
 & \((T_H-S_H)_{\rm odd}\)
 & Theory target from the termwise Wick rotation and the \(E\)--\(B\) check.\\
\texttt{TH\_plus\_SH}
 & \((T_H+S_H)_{\rm odd}\)
 & Legacy full combination; also the current OPE file's stored
   \texttt{combined} helicity channel.\\
\texttt{TH\_only}
 & \((T_H)_{\rm odd}\)
 & Temporal-only legacy diagnostic, corresponding to the intent of
   Eq.~(36) in the old note.\\
\bottomrule
\end{tabular}
\end{center}
All three were recomputed at
\(\tau/a^2=0.5\), \(z/a=2\), and \(P_z=3,4,5\) from the same 231
configurations.  Each selector uses the complete complex \texttt{pol35}
correlator in the numerator, \texttt{nopol} in the denominator, and the same
delete-one samples.  Vacuum subtraction and the denominator are recomputed
inside every sample.  Thus differences between the curves retain the
\(T_H\)--\(S_H\) covariance.

\begin{figure}[htbp]
  \centering
  \includegraphics[width=0.98\textwidth]{
    figures/helicity_three_operators_tau0p500_N231_z2_tsep8_P345.pdf}
  \caption{Insertion-time comparison of the three helicity selectors at
  \(t_{\rm sep}/a=8\).  The displayed quantity is
  \(\operatorname{Im}R_{\Delta g}=\operatorname{Re}[-\ii R_{\Delta g}]\);
  the saved arrays remain complex.  The plus channel reproduces the previous
  ratio-v2 product to \(3.7\times10^{-14}\) in the central value and
  \(5.4\times10^{-15}\) in its jackknife errors.  Endpoints are shown for
  diagnosis and are not included in a plateau fit.}
  \label{fig:helicity-three-selectors-insertion}
\end{figure}

\begin{figure}[htbp]
  \centering
  \includegraphics[width=0.98\textwidth]{
    figures/helicity_three_operators_tau0p500_N231_z2_midpoint_tsep5to11_P345.pdf}
  \caption{One- or two-slice midpoint averages for the three selectors.
  Only \(t_{\rm sep}/a\leq11\) is drawn so that the rapidly growing noise at
  larger separations does not hide the resolved short-separation comparison.
  These points are correlated diagnostics, not ground-state fits.}
  \label{fig:helicity-three-selectors-midpoint}
\end{figure}

At \(t_{\rm sep}/a=8\), the midpoint values for
\((T_H-S_H)_{\rm odd}\) are approximately
\(0.005,-0.002,-0.033\), whereas those for
\((T_H+S_H)_{\rm odd}\) are approximately
\(-0.045,-0.051,-0.053\), for \(P_z=3,4,5\), respectively.
The temporal-only values lie between them.  This demonstrates a sizable,
strongly correlated cancellation in the minus channel.  It does
\emph{not} select the plus channel as the physical operator: statistical
smoothness or signal size cannot override the tensor definition.

\subsection{Audit of the current OPE implementation}

The relevant source is \texttt{src/flowed\_gluon\_ope.cpp}.  The comparison
with the common convention above is
\begin{center}
\begin{tabular}{
  >{\raggedright\arraybackslash}p{0.31\textwidth}
  >{\raggedright\arraybackslash}p{0.18\textwidth}
  >{\raggedright\arraybackslash}p{0.42\textwidth}}
\toprule
Implementation item & Verdict & Reason \\
\midrule
Clover \(-\ii(Q-Q^\dagger)/8\) and traceless projection
 & consistent & Euclidean field is stored without an extra temporal phase.\\
\(\epsilon_E^{xyzt}=+1\) and \texttt{dual\_field}
 & consistent & It reproduces Eq.~\eqref{eq:code-dual-components}.\\
Unpolarized raw terms and \texttt{combined}
 & consistent & The code uses \(M_{ti;ti}^{U}-M_{ij;ij}^{U}\).\\
Helicity raw \(M_{ti;ti}^{H}\), \(M_{ij;ij}^{H}\)
 & consistent & They are the required \(F W\widetilde F\) contractions.\\
Helicity \texttt{odd\_difference}
 & consistent & It is stored as the unaveraged \(+z\) minus \(-z\) value.\\
Helicity \texttt{combined}
 & inconsistent & The present ternary uses
 \texttt{mtiti+mijij}; Eqs.~\eqref{eq:helicity-wick-result} and
 \eqref{eq:flowed-helicity-operator} require
 \texttt{mtiti-mijij}.\\
\bottomrule
\end{tabular}
\end{center}
The corresponding JSON metadata statements
`\texttt{helicity\_combined\_euclidean=Mtiti+Mijij}' and the plus-sign paper
target must also be corrected when the writer is updated.  Because every OPE
file already stores the two raw terms independently, no gauge-flow or OPE
rerun is required.  The three-selector diagnostic above reruns the helicity
disconnected ratio and delete-one jackknife from those raw terms.  It does
not overwrite the legacy ratio-v2 files.  Subtracting published means or
combining their errors in quadrature would discard the configuration-level
covariance and is not used.

\subsection{All-operator $C_3/C_2$ bootstrap product}
\label{sec:all-operator-c3c2}

The raw OPE files already contain the two independent tensor pieces needed by
both channels.  Therefore the requested three operator choices can be formed
without rerunning Wilson flow or the gluon OPE contraction.  The construction
uses the same configuration-level covariance estimator as
Eq.~\eqref{eq:flowed-gluon-disconnected-c3}, and reconstructs all six choices
from raw `Mtiti' and `Mijij'; the legacy `combined' helicity field is used only
to verify the source identity, not as a definition of the new product.

With $T_U=U_{tx}+U_{ty}$, $S_U=2U_{xy}$ and
$T_H=H_{tx}+H_{ty}$, $S_H=2H_{xy}$, the operator axis is ordered as follows:
\begin{center}
\begin{tabular}{p{0.24\textwidth}p{0.22\textwidth}p{0.42\textwidth}}
\toprule
Channel & Choice & Operator and $z$ projection\\
\midrule
Unpolarized & no $S_U$ & $T_U$, even in $z$\\
Unpolarized & $+S_U$ & $T_U+S_U$, even in $z$\\
Unpolarized & $-S_U$ & $T_U-S_U$, even in $z$\\
Helicity & no $S_H$ & $T_H$, odd in $z$\\
Helicity & $+S_H$ & $T_H+S_H$, odd in $z$\\
Helicity & $-S_H$ & $T_H-S_H$, odd in $z$\\
\bottomrule
\end{tabular}
\end{center}
The corresponding labels, in this order, are
\begin{center}
\begin{tabular}{p{0.43\textwidth}p{0.43\textwidth}}
\texttt{unpolarized\_TU\_even} &
\texttt{unpolarized\_TU\_plus\_SU\_even}\\
\texttt{unpolarized\_TU\_minus\_SU\_even} &
\texttt{helicity\_TH\_odd}\\
\texttt{helicity\_TH\_plus\_SH\_odd} &
\texttt{helicity\_TH\_minus\_SH\_odd}
\end{tabular}
\end{center}
The first three use the real
projection of the C3 numerator; the last three use the imaginary projection
only after forming the full complex helicity numerator.  In all six cases the
denominator is the real part of the unpolarized `nopol' two-point function;
`pol35' is used only for the helicity C3 covariance.

The new product is stored under
\begin{center}
\nolinkurl{/public/group/lqcd/donghx/Gradient_Flow_GluonPDF/ratio_bootstrap_all_operators_v1/results_v1/}.
\end{center}
It uses the frozen common manifest of $N=406$ configurations, 14 flow times
$\tau/a^2=0.1,0.2,0.3,0.4,0.5,0.6,1.0,1.4,1.8,2.2,2.6,3.0,3.4,3.8$,
$z/a=0,\ldots,24$, $t_{\rm sep}/a=5,\ldots,15$, cuts $1,2,3$,
$|P_z|=3,4,5$, the $+z$ Wilson-line direction, and the color-traceless
projection.  Each flow-time file contains
\begin{align}
 &C_3^{\rm pointwise},R^{\rm pointwise}:(6,25,11,16,3),
 &C_2^{\rm original}:(2,11,3),\\
 &C_{3,\rm bootstrap}^{\rm sum},R_{\rm bootstrap}^{\rm sum}
   :(1000,6,3,25,11,3),
 &C_2^{\rm bootstrap}:(1000,2,11,3).
\end{align}
Here the bootstrap axis contains exactly 1000 multinomial configuration
replicas; the deterministic central average is stored separately.  Vacuum
subtraction and the C2 denominator are recomputed in every replica, and the
seed is 1115.  The older N=406/Nboot=1500 ratio product and N=231 three-selector figures remain immutable
provenance products.  These new arrays are finite-flow bare disconnected
ratios only: no state isolation, flow matching, renormalization, mixing,
continuum limit, or PDF reconstruction has been applied.

\subsection{Disconnected three-point function and vacuum subtraction}

The gluon loop has no valence quark line joining it to the nucleon.  At fixed
source time \(t_0\), separation \(t_{\rm sep}\), and insertion offset
\(t_{\rm ins}\), the connected-in-the-statistical-sense disconnected
correlator is the gauge-ensemble covariance
\begin{align}
 C_{3,X}
 ={}&\left\langle
  O_X(t_0+t_{\rm ins})C_{2,\Pi_X}(t_0,t_0+t_{\rm sep})
 \right\rangle_{\rm cfg}\nonumber\\
 &-\left\langle O_X(t_0+t_{\rm ins})\right\rangle_{\rm cfg}
  \left\langle C_{2,\Pi_X}(t_0,t_0+t_{\rm sep})
  \right\rangle_{\rm cfg},
 \label{eq:flowed-gluon-disconnected-c3}
\end{align}
with
\begin{equation}
  \Pi_U={\tt nopol},
  \qquad
  \Pi_{\Delta g}={\tt pol35}.
  \label{eq:gluon-numerator-projectors}
\end{equation}
This subtraction removes the factorized vacuum expectation value.  Subtracting
only the loop mean or multiplying separately averaged objects after source
averaging is not generally the same estimator.

The production ratio is
\begin{equation}
 R_X(z,t_{\rm sep},t_{\rm ins};P_z,\tau)
 =\frac{\operatorname{mean}_{t_0} C_{3,X}}
        {\operatorname{mean}_{t_0}
          \langle C_{2,{\tt nopol}}\rangle_{\rm cfg}}.
 \label{eq:flowed-gluon-ratio}
\end{equation}
Crucially, \emph{both} channels use `nopol' in the denominator.  The
spin-projected `pol35' correlator appears only in the helicity numerator.  An
older ratio-v1 divided the helicity channel by the nearly zero ensemble mean
of `pol35', creating an artificial enhancement; those files are retained
only for provenance.

\subsection{Delete-one jackknife as part of the estimator}

For \(N\) configurations, sample \(k\) omits configuration \(k\) and
recomputes every nonlinear operation:
\begin{equation}
 R_X^{(k)}=
 \frac{\operatorname{mean}_{t_0}
   \left[\langle O_XC_{2,\Pi_X}\rangle_{(-k)}
   -\langle O_X\rangle_{(-k)}
    \langle C_{2,\Pi_X}\rangle_{(-k)}\right]}
 {\operatorname{mean}_{t_0}
    \langle C_{2,{\tt nopol}}\rangle_{(-k)}}.
 \label{eq:gluon-jackknife-estimator}
\end{equation}
For a real-valued component \(q\),
\begin{equation}
 \sigma_q^2=\frac{N-1}{N}\sum_{k=1}^{N}
  \left(q^{(k)}-\overline q_{\rm JK}\right)^2.
 \label{eq:gluon-jackknife-error}
\end{equation}
The implementation applies Eq.~\eqref{eq:gluon-jackknife-error} separately
to real and imaginary parts while retaining the complex ratio.  Reusing a
full-sample vacuum subtraction or denominator in every leave-one-out sample
would underestimate or distort the propagated covariance.

\subsection{State isolation: plateau, summation, and Feynman--Hellmann form}

Away from source and sink contact regions, the ratio has the spectral form
\begin{align}
 R_X(t_{\rm sep},t_{\rm ins})={}&h_X
 +A_Xe^{-\Delta E t_{\rm ins}}
 +B_Xe^{-\Delta E(t_{\rm sep}-t_{\rm ins})}
 +C_Xe^{-\Delta E t_{\rm sep}}+\cdots .
 \label{eq:gluon-ratio-spectral}
\end{align}
A plateau fit therefore requires both endpoint cuts and stability in
\(t_{\rm sep}\).  An alternative used in
Ref.~\cite{ChenEtAl2025GluonPDF} is the summed-difference or
Feynman--Hellmann estimator
\begin{align}
 \Delta R_X(t_{\rm sep})={}&
 \frac{\sum_{t=t_r}^{t_{\rm sep}+1-t_r}
 C_{3,X}(t_{\rm sep}+1,t)}{C_2(t_{\rm sep}+1)}
 -\frac{\sum_{t=t_r}^{t_{\rm sep}-t_r}
 C_{3,X}(t_{\rm sep},t)}{C_2(t_{\rm sep})},
 \label{eq:gluon-fh-ratio}\\
 \Delta R_X(t_{\rm sep})\simeq{}&h_X+c_1t_{\rm sep}
 e^{-\Delta E t_{\rm sep}}.
 \label{eq:gluon-fh-fit}
\end{align}
The current finite-flow bare gluon C3/C2 products precede this state-isolation layer.  A visible
interior signal is not yet a fitted matrix element.

\subsection{Flow-time window and order of limits}

The completed dimensionless flow times and radii are
\begin{center}
\small
\begin{tabular}{c@{\quad}c|c@{\quad}c|c@{\quad}c}
\toprule
\(\tau/a^2\)&\(r_{\rm flow}/a\)&\(\tau/a^2\)&\(r_{\rm flow}/a\)&
\(\tau/a^2\)&\(r_{\rm flow}/a\)\\
\midrule
0.1&0.894&0.2&1.265&0.3&1.549\\
0.4&1.789&0.5&2.000&0.6&2.191\\
1.0&2.828&1.4&3.347&1.8&3.795\\
2.2&4.195&2.6&4.561&3.0&4.899\\
3.4&5.215&3.8&5.514&&\\
\bottomrule
\end{tabular}
\end{center}
The useful window must suppress cutoff artifacts while avoiding substantial
distortion of the bilocal and hadron scales.  Indicative requirements are
\begin{equation}
  a\ll r_{\rm flow},
  \qquad r_{\rm flow}\ll |z|,\Lambda_{\rm QCD}^{-1},
  \qquad P_z r_{\rm flow}\ll1,
  \label{eq:gluon-small-flow-window}
\end{equation}
when a small-flow expansion is invoked.  They cannot all be imposed at
\(z=0\), and many large-\(\tau\), small-\(z\) points in the present grid do
not satisfy them parametrically.  Those points remain useful for diagnostics
but cannot be silently included in a small-flow fit.

A controlled program takes the continuum limit at fixed physical flow time
or otherwise specifies a joint \((a,\tau)\) trajectory, and only then performs
a justified flow-scheme or small-flow conversion.  Sending \(\tau/a^2\) to
zero on one lattice is not a continuum extrapolation.  Conversely, holding
\(\tau/a^2\) fixed as \(a\to0\) sends the physical smoothing radius to zero
and can reintroduce cutoff sensitivity.

\subsection{Helicity reduced pseudo-ITD and mixing caveats}

For orientation, the helicity reference defines a reduced pseudo-ITD with an
explicit phase,
\begin{equation}
 \widetilde{\mathfrak M}(\nu,z^2)
 =\ii\,
 \frac{[\widetilde{\mathcal M}_{00}(z,P_z)/(P_zP_0)]/Z_L(z/a)}
      {\mathcal M_{00}(z,P_z=0)/m_N^2}.
 \label{eq:helicity-reduced-pitd}
\end{equation}
Its invariant-amplitude content includes
\begin{equation}
 \widetilde{\mathfrak M}
 =\left[\widetilde{\mathcal M}_{sp}^{(+)}
       -\nu\widetilde{\mathcal M}_{pp}\right]
 -\frac{m_N^2}{P_z^2}\nu\widetilde{\mathcal M}_{pp}.
 \label{eq:helicity-target-mass-term}
\end{equation}
Equations~\eqref{eq:helicity-reduced-pitd}--
\eqref{eq:helicity-target-mass-term} specify the reduced-pseudo-ITD phase and
the need for a target-mass analysis, but they are not yet implemented in the
finite-flow bare C3/C2 products.  This normalization phase must be applied only after the corrected
Euclidean operator and the complete complex correlator have been formed.
Furthermore, perturbative evolution/matching of a gluon observable couples to
the flavor-singlet quark distribution.  Omitting that mixing is an explicit
approximation, not a property of gradient flow.

\subsection{Current \texorpdfstring{\(32^3\times96\)}{32x96} production}

The OPE products are under
\begin{center}
\nolinkurl{/public/group/lqcd/donghx/Gradient_Flow_GluonPDF/output_v5},
\end{center}
and the all-operator C3/C2 products are under
\begin{center}
\nolinkurl{/public/group/lqcd/donghx/Gradient_Flow_GluonPDF/ratio_bootstrap_all_operators_v1/results_v1}.
\end{center}
The new all-operator Slurm array (job 3575896) constructs six
covariance-preserving choices at all 14 flow times and resamples each to
Nboot=1000.  The files use one frozen common set of 406 configurations with
manifest SHA-256
\begin{center}
\ttfamily\small
1a2161d754493f03a4c5e26d5fe20a0a83446f762f3a8bc159cc282f866b53e.
\end{center}
The older ratio-v2 products and the N=231 three-selector figures remain
legacy provenance.  The collection validator now reports
\texttt{complete\_and\_validated}; its receipt and SHA-256 are
\begin{center}
\ttfamily\small
ratio\_bootstrap\_all\_operators\_v1/results\_v1/collection\_validation.json\\
ef79c67b7dbaf3ab895a4d7bdec53c70b9ce9260a8c622f84f00f8ed8c81e8b0.
\end{center}
This establishes product and schema completion, not ground-state isolation,
renormalization, matching, or PDF correctness.

The current kinematics and storage contract are
\begin{center}
\begin{tabular}{>{\raggedright\arraybackslash}p{0.29\textwidth}>{\raggedright\arraybackslash}p{0.62\textwidth}}
\toprule
Item & Value\\
\midrule
Gauge input & thin links, no HYP or other pre-smearing\\
Flow & Wilson flow, RK3, \(\epsilon=0.01\)\\
Momentum & \(P_z=+3,+4,+5\); only the \(+z\) direction\\
Separations & \(z/a=0,\ldots,24\), \(t_{\rm sep}/a=5,\ldots,15\)\\
OPE projection & color-traceless\\
Saved components & combined, \(M_{ti;ti}\), \(M_{ij;ij}\); new products
recombine raw \texttt{Mtiti}/\texttt{Mijij}\\
Operator choices & in each channel: no \(S\), \(T+S\), and \(T-S\),
with even/odd $z$ projection as appropriate\\
Ratio schema & {\scriptsize\texttt{gradient\_flow\_gluon\_c3\_c2\_bootstrap\_all\_operators\_v1}}\\
Status & finite-flow bare disconnected C3/C2; six operators, N=406,
Nboot=1000, all 14 flow times\\
\bottomrule
\end{tabular}
\end{center}
Only one Wilson-line direction was generated, so the result must not be
described as a six-direction average.

One OPE file has axes
\begin{equation}
 (\text{field projection},\text{channel},\text{orientation},
  \text{component},z,t)=(2,2,4,6,25,96),
 \label{eq:gluon-ope-shape}
\end{equation}
The legacy ratio-v2 file has axes
\begin{align}
 &(\text{channel},\text{orientation},\text{component},\text{direction},
 z,t_{\rm sep},t_{\rm ins},|P_z|)\nonumber\\
 &\hspace{3cm}=(2,4,3,1,25,11,16,3).
 \label{eq:gluon-ratio-shape}
\end{align}
The new all-operator pointwise file instead has axes
\begin{equation}
 (\text{operator},z,t_{\rm sep},t_{\rm ins},|P_z|)=(6,25,11,16,3),
 \qquad
 (\text{channel},t_{\rm sep},|P_z|)=(2,11,3),
 \label{eq:gluon-all-operator-shape}
\end{equation}
and its bootstrap summed ratio has axes
\begin{equation}
 (\text{bootstrap},\text{operator},\text{cut},z,t_{\rm sep},|P_z|)
 =(1000,6,3,25,11,3).
 \label{eq:gluon-all-operator-bootstrap-shape}
\end{equation}
The valid-insertion and valid-cut masks are mandatory; padded entries are not
measurements.

\begin{figure}[htbp]
  \centering
  \includegraphics[width=0.98\textwidth]{figures/ratio_plateau_tau0p500_z2_P345.pdf}
  \caption{Diagnostic finite-flow bare ratios at \(\tau/a^2=0.5\) and
  \(z/a=2\).  The unpolarized plot uses the even `combined' channel.  The
  helicity panel is retained only as a legacy sign diagnostic: it used the
  incorrect odd \(M_{ti;ti}^{H}+M_{ij;ij}^{H}\) combination and is not the
  target in Eq.~\eqref{eq:flowed-helicity-operator}.  Endpoint values can contain
  contact effects; no plateau, summed-ratio, Feynman--Hellmann, or multi-state
  fit is represented by these points.}
  \label{fig:flowed-gluon-ratio-diagnostic}
\end{figure}

At this diagnostic point the unpolarized interior is approximately \(-0.5\).
The legacy helicity panel gave values near \(-0.08\) to \(-0.10\).  The
covariance-preserving comparison in
Figs.~\ref{fig:helicity-three-selectors-insertion} and
\ref{fig:helicity-three-selectors-midpoint} confirms substantial
\(T_H-S_H\) cancellation at small source--sink separations.  No statement
about a helicity plateau should be made before state-isolation fits and
spin/momentum-reversal phase checks.

\subsection{Required path from finite-flow bare C3/C2 to a gluon PDF}

The remaining physics chain is:
\begin{enumerate}
  \item impose source/sink endpoint cuts and extract the ground-state matrix
  element with plateau, summed/Feynman--Hellmann, and preferably multi-state
  stability comparisons;
  \item propagate correlated jackknife samples through every fit and compare
  fit windows without point-by-point optimization;
  \item select a controlled flow-time window and perform continuum analysis
  at a stated physical flow prescription;
  \item determine the finite-flow operator normalization and flow-scheme
  conversion, including gluon--singlet-quark mixing;
  \item control target-mass and higher-twist corrections and test momentum
  dependence;
  \item apply the appropriate LaMET or pseudo-PDF matching kernel and scale
  evolution;
  \item reconstruct the finite-\(z\) data in \(x\), with inverse-problem and
  truncation systematics.
\end{enumerate}
Until these steps are complete, the correct label is
``finite-flow bare disconnected gluon ratio'', not ``gluon PDF''.

\section{Connected nucleon contractions}

\subsection{Proton interpolator and spin projection}

A local proton interpolator may be written schematically as
\begin{equation}
  N_\alpha(x)
  =
  \epsilon^{abc}
  \left[u^{aT}(x)\,C\gamma_5\gamma_4\,d^b(x)\right]
  u^c_\alpha(x),
  \label{eq:proton-interpolator}
\end{equation}
with color and space represented by distillation tensors. In the implemented
Euclidean basis,
\begin{align}
  C&=\gamma_4\gamma_2,
  &
  \gamma_5&=\gamma_1\gamma_2\gamma_3\gamma_4,
  \\
  \Gamma_{\rm diq}
  &=C\gamma_5\gamma_4=(\gamma_3\gamma_1)\gamma_4,
  \\
  \Pi_0
  &=\frac{1+\gamma_4}{2},
  &
  \Pi_k
  &=\Pi_0\,\ii\gamma_k\gamma_5,
  \qquad k=1,2,3.
  \label{eq:spin-projectors}
\end{align}
The outputs \texttt{nopol}, \texttt{pol15}, \texttt{pol25}, and
\texttt{pol35} correspond to these projectors. Equation~\eqref{eq:spin-projectors}
fixes the basis-dependent convention.

\subsection{Connected Wick terms}

Combining direct and exchange contractions gives
\begin{align}
  C_{3,u}^{\rm conn}
  &=-\mathcal U_1+\mathcal U_2+\mathcal U_3-\mathcal U_4,
  \\
  C_{3,d}^{\rm conn}
  &=+\mathcal D_1-\mathcal D_2.
  \label{eq:wick-signs}
\end{align}
The symbols include spectator perambulators, baryon tensors, diquark matrices,
the generalized vertex, and both flow-time-source sums. The six explicit
contractions and their axis order define the regression against the unflowed
code.

The saved three-point functions keep current-spin indices \(m,o\) open.
Contracting them later with \(\Gamma_{mo}\) forms unpolarized, helicity, or
transversity operators and reuses the costly flow contractions.

\subsection{Connected versus disconnected flavor contributions}

The \(u\)- and \(d\)-insertion arrays are connected only. In the isospin
limit, disconnected terms cancel in \(u-d\) if both flavors use identical
operators and normalization, but not in \(u+d\). A strange PDF requires
disconnected loops and is absent from these arrays.

\section{The present \texorpdfstring{\(32^3\times96\)}{32x96} implementation}

\subsection{Numerical definition}

Table~\ref{tab:implementation} lists the code defaults; they do not establish
a controlled continuum window.

\begin{longtable}{
  >{\raggedright\arraybackslash}p{0.29\textwidth}
  >{\raggedright\arraybackslash}p{0.65\textwidth}}
  \caption{Default parameters of the current light-flavor connected
  calculation.}\label{tab:implementation}\\
  \toprule
  Item & Present choice \\
  \midrule
  \endfirsthead
  \toprule
  Item & Present choice \\
  \midrule
  \endhead
  Gauge ensemble label
  & \texttt{beta6.41\_mu-0.2295\_ms-0.2050\_L32x96} \\
  Geometry
  & \(32^3\times96\), with periodic spatial directions \\
  Initial links supplied to flow
  & Original thin links (\texttt{Nosmear}); no HYP pre-smearing \\
  Flow evolution
  & Canonical thin-link 4D Wilson gauge flow and 4D covariant fermion flow
    through QUDA \texttt{performGFlowQuda} \\
  Flow parameters
  & \(\tf/a^2=0.5\), \(\epsilon=0.01\), \(N_{\rm step}=50\), and
    \(\rflow/a=\sqrt{8\tf/a^2}=2\) \\
  Temporal boundary condition
  & Anti-periodic fermion kernel; its temporal center phase is removed before
    saving the flowed links \\
  Distillation dimensions
  & \(N_v^{\rm current}=400\) for the flowed current and
    \(N_v^{N}=200\) for source and sink baryon tensors \\
  Flow batching
  & 32 modes per batch; four modes are packed into the spin components of one
    Wilson spinor and unpacked after flow.  Thus one source time uses 13
    multi-right-hand-side QUDA calls: twelve with eight spinors and one with
    four spinors \\
  Temporal support retained
  & \(|t_J-\tau|_{\rm periodic}\leq4\), hence nine source slices for each
    current time and \(9^2\) pairs \((\tau_a,\tau_b)\) \\
  Wilson-line geometry
  & Production direction \(z\), both line orientations \(-,+\), and stored
    lengths \(|z|/a=0,\ldots,20\) \\
  Momentum phase
  & \(\exp[-2\pi\ii(n_z z+n_y y+n_x x)/32]\), with one nonzero component and
    the production choices \(n_z=+3,+4,+5\) \\
  Source and sink times
  & All 96 source times; production source--sink separations
    \(T/a=6,7,8,9,10\) \\
  Flow convention for correlators
  & Unflowed \(C_2\); only the \(C_3\) insertion is flowed \\
  Numerical precision
  & Flowed basis and vertex default to \texttt{complex64}; contractions and
    final accumulations use \texttt{complex128} \\
  \bottomrule
\end{longtable}

At each \(t_J\), direction, and \((\tau_a,\tau_b)\), the stored vertex has
axes
\begin{equation}
  (s_z,\,|z|,\,n_{\rm left},\,n_{\rm right}),
  \qquad
  s_z=(\text{minus},\text{plus}).
\end{equation}
The contraction sums every retained \((\tau_a,\tau_b)\). Restricting to
\(\tau_a=\tau_b=t_J\) removes 4D temporal diffusion.

As in the reference code, \(C_2\) gains \(-1\) for
\(t_{\rm sink}<t_{\rm src}\), and \(C_3\) gains \(-1\) for
\(t_{\rm src}+T\geq N_t\). Zero-flow tests must verify these anti-periodic
wrap signs.

\subsection{Temporal-window approximation}

A 4D heat kernel has nonzero support on every time slice. Define the discarded
norm for source slice \(\tau\) as
\begin{equation}
  R_{\rm disc}(\tau;w)
  =1-
  \frac{\displaystyle
    \sum_{|t-\tau|_{\rm periodic}\leq w}
    \|\Phi_{\tau}^{(\tf)}(t)\|^2}
  {\displaystyle
    \sum_{t=0}^{N_t-1}\|\Phi_{\tau}^{(\tf)}(t)\|^2}.
  \label{eq:discarded-norm}
\end{equation}
For configuration 10000, \(\tf/a^2=0.5\), \(\epsilon=0.01\), and a
four-mode pilot, the observed discarded fractions were
\begin{center}
\begin{tabular}{ccc}
  \toprule
  Temporal window \(w\) & Retained slices & Discarded norm fraction \\
  \midrule
  2 & 5 & \(2.87\times10^{-4}\) \\
  3 & 7 & \(3.65\times10^{-6}\) \\
  4 & 9 & \(3.01\times10^{-8}\) \\
  \bottomrule
\end{tabular}
\end{center}
This pilot supports \(w=4\) for a \(10^{-6}\) norm threshold, but does not
bound correlator errors: contractions weight the discarded tails
nonuniformly. Selected correlators must compare \(w=3,4\) and a wider window.

\subsection{Dimensionless scale diagnostics}

For a momentum integer \(n\) along a length-32 direction,
\begin{equation}
  aP_z=\frac{2\pi n}{32}.
\end{equation}
The present flow time gives the following dimensionless combinations:
\begin{center}
\begin{tabular}{cccc}
  \toprule
  \(n\) & \(aP_z\) & \(\tf P_z^2\) & \(\rflow P_z\) \\
  \midrule
  3 & 0.5890 & 0.1735 & 1.1781 \\
  4 & 0.7854 & 0.3084 & 1.5708 \\
  5 & 0.9817 & 0.4819 & 1.9635 \\
  \bottomrule
\end{tabular}
\end{center}
Neither \(\tf P_z^2\) nor \(a/\rflow\) is uniformly small. The parameters are
a starting point, not a controlled matching window; flow-time and
lattice-spacing scans are required. Zero momentum is useful for checks, not
LaMET.

On a periodic length-32 lattice, \(+20a\) and \(-12a\) share an endpoint but
follow different paths. Fits should normally use nonwinding \(|z|<L/2\)
unless winding effects are modeled; stored points need not all enter the PDF.

\section{Scheme qualifications}

\subsection{Thin-link matching scheme}

The production calculation uses the canonical initial condition
\begin{equation}
  U_{\rm thin}
  \xrightarrow{\text{4D Wilson/fermion flow}}
  \bigl[B_\mu(\tf),K_{U_{\rm thin}}(\tf)\bigr],
  \qquad
  W=W_{B(\tf)}.
  \label{eq:thin-link-flow-scheme}
\end{equation}
Thus the flowed Wilson line and both flowed fermion legs use one common gauge
history. The code records \texttt{gauge\_input\_scheme=thin},
\texttt{gauge\_pre\_smearing=none}, and the input path in every v2 manifest;
the cache tag begins with \texttt{thin\_}. Legacy HYP+flow files are a distinct
finite-\(a\) discretization and are rejected rather than reused. Consequently,
the flowed gauge, basis, vertex, and \(C_3\) must be regenerated, while the
unflowed perambulators and baryon vertices remain reusable.

\subsection{Distillation and compression}

The current uses exact low modes, compressed intermediate modes, and an
orthogonalized stochastic high-mode completion.  It should therefore be
interpreted as a stochastic estimator of the conventional coordinate-space
propagator, not merely as a lowest-400-mode LapH truncation.  Orthonormality
makes the sampled-basis algebra valid, while \(\omega_{mn}\) supplies the
sector-size and repeated-index corrections.  Finite-noise, compression,
shared-noise, and \(N_v\) convergence tests remain necessary; sample-wise
orthogonality alone is not a proof of completeness or absence of stochastic
bias.

\section{What the current calculation does and does not determine}

The workflow computes:
\begin{itemize}
  \item flowed links and the 4D fermion kernel acting on all-time LapH vectors;
  \item flowed Wilson-line vertices with both kernel source times;
  \item unflowed proton \(C_2\) and flowed-insertion connected \(C_3\) for
        light \(u,d\);
  \item open current-spin components and the kinematics in
        Table~\ref{tab:implementation}.
\end{itemize}

These are bare correlators at finite flow time. They use the implemented
thin-link flow and distillation scheme; they are not physical PDFs. Missing
steps are
\begin{itemize}
  \item physical acceptance of the measured ringed normalization
        (Section~\ref{sec:ringed-tenconf-results});
  \item ground-state extraction from several source--sink separations;
  \item quark-line disconnected contractions;
  \item continuum, volume, flow-time, and large-momentum control;
  \item conversion to \(\MSbar\), LaMET matching, and power corrections;
  \item \(x\)-space reconstruction from finite, noisy separations.
\end{itemize}
The ratio \(C_3/C_2\) cancels nucleon overlaps, not flowed-field normalization
or scheme matching.

\section{Current quark ratio fit and small-flow checkpoints}
\label{sec:current-quark-ratio-fit}

This section records the September 2026 ten-configuration pilot so that the
figures and the flow production are not confused with the gluon \(N=406\)
analysis.  The input ensemble is \(32^3\times96\), with
\(a\simeq0.0775\,\mathrm{fm}\), thin-link four-dimensional flow, and the
connected \(u-d\) flavor combination.  The currently assembled pilot product
contains ten configurations,
\begin{equation}
  {\cal C}=\{10000,10100,\ldots,10900\},
\end{equation}
and four flow times
\(\tf/a^2=0.2,0.5,0.8,1.0\).  It is an early \(N_{\rm boot}=50\)
diagnostic; it must not be mixed with the later gluon \(N=406\),
\(N_{\rm boot}=1000\) or \(1500\) products.

\subsection{Ratio estimator and axes}

For every bootstrap replica \(b\), the code first averages the raw
configuration-level three- and two-point functions over the resampled
configuration list and then forms the ratio.  The unpolarized projection is
the even \(z\) combination and the stored estimator is
\begin{equation}
 R_U^{(b)}(T,\tau,z,P_z)
 =
 \frac{\operatorname{Re} C_{3,U}^{(b)}(T,\tau,z,P_z)}
      {\operatorname{Re} C_{2,\mathrm{nopol}}^{(b)}(P_z)}.
\end{equation}
The ratio is formed before the final real projection; this preserves the
nonlinear denominator and the shared bootstrap covariance.  The assembled
array has axes
\begin{equation}
(\tf,\mathrm{boot},P_z,z,T,\tau),
\end{equation}
with \(P_z=3,4,5\), \(T=6,\ldots,10\), and insertion times
\(\tau=0,\ldots,T\).  Invalid insertion points are retained as NaN masks,
not interpreted as zero measurements.

\subsection{Direct-ratio AIC fit}

The accompanying fit product uses a constant direct-ratio model in the
physical separation region.  With the strict threshold
\(T_{\min}=\lfloor0.6\,\mathrm{fm}/a\rfloor+1=8\), the fit separations are
\(T=8,9,10\).  Symmetric insertion cuts \(c=1,2,3\) remove the source and
sink contact regions.  The two registered windows \(8\)--\(9\) and
\(8\)--\(10\) are AIC weighted inside each cut, then the usable cuts are
mixed with equal weight.  This is a direct-ratio plateau diagnostic, not the
production correlated Sumratio fit:
\begin{equation}
 R_U(T,\tau)=M_U+\text{excited-state/contact corrections},
\qquad
 \mathrm{AIC}=\chi^2+2k+2N_{\rm omit},\quad k=1.
\end{equation}
The reported uncertainty is the bootstrap standard deviation with
\(\mathrm{ddof}=1\), without division by \(\sqrt{N_{\rm boot}}\).  The fit
does not apply \(Z_\chi\), SFTE, LaMET matching, disconnected diagrams, or a
continuum extrapolation.

\subsection{Ratio and fit figure}

The reproducible plotting entry point is
\begin{center}
\nolinkurl{/public/group/lqcd/donghx/Gradient_Flow_GluonPDF/plot_quark_ratio_fit_v1.py}.
\end{center}
It produces one page for each pair of flow time and momentum, with separate
subpanels for \(z/a=0,\ldots,15\).  Colored, unconnected points show the
different \(T\) ratios; a black horizontal line and gray band show the
direct-ratio AIC result and its \(1\sigma\) uncertainty.  The full output is
\begin{center}
\nolinkurl{/public/group/lqcd/donghx/Gradient_Flow_GluonPDF/quark_unpolarized_ratio_fit_aic_v1/ratio_C3_over_C2_with_direct_AIC_fit_Nboot50.pdf}.
\end{center}
The input hashes, axes, and fit rule are frozen in the accompanying
\nolinkurl{ratio_C3_over_C2_with_direct_AIC_fit_manifest.json}.  Figure
\ref{fig:quark-ratio-aic-current} shows its first page as an example.

\begin{figure}[htbp]
  \centering
  \includegraphics[width=0.98\textwidth,page=1]{figures/quark_ratio_C3C2_direct_AIC_Nboot50.pdf}
  \caption{Example page of the connected \(u-d\) flowed ratio and direct-ratio
  AIC fit.  The full twelve-page PDF contains all four flow times and
  \(P_z=3,4,5\).  This is a finite-flow bare diagnostic; it is not a
  renormalized quasi-PDF.}
  \label{fig:quark-ratio-aic-current}
\end{figure}

\subsection{New \(t_f/a^2=0.05\) and \(0.10\) checkpoints}

Ten independent GPU jobs were submitted with the same pilot configuration
list \({\cal C}\):
\begin{center}
\nolinkurl{/public/group/lqcd/donghx/Gradient_Flow_GluonPDF/submit_quark_flow_005_01_all10.sh}.
\end{center}
The driver uses \texttt{FLOW\_TIME\_SPECS=0.05:2;0.10:2},
\texttt{FLOW\_EPSILON=0.01}, \(P_z=3,4,5\) in the \(z\) direction, and
\(T=6,\ldots,10\).  The first checkpoint requires five flow steps.  The
\(0.10\) checkpoint restarts from \(0.05\) and advances five additional
steps, giving ten steps from the original thin links.  Each job requests four
A800 GPUs and four MPI ranks on \texttt{na800-sug} under reservation
\texttt{donghx\_223}.  The submission receipt is
\begin{center}
\nolinkurl{/public/group/lqcd/donghx/Gradient_Flow_GluonPDF/jobs/submitted_005_01_20260913_151237.md}.
\end{center}

At the 2026-09-14 status check, jobs 3629814--3629818 had exit code \(0:0\),
3629819--3629822 were running, and 3629823 was pending for resources.  The
checkpoint trees contained eight conf directories for \(0.05\) and six for
\(0.10\) at that time.  These counts are a dated status snapshot, not a
completion claim.  Final acceptance requires the production validator with
the same \texttt{SAVE\_ROOT}, momentum filters, tsep list, output hashes, and
raw \(C_2/C_3\) checks.  Until that validation is complete, the new checkpoints
must not be used for a small-flow extrapolation or called a physical PDF.

\clearpage
\section{Results inventory: lambda reconstruction, Fourier and LaMET
(9 October 2026)}
\label{sec:results-october}

This section updates the historical status snapshots without changing the
saved numerical products. The evidence is the existing source code,
READMEs, manifests, validation reports and plots inspected on 9 October.
No ratio, bootstrap, fit or matching calculation was repeated to produce
this documentation update. The processing order is
\begin{equation}
 \begin{aligned}
 M_{\rm flow}(t,z,P_z)
 &\longrightarrow M_{\rm converted}(t,z,P_z)
 \longrightarrow M_0(z,P_z),\\
 M_0(z,P_z)
 &\longrightarrow H_{\rm reconstructed}(\lambda,P_z)
 \longrightarrow \widetilde f(x,P_z)
 \longrightarrow f_{\rm LaMET}(x,\mu).
 \end{aligned}
 \label{eq:october-chain}
\end{equation}
The last arrow is a momentum-space convolution, hence it follows Fourier
transformation. \emph{Lambda extrapolation} here means a model for the
unmeasured tail of an Ioffe-time correlator, not a physical
\(\lambda\to\infty\) limit or a continuum extrapolation.

All relative analysis paths below are rooted at
\begin{center}
\nolinkurl{/public/group/lqcd/donghx/Gradient_Flow_GluonPDF}.
\end{center}
For L32x96, \(a=0.0775\,\mathrm{fm}\), \(L/a=32\), and
\begin{equation}
 P_z=\frac{2\pi n_z}{32a},\qquad
 \lambda=\frac{z_{\rm fm}P_{z,\rm GeV}}{\hbar c}
         =\frac{z}{a}\frac{2\pi n_z}{32},\qquad
 \hbar c=0.1973269804\,\mathrm{GeV\,fm}.
\end{equation}
The labels \(n_z=3,4,5\) correspond approximately to
\(P_z=1.50,2.00,2.50\,\mathrm{GeV}\). A file named \texttt{Pz3}
does not mean \(3\,\mathrm{GeV}\).
The original ensemble estimate is distinguished from its bootstrap
replicas. Statistical errors are replica standard deviations with
\(\mathrm{ddof}=1\), never divided by \(\sqrt{N_{\rm boot}}\).

\subsection{Flow-window restrictions and the origin of \(a^2/t_f\)}
\label{sec:flow-artifact-derivation}

The continuum small-flow expansion and the lattice cutoff expansion must
both be useful. For a nonlocal point, a candidate hierarchy is
\begin{equation}
 a\ll \rflow=\sqrt{8\tf}
 \ll\min\{|z|,\,|P_z|^{-1},\,\Lambda_{\rm QCD}^{-1}\},
 \qquad z\ne0.
 \label{eq:october-window}
\end{equation}
Numerical inequalities need observable-specific tests: the factor eight
in \(\rflow^2=8\tf\) cannot be hidden when comparing criteria.
To see why flow-enhanced cutoff terms can appear, linearize the transverse
flow and use the nearest-neighbor lattice Laplacian:
\begin{align}
 B_\mu(t,p)&=e^{-tp^2}A_\mu(p),\\
 \widehat p^2
 &=\frac4{a^2}\sum_\nu\sin^2(ap_\nu/2)
 =p^2-\frac{a^2}{12}\sum_\nu p_\nu^4+O(a^4p^6),\\
 e^{-t\widehat p^2}
 &=e^{-tp^2}
   \left[1+\frac{a^2t}{12}\sum_\nu p_\nu^4+\cdots\right].
 \label{eq:lattice-flow-kernel-expansion}
\end{align}
A flow-damped ultraviolet momentum integral samples \(p^2\sim1/t\);
its relative correction can therefore scale as
\(a^2tp^4\sim a^2/t\).
For a fixed infrared Fourier mode the correction is instead \(a^2tp^4\);
one must not assign \(a^2/t\) to every individual mode.

For example, extending the Brillouin zone in the regime \(t/a^2\gg1\),
the free four-dimensional heat trace gives
\begin{align}
 I(t,a)&=\int_{\rm BZ}\frac{\dd^4p}{(2\pi)^4}
                       e^{-2t\widehat p^2}\\
 &=\frac1{(8\pi t)^2}
   \left[1+\frac{a^2}{8t}+O\!\left(\frac{a^4}{t^2}\right)\right].
 \label{eq:heat-trace-cutoff}
\end{align}
Here the Gaussian average of each \(p_\nu^4\) is
\(3/(16t^2)\), so
\((a^2t/6)\sum_{\nu=1}^4\langle p_\nu^4\rangle=a^2/(8t)\).
The coefficient \(1/8\) belongs to this toy integral, not to the measured
gluon OPE. Improved actions, flow kernels and observables can cancel some
leading terms; additional \(a^2/z^2\) and \(a^2P_z^2\) corrections may remain.
Schematically,
\begin{equation}
 M_{\rm converted}(t,a)=M_R+tA+\frac{a^2}{t}B+\cdots .
\end{equation}
Thus a fixed-\(a\) intercept as \(t\to0\) is a diagnostic, not automatically
the continuum SFTE limit. Halving the integrator step \(\epsilon\) tests
flow integration, not these cutoff artifacts.

\begin{center}
\begin{tabular}{rrrr}
\toprule
\(\tau=\tf/a^2\) & \(\rflow/a\) & \(a^2/\tf\) &
\(\tf P_z^2\), \(n_z=5\)\\
\midrule
0.1 & 0.894 & 10.000 & 0.096\\
0.2 & 1.265 & 5.000 & 0.193\\
0.5 & 2.000 & 2.000 & 0.482\\
1.0 & 2.828 & 1.000 & 0.964\\
1.4 & 3.347 & 0.714 & 1.349\\
1.8 & 3.795 & 0.556 & 1.735\\
3.8 & 5.514 & 0.263 & 3.663\\
\bottomrule
\end{tabular}
\end{center}
In particular, the exploratory gluon cut \(\tau\ge1\) does not establish
a controlled scale hierarchy. AIC weights cannot repair an absent
hierarchy. Window scans should use declared physical criteria and preserve
the \(z,P_z\)-dependent masks. A controlled limit requires several lattice
spacings at fixed physical \(t,z,P_z\), or a justified joint fit with
identifiable cutoff terms.

\subsection{Gluon unpolarized: the \(406\)-configuration full chain}

The unified entry point is
\begin{center}
\nolinkurl{gluon_unpolarized_stages_CT18NLO_20260921/README.md}.
\end{center}
This branch uses \((T_U-S_U)_{\rm even}\), its real part, 406 common
configurations, seed 1115, and 1000 paired bootstrap replicas. The actual
ratio estimator averages \(C_2,C_3\) inside the ensemble or replica,
performs the vacuum subtraction and forms
\(\operatorname{Re}C_3/\operatorname{Re}C_{2,\rm nopol}\).
It is not an average of configuration-level ratios.
The bare fit uses a Sumratio forward-difference plateau with
\(T_{\min}=8\), hence \(aT_{\min}=0.62\,\mathrm{fm}>0.6\,\mathrm{fm}\).
The relaxed selection is \(Q\ge10^{-6}\) and
\(\chi^2/{\rm dof}\le10\), with flagged minimum-error fallbacks.
All 1050 finite-flow centers pass; 1598 of 1,050,000 replica positions
use a fallback. Complete numerical coverage is not a quality guarantee.

The eight zero-flow inputs are
\(\tau=1,1.4,1.8,2.2,2.6,3,3.4,3.8\).
The 21 contiguous windows containing at least three points use the full
paired cross-flow covariance, a linear intercept and replica AIC weights.
All 75 centers and 75,000 intercept replicas have acceptable windows
without new flow-fit fallback; upstream flags remain attached.

Relative to the \(z\)-directed line, \(F_{tx},F_{ty},F_{xy}\) are all
perpendicular--perpendicular endpoints. The archived nonlocal conversion is
\begin{align}
 M_{\MSbar}(t,z)&=\frac{e^{-\delta m|z|}}{c_\perp^2}M_{\rm flow}(t,z),\\
 c_\perp&=1+\frac{\alpha_s C_A}{4\pi}\ln(2\mu^2te^{\gamma_E}),&
 \delta m&=-\frac{\alpha_s C_A}{4\pi}\sqrt{\frac{2\pi}{t}},
 \label{eq:october-gluon-conversion}
\end{align}
with \(\mu=2\,\mathrm{GeV}\), \(\alpha_s=0.25\), \(C_A=3\),
and \(t\) in \(\mathrm{GeV}^{-2}\).
The stated \(\kappa_F=1\) normalization assumption still needs an
independent tree-level audit. These coefficients, their signs and the
endpoint sector must not be changed during plotting.

The current amplitude-preserving branch is
\begin{center}
\nolinkurl{matching/unpolarized_relaxed_q1em6_chi10_minerr_v3/}\\
\nolinkurl{lambda_fourier_quasi_no_z0_v1/quasi_support_refined/}.
\end{center}
It replaces the older \(M(z)/M(0)\) shape branch as the entry point for
absolute-amplitude comparisons. The Clover field represents \(a^2gF\).
The fundamental color-trace factor \(1/2\) cancels the even-sum factor two;
the ratio carries \(1/(2E)\). No extra volume or color factor is inserted:
\begin{equation}
 H(z)=-\frac{2R_{\MSbar}(z)}{4\pi\alpha_s(aE)},\qquad
 Q(x)=\frac1\pi\int_0^\infty\dd\lambda\,
                          \cos(x\lambda)H(\lambda).
 \label{eq:october-gluon-normalization}
\end{equation}
For this convention \(Q(x)=xg(x)\) at positive \(x\), and
\(H(0)=2\langle x\rangle_g\). The local point is matched independently:
\begin{align}
 C_{\rm loc}(t)&=1+\frac{\alpha_sC_A}{4\pi}
 \left[\frac{11}{3}\{\ln(2\mu^2t)+\gamma_E\}+\frac73\right],\\
 R_{{\rm loc},\MSbar}&=(2-C_{\rm loc})R_{{\rm loc},{\rm flow}} .
\end{align}
This is the strict one-loop inverse. The nonlocal endpoint formula must
not be substituted at \(z=0\). The unexpanded inverse is retained only
as a scheme-truncation sensitivity.

\begin{center}
\begin{tabular}{rrrr}
\toprule
\(n_z\) & \(aE\), statistical & \(\langle x\rangle_g\), statistical &
tail fallbacks/1000\\
\midrule
3 & \(0.74630\pm0.00382\) & \(0.26905\pm0.04288\) & 10\\
4 & \(0.91096\pm0.00995\) & \(0.32487\pm0.06298\) & 31\\
5 & \(1.10166\pm0.02216\) & \(0.21875\pm0.06807\) & 9\\
\bottomrule
\end{tabular}
\end{center}
The energy and local moment are paired with the same replicas.
The tail family is
\begin{equation}
 H_{\rm tail}(\lambda)=A(\lambda/8)^{-p}
                   e^{-(\lambda-8)/L},\qquad
 p\in\{0.5,1,2\},\quad L\in\{4,8,16\}.
\end{equation}
Three lower cuts \(\lambda_{\min}=4,5,6\), with
\(\lambda\le12\) and \(z/a\le16\), give 27 candidates.
Only the amplitude is fitted; replica quality selection and AIC are
preserved. The join is smooth over \(\lambda=7\)--8. The interval between
the local and first nonlocal point uses the declared quadratic completion.
Neither tail choice nor short-distance completion is a measured theorem.

The quasi-to-light-cone step uses the \(gg\) kernel associated with
arXiv:2510.26425v2 and a strict one-loop inverse. The scheme conversion
uses a perturbative hybrid proxy, not a measured \(P_z=0\) denominator.
The refined quasi support is \([-12,12]\), with \(\Delta x=0.01\) and
\(\xi_{\max}=200\). Points \(|x|<0.02\) remain masked, not filled with zero.
The archived support refinement changes the result by at most
about \(0.0254\) statistical standard deviations over \(0.2\le x\le0.8\);
this is a numerical test, not a bound on the model or physics uncertainty.

\begin{figure}[p]
 \centering
 \includegraphics[page=1,width=\textwidth]{figures/gluon_lambda_tail_N406_Nboot1000_20260921.pdf}
 \par\medskip
 \includegraphics[width=\textwidth]{figures/gluon_Fourier_quasi_N406_Nboot1000_20260921.pdf}
 \par\medskip
 \includegraphics[width=\textwidth]{figures/gluon_LaMET_N406_Nboot1000_20260921.pdf}
 \caption{Gluon unpolarized \(N_{\rm conf}=406,N_{\rm boot}=1000\):
 lambda-tail reconstruction (top), Fourier quasi distribution (middle),
 and quasi/LaMET comparison (bottom), with \(n_z=3,4,5\) in separate panels.
 These are fixed-spacing, \(gg\)-only diagnostics from the
 amplitude-preserving branch. Bands propagate paired statistical
 uncertainty; they do not include the scale-window, reconstruction or
 singlet-mixing systematics.}
 \label{fig:gluon-lambda-fourier-lamet}
\end{figure}

The CT18NLO reference uses \(Q=2\,\mathrm{GeV}\), PDG 21, and
\texttt{xfxQ}, which already returns \(xg\). Its 59 members are a central
set and 29 Hessian pairs, not bootstrap samples. The native 90\% confidence
band is converted to 68.268949\% by the factor \(1/1.64485362695\),
with asymmetric Hessian errors. Its coupling
\(\alpha_s(2\,\mathrm{GeV})\simeq0.29821\) differs from the archived lattice
choice 0.25 and has not been silently substituted.
For \(xg/\langle x\rangle_g\), each lattice replica is divided by its own
paired local moment; every CT18 member uses its own
\(\int_0^1xg(x)\dd x\), whose central value is approximately 0.41769.
No unit-area constraint is imposed on the measured finite \(x\) interval.

\begin{figure}[htbp]
 \centering
 \includegraphics[page=1,width=0.92\textwidth]{figures/gluon_normalized_CT18NLO_N406_20260921.pdf}
 \caption{Amplitude-preserving gluon LaMET diagnostic, subsequently plotted
 as \(xg/\langle x\rangle_g\), compared with CT18NLO at 2 GeV.
 The three momenta share a plot; the requested vertical range is
 \([-0.5,3]\). Normalizing this presentation does not validate the
 untested absolute normalization, hybrid proxy or small-flow hierarchy.}
 \label{fig:gluon-normalized-ct18-current}
\end{figure}
\clearpage

\subsection{Quark \(20\)-configuration, \(100\)-replica three-channel chain}

The upstream directories are
\begin{center}
\nolinkurl{quark_ratio_fit_20_ud_0p1_relaxed_20260927/}\\
\nolinkurl{quark_ratio_fit_20_ud_0p1_relaxed_matching_20260927/}\\
\nolinkurl{quark_ratio_fit_20_ud_0p1_relaxed_flow_extrap_20260927/}.
\end{center}
They cover \(U,H,T\), flavors \(u,d,u-d\), and \(n_z=3,4,5\).
The direct-ratio fit imposes \(aT>0.6\,\mathrm{fm}\), hence \(T_{\min}=8\),
with the relaxed \(Q>10^{-6}\), \(\chi^2/{\rm dof}<10\) selection and
flagged minimum-error replacements. The corrected four-dimensional SFTE
factor is \(Z_\chi e^{-\delta m|z|}/C_\psi\).
This branch assumes \(Z_\chi=1\), still unvalidated; it is not the
ten-configuration measured-ringed product.

The paired input flow times are \(0.1,0.2,0.5,0.8\).
The primary intercept averages three linear windows:
\([0.1,0.2,0.5]\), \([0.1,0.2,0.5,0.8]\), and \([0.2,0.5,0.8]\).
The declared score is
\(\chi^2+2k+2N_{\rm omitted}\), with \(k=2\).
The omitted-point penalty is a registered diagnostic choice, not a universal
correlated-data AIC theorem. A three-point \(\tau\ge0.2\) line and a
four-point quadratic are separate sensitivities.
Ratio-window systematics and flow-window spread are not absorbed into
the downstream statistical bands. Analytic odd zeros for \(H,T\) at \(z=0\)
are imposed during reconstruction; other failed data are never filled this way.

The post-processing root is
\begin{center}
\nolinkurl{flow_zero_lambda_fourier_quasi_matching_20260927/}.
\end{center}
Its \texttt{code/run\_pipeline.py} fits the retained component at
\(\lambda\ge7\), replica by replica, to a five-parameter projection:
\begin{align}
 f_R(\lambda)&=\left[\frac{A_1}{\lambda^a}
    +\frac{A_2\cos(\lambda-\pi b/2)}{\lambda^b}\right]
                          e^{-\lambda/\lambda_0},\\
 f_I(\lambda)&=\left[\frac{A_1}{\lambda^a}
    +\frac{A_2\sin(\pi b/2-\lambda)}{\lambda^b}\right]
                          e^{-\lambda/\lambda_0}.
\end{align}
The amplitudes are effective projected amplitudes, not the two amplitudes
of a complete complex correlator. Full bootstrap covariance is used with
5\% diagonal shrinkage and an SVD inverse.
PCHIP supplies the region below seven; the tail extends to 24 with
\(\Delta\lambda=0.05\). For the stored convention,
\begin{equation}
 \widetilde q_{\rm proj}(x)
 =\frac1\pi\int_0^{24}\dd\lambda\,
 [\operatorname{Re}H\cos(x\lambda)
       -\operatorname{Im}H\sin(x\lambda)] .
 \label{eq:quark-single-component-fourier}
\end{equation}
The grid is \(x\in[-2,2]\) with 401 points.
Only \(\operatorname{Re}H\) is supplied for \(U\), and only
\(\operatorname{Im}H\) for \(H,T\).
The orthogonal component is set to zero solely for this explicitly labeled
single-component diagnostic. This cannot determine a complete quark/antiquark
distribution or a final polarized PDF. The Wick-phase interpretation and
the transversity relative sign still require external physics confirmation.

The quasi kernel is the historical
\begin{center}
\nolinkurl{/public/group/lqcd/donghx/Diagram_GluonPDF/Compare_YaoFei/Matching/Matching_restore.py}.
\end{center}
The archived fixed-order hybrid parameters are
\(\mu=2\,\mathrm{GeV}\), \(\alpha_s=0.296\),
\(z_s=0.25\,\mathrm{fm}\), \(\xi_{\max}=50\), Gauss order 48,
and the parton-momentum logarithm prescription.
No local-charge division or charge restoration is made.
A complete theoretical audit closing the GF output scheme with this
historical hybrid kernel remains pending; numerical execution is not that audit.

\begin{center}
\begin{tabular}{lrrr}
\toprule
\(u-d\) channel & \(n_z=3\) & \(n_z=4\) & \(n_z=5\)\\
\midrule
\(U\): tail-boundary replicas/100 & 100 & 100 & 100\\
\(H\): tail-boundary replicas/100 & 99 & 100 & 77\\
\(T\): tail-boundary replicas/100 & 99 & 100 & 98\\
\bottomrule
\end{tabular}
\end{center}
The existing validation reports 27 Fourier and 27 matched HDF5 files,
all 100 replicas finite and numerically successful, with status
\texttt{PASS\_WITH\_TAIL\_MODEL\_WARNING}.
The boundary counts above must accompany the plots. An optimizer success
flag is not acceptance of a physical tail model.

\begin{figure}[!htbp]
 \centering
 \includegraphics[width=\textwidth]{figures/quark_lambda_unpolarized_ud_N20_Nboot100_20260927.pdf}
 \par\bigskip
 \includegraphics[width=\textwidth]{figures/quark_LaMET_unpolarized_ud_N20_Nboot100_20260927.pdf}
 \caption{Quark unpolarized connected \(u-d\), 20 configurations and
 100 paired replicas: real-component lambda reconstruction (top) and
 Fourier/quasi matching (bottom). \(Z_\chi=1\) is assumed.
 All 100 tail fits for each momentum reach a parameter boundary.
 The missing imaginary component and the kernel-scheme audit prevent
 interpretation as a complete \(u-d\) PDF.}
 \label{fig:quark-unpolarized-chain-current}
\end{figure}
\clearpage
\begin{figure}[p]
 \centering
 \includegraphics[width=\textwidth]{figures/quark_lambda_helicity_ud_N20_Nboot100_20260927.pdf}
 \par\bigskip
 \includegraphics[width=\textwidth]{figures/quark_LaMET_helicity_ud_N20_Nboot100_20260927.pdf}
 \caption{Quark helicity connected \(u-d\) single-imaginary-component
 diagnostic: lambda continuation (top) and Fourier/quasi matching (bottom).
 Tail boundaries occur in 99, 100 and 77 replicas out of 100 for
 \(n_z=3,4,5\), respectively. The physical complex projection and
 normalization remain to be verified.}
 \label{fig:quark-helicity-chain-current}
\end{figure}
\clearpage
\begin{figure}[p]
 \centering
 \includegraphics[width=\textwidth]{figures/quark_lambda_transversity_ud_N20_Nboot100_20260927.pdf}
 \par\bigskip
 \includegraphics[width=\textwidth]{figures/quark_LaMET_transversity_ud_N20_Nboot100_20260927.pdf}
 \caption{Quark transversity connected \(u-d\) single-imaginary-component
 diagnostic: lambda continuation (top) and Fourier/quasi matching (bottom).
 Tail boundaries occur in 99, 100 and 98 replicas out of 100.
 The operator relative sign is pending external physics confirmation.
 Gray endpoint regions on matching plots do not imply that the remaining
 region has passed the missing theory and systematics checks.}
 \label{fig:quark-transversity-chain-current}
\end{figure}
\clearpage

The unpolarized reference is CT18NLO. The helicity comparison actually used
\texttt{NNPDFpol20\_lo\_as\_01180}, the installed LO fallback, not an NLO set.
The dedicated transversity comparison uses
\begin{center}
\nolinkurl{/public/group/lqcd/donghx/Diagram_GluonPDF/Matching/JAM2022.txt}.
\end{center}
Its three columns are \(x\), central value and one-sigma error. The recorded
scale is 2 GeV, and the plotted quantity is \(q\), not \(xq\); neither
curve is divided by the tensor charge. The provenance/convention connection
and the lattice operator sign are still pending confirmation.
Since the lattice input lacks the orthogonal complex component, comparisons
with complete phenomenological flavor PDFs are diagnostic overlays.

\begin{figure}[htbp]
 \centering
 \includegraphics[width=\textwidth]{figures/quark_transversity_JAM2022_N20_20260928.pdf}
 \caption{Recorded \(u-d\) transversity/JAM2022 diagnostic at 2 GeV,
 without tensor-charge normalization. This preserves the existing plot,
 not a validation of its single-component reconstruction or relative sign.
 Statistical bands omit the upstream flow-window and tail-model systematics.}
 \label{fig:quark-jam-diagnostic-current}
\end{figure}

\subsection{Keep the measured-\(Z_\chi\) ten-configuration branch separate}

The earlier 50-replica branches remain available under
\begin{center}
\nolinkurl{quark_three_stages_10conf_Nboot50_20260920/}\\
\nolinkurl{quark_Zchi1_control_10conf_Nboot50_20260920/}.
\end{center}
Each has
\nolinkurl{lambda_fourier_lamet_valence_v2/{sixflow_sumdiff,fourflow_ratio}/}.
They use \(\tau=0.2,0.5,0.8,1.0\) linear flow intercepts, with
quadratic sensitivity. Unlike the twenty-configuration single-component
prototype, these branches explicitly interpret the unpolarized real part
as a valence projection. With \(q(-x)=-\bar q(x)\),
\begin{equation}
 v_{u-d}=(u-\bar u)-(d-\bar d),\qquad
 \widetilde v(x)=\frac2\pi\int_0^\infty\dd\lambda\,
                         \cos(x\lambda)\operatorname{Re}h(\lambda).
\end{equation}
The factor two cannot be dropped, and this output is not the full
positive-\(x\) \(u-d\) distribution.

This branch uses the direct \(\MSbar\), \(\gamma_0\) kernel of
arXiv:1801.03917v2, Eq.~(68), with \(\mu=2\,\mathrm{GeV}\),
\(\alpha_s=0.30\) and a strict one-loop inverse.
The nonlocal \(\MSbar\) short-distance OPE logarithm is modeled explicitly;
a finite local charge must not simply be inserted as nonlocal \(h(0)\).
Version two includes the large-\(|y|\), \(1/|y|\) Fourier tail induced
by that logarithm; the hard-support version one is historical.
There is no division by \(h(0)\) or enforced valence integral.

The four branches are numerically complete. The archived checks include
paired replica identities, source preservation, a Fourier analytic test
with maximum discrepancy \(3.33\times10^{-16}\), and an independent
high-\(|y|\) integral check at \(1.1\times10^{-13}\).
However, all registered tail models for \(n_z=3,4\) have
nominal \(Q<0.05\), in every branch.
Over \(0.2\le x\le0.8\), short-distance variations reach about
18.46 main statistical standard deviations and linear-to-quadratic
flow variations reach about 17.39. These are model sensitivities, not
new confidence intervals. The narrow statistical bands cannot be
reported as total PDF precision.
The corresponding CT18 comparisons use \(u_v-d_v\).
An additional plot labeled with \(\bar u=\bar d=0\) is a separate
assumption, not a measured sea separation.

\subsection{Data formats, paths and numerical evidence}

For the gluon unified directory, the principal arrays are
\begin{center}
\begin{tabular}{p{0.48\textwidth}p{0.43\textwidth}}
\toprule
File under \texttt{data/} & Axes or shape\\
\midrule
\nolinkurl{finite_flow_stages.npz} &
\((14,1001,3,25)\): flow, original center then bootstrap, momentum, \(z\)\\
\nolinkurl{flow_zero_selected.npz} &
center \((3,25,2)\), boot \((1000,3,25,2)\); intercept/slope parameter axis\\
\nolinkurl{lightcone_and_quasi.npz} &
matched and quasi: \((3,1001,201)\), plus computed masks\\
\nolinkurl{lambda_reconstruction.npz} &
\((3,1001,801)\), tree and hybrid-proxy amplitudes\\
\bottomrule
\end{tabular}
\end{center}
Evidence is stored in
\nolinkurl{validation/chain_checks.json},
\nolinkurl{validation/collection_readback.json} and
\nolinkurl{provenance/source_LaMET_validation.json}.
These archived reports pass the declared numerical checks;
\texttt{physical\_acceptance=false} remains in the manifest.
The source numerical entry points are
\nolinkurl{matching/gluon_absolute_normalization_v1.py},
\nolinkurl{matching/lambda_fourier_quasi_no_z0_v1.py} and
\nolinkurl{matching/refine_no_z0_quasi_support_v1.py}.

For the twenty-configuration quark root, \texttt{results/} contains
\begin{itemize}
 \item \nolinkurl{fourier/<channel>/<flavor>/fourier_L32x96_*.h5},
 with \texttt{quasi\_pdf/samples} of shape \((401,100)\);
 \item \nolinkurl{matching/<channel>/<flavor>/matched_L32x96_*.h5},
 with \texttt{matched\_pdf/samples}, preserving the original replica order;
 \item \nolinkurl{plots/lambda/lambda_extrapolation_<channel>_<flavor>.pdf}
 and \nolinkurl{plots/matching/matching_<channel>_<flavor>.pdf};
 \item \nolinkurl{plots/lhapdf/matched_vs_lhapdf_<channel>_<flavor>.pdf}
 and the dedicated \nolinkurl{plots/jam2022/} transversity comparison;
 \item \nolinkurl{manifest.json}, \nolinkurl{summary.csv} and
 \nolinkurl{validation_report.json}.
\end{itemize}
Here \texttt{channel} is unpolarized/helicity/transversity and
\texttt{flavor} is u/d/u\_minus\_d. Thus the \(u,d\) outputs not embedded
in this note are available with the same layout.
Submit \texttt{bash slurm/submit\_pipeline\_cpu.sh} from the pipeline
directory for CPU processing, after selecting a new versioned output
for any rerun. Analysis and numerical reconstruction are not to be
launched as heavy login-node jobs.
The ten-configuration direct-\(\MSbar\) code and detailed numerical
checks are in \nolinkurl{quark_lambda_fourier_lamet_20260920/};
its \texttt{valence\_results.npz} arrays have shape \((3,51,99)\).

\subsection{L32x64 production snapshot and next checks}

The L32x64 ensemble has \(a=0.0897\,\mathrm{fm}\) and uses thin LIME
inputs from the beta6.308 configuration directory.
The CPU gluon campaign
\nolinkurl{gluon_l32x64_allconf_20261009/} enumerated 1700 valid inputs,
configurations 1000 through 17990, and 23,800 requests for the same
14 gluon flow times as L32x96. It uses \(\epsilon=0.01\),
fine-step test \(0.005\), \(z\)-direction 2 and 25 separations.
Initial audit job 3825824 completed with exit zero:
23,038 existing raw products were validated/exported and 762 missing
or invalid requests were identified in 600 configuration groups.
All 14 step-halving gates passed, with maximum relative OPE \(L^2\)
change \(2.18071\times10^{-5}\) and plaquette difference
\(7.51153\times10^{-6}\).

Automatic dispatch 3825825 failed because the account permits at most
200 submitted jobs. The missing-only recovery array 3826064 therefore
has 60 items, each processing ten groups sequentially, without a script
concurrency cap. Final audit 3826065 depends on its successful completion.
At this dated snapshot production was still active; all-request completion
requires zero missing entries in the final audit, not merely existing
directories. The four CPU partitions are
\texttt{a192c1t}, \texttt{cpu6248R}, \texttt{cpueicc}, \texttt{i72c512g}.
The primitive terms are saved separately as \(T_U,S_U,T_H,S_H\),
with target \((T_U-S_U)_{\rm even}\) and \((T_H-S_H)_{\rm odd}\).
The plus helicity combination and \(T_H\) remain diagnostics.
For example, the separated primitives have shape \((2,4,25,64)\),
with projection/orientation/separation/time axes; full complex values
and the declared traced-field variants are retained.

L32x64 quark production in
\nolinkurl{quark_l32x64_flow_20261007/} is a separate
20-configuration, eight-flow campaign. Array 3816972 was still running
or pending on \texttt{na800-sug}, with its existing concurrency limit four.
Future GPU submissions use reservation \texttt{donghx\_223}.
Basis caches are retained; recoverable vertex/VVV caches need not persist.
Neither production counts nor the eventual 30-configuration campaign
change the twenty-configuration analysis scope retroactively.

The next physics checks are the quark tail residual/boundary scan,
completion of the missing complex projections and sign audit,
paired \(Z_\chi\) measurements on the enlarged common configuration set,
scheme closure between GF conversion and the chosen quasi kernel,
and separate short-distance/tail/flow-window systematic variations.
Gluons additionally need independent normalization and local-point checks,
the hybrid denominator and singlet mixing; helicity must not reuse
an unpolarized kernel by assumption.
Only additional spacings and a demonstrated hierarchy can turn the
present fixed-spacing diagnostics into controlled continuum PDFs.
\clearpage

\section{Blending weights, all-time construction and temporal-window tests}
\label{sec:blending-oct15}

This section records the audit completed on 15 October 2026. It separates
an estimator definition, an implementation regression and a physical
validation. The historical production estimator is called A. The user has
subsequently selected B for future C3 calculations; its implementation and
window test are isolated successors. No all-current, all-separation,
all-flow B production rerun has been submitted as part of this documentation
update. Existing fits, flow conversions and reconstructed PDFs remain
historical A products unless explicitly regenerated and labelled otherwise.

\subsection{Time labels, sampling identity and weights}

We use \(\tf\) for physical flow time,
\(\tau=\tf/a^2\) for dimensionless flow time,
\(\epsilon=\Delta\tau\), \(\sigma\) for an original basis time, and
\(t_J\) for the output/current time. Thus \path{basis/tau000/} means
\(\sigma=0\), not zero flow. A sampled vector is identified by the full
label \((\sigma,r,m)\): original time, spectral sector and saved slot.
Equal slot numbers at different original times are not the same
four-dimensional input vector.

For a sector of dimension \(d_r\) with \(n_r\) retained vectors, the
one-index weight is \(w_{\sigma m}=d_r/n_r\). Under the conditional model
of a uniform subset without replacement from the same orthonormal frame,
the second-order table is
\begin{equation}
 \omega^{(2)}_{mn}=
 \begin{cases}
 d_r/n_r,&m=n\text{ in sector }r,\\
 d_r(d_r-1)/[n_r(n_r-1)],&m\ne n\text{ in the same sector }r,\\
 (d_r/n_r)(d_s/n_s),&m\in r,\ n\in s\text{ independent sectors}.
 \end{cases}
 \label{eq:omega-second-oct15}
\end{equation}
The first and second cases encode different inclusion probabilities.
They must not be replaced by the square root of the diagonal table to
manufacture a one-index factor.

The two source-time prescriptions are
\begin{align}
 \Omega^A_{(\sigma_a,m),(\sigma_b,n)}&=\omega^{(2)}_{mn},\\
 \Omega^B_{(\sigma_a,m),(\sigma_b,n)}&=
 \begin{cases}
 \omega^{(2)}_{mn},&\sigma_a=\sigma_b\pmod {N_t},\\
 w_{\sigma_a m}w_{\sigma_b n},&\sigma_a\ne\sigma_b\pmod {N_t}.
 \end{cases}
 \label{eq:scheme-B-oct15}
\end{align}
Flow is a deterministic map on the already sampled columns at fixed gauge
field. It does not resample them; neither loss of norm nor a nontrivial
flowed Gram matrix changes the original inclusion probabilities.
Both prescriptions are independent of \(\tf\), \(t_J\), and separation
at fixed sampling rule. B depends on the equality of the \emph{original}
times, not on the output time.

For L32x64 the exact sector has \((d,n)=(100,100)\); ten compressed
subsectors have \(n=10\) each, with dimensions
\begin{equation}
 (d_1,\ldots,d_{10})=(50,50,100,100,200,200,350,350,550,550).
\end{equation}
The high sector has \((d,n)=(95704,200)\), so its first weight is
478.52. The full space is \(32^3N_c=98304\), but only 400 columns are
retained. Source/sink VVV uses the first 200 columns, including compressed
random columns; these are not 200 exact modes. L32x96 uses different sector
dimensions, so its weight table cannot be copied numerically.

\subsection{What is proved conditionally, and what is not proved}

For two independent source-time frames, suppose every random slot in one
sector satisfies \(\mathbb E R_{\sigma m}=P_\sigma/d\), and let the
coefficient \(M\) be independent of these draws. Then
\begin{equation}
 \mathbb E\sum_{m,n}C_{mn}R_{\sigma_a m}M R_{\sigma_b n}
 =\frac{\sum_{m,n}C_{mn}}{d^2}P_{\sigma_a}MP_{\sigma_b}.
\end{equation}
Both A and B have the required table sum:
\begin{equation}
 \sum_{m,n}\omega^{(2)}_{mn}
 =n\frac dn+n(n-1)\frac{d(d-1)}{n(n-1)}=d^2
 =\sum_{m,n}w_mw_n.
\end{equation}
Thus a difference between individual cross-time table entries does not by
itself prove a bias in A: under these assumptions the complete mode sums
can have the same mean and different finite-sample values or variances.
For example, \(d=50,n=10\) gives A diagonal 5, off-diagonal
\(27.222\ldots\), and B entries 25; both tables sum to 2500.
Exact sectors require a deterministic complete sum instead of a random
slot assumption.

The complete baryon coefficient also contains source/sink VVV, spectators
and other lines which may share noise, especially when original current
times coincide with source or sink. The actual frozen generator, seeds,
slot moments and higher joint-noise moments are therefore \textbf{still
to be confirmed}. Choosing B is not a proof of its full-baryon
unbiasedness. Zero-flow and factorization regressions cannot establish
this missing stochastic statement. A measured A/B difference is neither
a bias estimate nor a final PDF systematic uncertainty.

The completed A/B audit is in
\nolinkurl{quark_l32x64_omega_compare_20261014/} below the project root
\path{/public/group/lqcd/donghx/Gradient_Flow_GluonPDF/}.
It contains ten configurations and four current times per configuration.
Equal-time blocks and their Wick terms agree exactly; A reproduces the
old production with relative \(L^2\) error about \(9.1\times10^{-15}\).
Total C3 relative A/B differences range from 1.3827--2.4865\% for \(u\)
and 1.7738--3.3757\% for \(d\) in the tested cases. All 40 cases pass
implementation/completeness checks. The corrected CPU summary job
3851153 completed after the earlier environment-startup failure; the
failure was not a GPU numerical failure.

\subsection{From the four-dimensional flow to basis and vertex}

The continuum equations and their kernel interpretation are
\begin{align}
 \partial_{\tf} B_\mu&=D_\nu G_{\nu\mu},&
 \partial_{\tf}\chi&=\Delta_{B(\tf)}\chi,\\
 K(\tf)&=\mathcal P_f\exp\!\left(\int_0^{\tf}\dd s\,
       \Delta_{B(s)}\right),&
 \Phi&=K\mathcal V.
\end{align}
The input embedding and its matrix shape are
\begin{equation}
 \eta_{\sigma n}(x)=\delta_{x_4,\sigma}V_n(\bm x;\sigma),\qquad
 \mathcal V\in\mathbb C^{(N_tN_s^3N_c)\times(N_tN_v)}.
\end{equation}
Spin is an identity factor suppressed here. Embedding adds zero entries,
not independent columns; hence there are \(N_tN_v\), not
\(N_tN_s^3N_c\), input columns. The implementation evaluates \(K\mathcal V\)
and does not build dense \(K\), solve the full spectrum, or project the
evolution back to 400 dimensions at every step.

The dimensionless lattice Laplacian is
\begin{equation}
 (L_U\phi)(x)=\sum_{\mu=1}^4\bigl[
 U_\mu(x)\phi(x+\hat\mu)+U_\mu^\dagger(x-\hat\mu)
 \phi(x-\hat\mu)\bigr]-8\phi(x).
\end{equation}
The coupled RK3 stages are given in
Section~\ref{sec:pyquda-gflow}, with \(\epsilon=\Delta(\tf/a^2)\),
\(N_{\rm step}=\tau/\epsilon\) and each Laplacian using its matching
intermediate gauge. There is no mass term or Dirac solve. The heat map is
not unitary: even if \(\mathcal V^\dagger\mathcal V=I\), its flowed
Gram matrix is \(\mathcal V^\dagger K^\dagger K\mathcal V\), not generally
\(I\). Neither a Gram inverse nor keeping only Q from a QR decomposition
is part of this algorithm.

The implementation order is as follows.
\begin{enumerate}
 \item Validate the original basis/perambulator identity, sector ordering
 and unsmeared LIME gauge. The historical basis filename's HYP label does
 not change the thin-link starting field of the flowed operator.
 \item Embed each \(\sigma\). Pack four separate input columns into four
 spin lanes, possible because the Laplacian is spin diagonal. A batch of
 32 columns uses eight Wilson spinors; 400 columns require 13 batches,
 the last containing 16 columns. Separate spin lanes mean no spin mixing,
 not statistical independence of the input columns. This packing is not resampling.
 \item Evolve on the full four-dimensional lattice, with fermion temporal
 antiperiodicity and a common gauge trajectory for both endpoints and
 their Wilson line. Remove the auxiliary fermion center phase from the
 gauge saved for the operator. Multi-flow checkpoints continue the full
 resident field, never a truncated saved slice.
 \item Unpack all output times, accumulate full norm diagnostics, then
 retain the selected slices. A basis slice has shape
 \((400,32,32,32,3)\), complex64; gauge links are complex128. Coordinate
 storage is \((t,z,y,x)\), link directions are \((x,y,z,t)\).
 \item At fixed \(t_J\), enumerate all retained ordered source-time pairs
 and construct both signed straight Wilson chains. Compute
 \begin{equation}
 J(t_J;\sigma_a,\sigma_b;z)
 =\mathcal V_{\sigma_a}^\dagger K^\dagger
    \mathcal O(t_J,z;B)K\mathcal V_{\sigma_b},\qquad
 \widetilde J=\Omega^B\odot J.
 \end{equation}
 Geometry has shape \((2,21,400,400)\): sign, separation 0--20, left
 mode, right mode. Apply the Hadamard table exactly once after geometry;
 do not weight basis or add a second table at C3 read time. The current
 spin matrix \(\Gamma\) is applied later.
 \item Combine \(\widetilde J\) with ordinary perambulators and unflowed
 source/sink VVV, using complex128 accumulation and the unchanged six
 Wick terms. The factorization
 \texttt{tuSLnp,tambn->uSLpamb} sums one original time first and retains
 the other; it is a finite-sum reordering, not another sampling step.
 The external antiperiodic sign is \(-1\) if
 \(t_{\rm src}+t_{\rm sep}\ge N_t\).
\end{enumerate}
The two time sums arise from the two endpoints. Reduce time labels modulo
\(N_t\), deduplicate, and enumerate ordered pairs. Equal-time pairs appear
once, reversed unequal pairs are distinct. There is no division by
\(N_t\) or the number of window pairs and no extra factor \(1/2\).

\subsection{Full integration versus finite temporal storage}

Define periodic distance
\(d_{N_t}(t_J,\sigma)=\min_{k\in\mathbb Z}|t_J-\sigma+kN_t|\).
Full flow is solved first. Each input-time block has its own output
projector, \(\Phi_{\sigma,w}=M_{\sigma,w}\Phi_\sigma\), keeping
\(d_{N_t}(t_J,\sigma)\le w\). Concatenate these blocks to define
\(\Phi_w\); this is not a single common output-time projector applied
to all input columns. For \(E=\Phi-\Phi_w\), the aggregate diagnostic is
\begin{equation}
 f_{\rm discard}=\frac{\|E\|_F^2}{\|\Phi\|_F^2},\qquad
 \frac{\|E\|_F}{\|\Phi\|_F}=\sqrt{f_{\rm discard}}.
 \label{eq:window-tail-oct15}
\end{equation}
Aggregation is over modes, space, color and all output times for
\emph{each configuration and each original time}, not an ensemble or
bootstrap average. A \(10^{-6}\) squared-norm gate is a \(10^{-3}\)
amplitude-norm scale, not a \(10^{-6}\) C3 error bound.
For the assembled geometric operator,
\begin{align}
 J-J_w&=\Phi_w^\dagger\mathcal O E+E^\dagger\mathcal O\Phi_w
        +E^\dagger\mathcal O E,\\
 \|J-J_w\|_F&\le\|\mathcal O\|_2
       \bigl(2\|\Phi_w\|_F\|E\|_F+\|E\|_F^2\bigr).
 \label{eq:window-bound-oct15}
\end{align}
This absolute bound does not give a relative bound on a selected output
time, mode, weighted vertex or contracted C3: weights, perambulators and
cancellations still matter. A local Gaussian free heat kernel motivates
checking a finite window, but does not replace a paired C3 test.
Discarded norms cannot recover missing complex amplitudes. One must not
renormalize weights to compensate for a missing tail.

\subsection{Completed L32x64 scheme-B window experiment}

The isolated experiment is
\nolinkurl{quark_l32x64_schemeB_window_20261015/} below the project root.
Its frozen \texttt{manifest.json} uses configurations
10000, 10040, \ldots, 10360, \(a=0.0897\,\mathrm{fm}\),
\(\tau=0.5\), \(\epsilon=0.01\), 50 steps,
\(t_J=0,16,32,63\), \(t_{\rm sep}=7a=0.6279\,\mathrm{fm}\),
and insertion offset \(dt=3\). Integer momenta are 3, 4, 5 in the
z direction, separations 0--20 with both signs, four nucleon
polarizations, full open \(4\times4\) current spin and all six Wick terms.

The test integrates all 400 columns at the 40 required original times
per configuration over the full 64-time lattice. It stores only the
52 required \((\sigma,t_J)\) slices, about 15.24 GiB per configuration,
while norm diagnostics still cover all output times. All windows 4, 5, 6
use this same complete response. Its selected-output cache schema is not
a standard complete production-basis schema.

All ten configurations and 40 cases pass, with no missing products.
The 360 overlapping old/new \(w=4\) basis slices and C2 agree exactly.
The largest relative \(L^2\) regression errors are
\(4.744\times10^{-15}\) for B total C3 and
\(9.327\times10^{-15}\) for ordered-block Wick terms.
The physical-data factorized-versus-explicit contraction check at
configuration 10000, \(t_J=0\), covers all three windows and has maximum
\(5.449\times10^{-15}\).

\begin{table}[htbp]
 \centering
 \caption{Maximum discarded squared-norm fractions over tested
 configurations and original times. The one-body weighted norm is an
 auxiliary first-weight diagnostic, not the full B-weighted C3.}
 \label{tab:window-norm-oct15}
 \begin{tabular}{rcccc}
 \toprule
 \(w\) & Aggregate & Per-mode max & Weighted one-body & Sector max\\
 \midrule
 4 & \(2.78755\,10^{-8}\) & \(5.86455\,10^{-8}\)
   & \(4.65347\,10^{-8}\) & \(5.77024\,10^{-8}\)\\
 5 & \(1.59311\,10^{-10}\) & \(3.89139\,10^{-10}\)
   & \(2.98558\,10^{-10}\) & \(3.82405\,10^{-10}\)\\
 6 & \(6.75419\,10^{-13}\) & \(1.90064\,10^{-12}\)
   & \(1.41614\,10^{-12}\) & \(1.86526\,10^{-12}\)\\
 \bottomrule
 \end{tabular}
\end{table}
\begin{table}[htbp]
 \centering
 \caption{Largest relative Frobenius change of total complex C3 among
 the 40 cases, at fixed scheme B. Entries are percentages, not fractions.}
 \label{tab:window-c3-oct15}
 \begin{tabular}{rccc}
 \toprule
 Window pair & \(u\) & \(d\) & \(u-d\)\\
 \midrule
 \(4\to5\)&0.0126696\%&0.0146414\%&0.0134355\%\\
 \(5\to6\)&0.00080954\%&0.00093124\%&0.00082711\%\\
 \(4\to6\)&0.0130625\%&0.0149632\%&0.0137709\%\\
 \bottomrule
 \end{tabular}
\end{table}

For the open-spin ratio diagnostic the order is: form complex
\(C_3/\operatorname{Re}C_{2,\mathrm{nopol}}\) in each configuration/case,
average four cases \emph{inside} each configuration, then compute real
and imaginary means and SEM across ten configurations. Four current times
are not 40 independent configurations. No new bootstrap, physical
\(\Gamma\) projection or bare-element fit is performed in this test.
For \(w=4\to6\), the norm of the mean ratio shift divided by the norm
of the baseline SEM ranges from \(1.60718\times10^{-4}\) to
\(4.71277\times10^{-4}\) across flavors and momenta. This is descriptive,
not a pointwise significance or covariance test.

At the tested statistical precision these results support \(w=4\) for
B pilot production in this restricted range. \(w=6\) is still a finite
window, not an exact full-time C3 reference. This is not an approval of
all flow times, source/sink separations, physical channels or L32x96.

\begin{figure}[htbp]
 \centering
 \includegraphics[width=\linewidth]{figures/basis_tail_diagnostics.pdf}
 \caption{Full-integration norm diagnostics for windows 4, 5 and 6.
 The gate applies to the aggregate squared norm only. Input is the
 existing experiment's \texttt{summary/basis\_tail\_metrics.csv}.}
 \label{fig:window-norm-oct15}
\end{figure}
\begin{figure}[htbp]
 \centering
 \includegraphics[width=\linewidth]{figures/C3_window_convergence.pdf}
 \caption{Paired complex C3 convergence at fixed B for \(u,d,u-d\).
 Every pair uses identical basis, gauge, perambulators and noise;
 it does not mix A/B variation with temporal truncation.}
 \label{fig:window-c3-oct15}
\end{figure}
\begin{figure}[htbp]
 \centering
 \includegraphics[width=\linewidth]{figures/paired_open_ratio_window_shift.pdf}
 \caption{Norm of paired mean open-spin ratio change relative to a
 baseline SEM norm. This global diagnostic is neither a projected
 physical PDF ratio nor a pointwise significance measure.}
 \label{fig:window-ratio-oct15}
\end{figure}
\clearpage

\subsection{Code, data contract and concrete next steps}

All following paths are relative to the B experiment root above.
\begin{longtable}{>{\raggedright\arraybackslash}p{0.43\textwidth}
                  >{\raggedright\arraybackslash}p{0.48\textwidth}}
 \toprule
 Entry or product & Role / contract\\
 \midrule
 \endhead
 \path{reference/flow_common_l32x64.py} &
 \texttt{first\_order\_weights}, \texttt{create\_time\_pair\_omega},
 source-time reduction and geometric vertex.\\
 \path{reference/generate_flowed_vertex_l32x64.py} &
 Standard B generator; \texttt{vertex\_schemeB} cache, schema
 \nolinkurl{generalized_flowed_vertex_l32x64_schemeB_v1}.\\
 \path{generate_selected_basis.py}, \path{compare_windows.py} &
 Full-grid integration/selected output for this test; paired
 window reconstruction and Wick/factorization regression.\\
 \path{summary/report.json}, \path{results/<conf>/tJ###/receipt.json} &
 PASS, limits, regressions and missing-product audit.\\
 \path{summary/paired_window_samples.npz} &
 Paired C3/C2, ratios, real/imaginary SEM and per-conf differences.\\
 \path{summary/C3_window_metrics.csv},
 \path{summary/basis_tail_metrics.csv} &
 All case-by-case C3 changes and norm diagnostics.\\
 \path{run_B_from_cache.py},
 \path{run_B_from_cache_na800_sug.slurm} &
 Prepared full-current B entry, not yet submitted in bulk; requires
 explicit \texttt{B\_CONF} and standard complete basis/gauge.\\
 \bottomrule
\end{longtable}
The single-case arrays are in
\path{results/<conf>/tJ###/window_comparison.npz}.
\texttt{c3\_by\_window\_flavor} has shape
\((3,2,2,4,3,21,4,4)\): window, flavor \((u,d)\), sign, nucleon
polarization, momentum, separation and open spins \((m,o)\).
\texttt{signed\_wick\_by\_window} has shape
\((3,6,2,4,3,21,4,4)\), with signs already included;
\texttt{signed\_wick\_by\_radius} has leading radius length 7;
\texttt{c2\_common} has shape \((4,3)\).
\(u-d\) is constructed using the same paired arrays in the summary.

CPU gate 3851836, GPU pilot 3851837\_0, remaining GPU array
3851838\_1--9 and CPU summary 3851839 all completed with exit 0:0.
GPU times were approximately 30:42--32:42; the final CPU summary took
1:03. Frozen-source hashes and 40 NPZ primary-array headers passed
the recorded audit. Logs contain no numerical failures; locale and font
cache messages are not physics diagnostics.
To regenerate only the existing CPU summary, submit
\texttt{sbatch --parsable summarize\_cpu.slurm} from this experiment
directory; do not rerun GPU flow just for plotting.

Standard B C3 will be saved separately under
\path{C3_schemeB/<flow-tag>/Nev_200/<conf>/}, with
\texttt{flow\_weight\_scheme} and
\texttt{flow\_weight\_applied\_count=1} metadata.
Retain basis/gauge and final C2/C3; vertex/VVV can be rebuilt in job-local
scratch and removed there. GPU jobs use \texttt{na800-sug} and reservation
\texttt{donghx\_223}; analysis/fit jobs use CPU compute nodes, not login
nodes. No new physics job was submitted by this documentation update.

Next, trace the actual random generator/seed and full baryon moments;
project the paired window samples with the frozen physical spin
conventions; extend the test to other flows, currents and source/sink
separations before claiming a general truncation bound. Any downstream
ratio/fit/matching refresh needs a new B-labelled product and paired
configuration/bootstrap identities. The old A analysis and its
normalization/matching caveats are not overwritten.

\section{Wilson-loop relative masses, static potentials and candidate
line-mass coefficients: 20 October 2026}
\label{sec:line-mass-oct20}

\subsection{Scope, frozen inputs and current conclusion}

This update documents completed numerical products, not a new measurement
or a new fit performed during documentation. The ensemble is L32x64,
\(a=0.0897\,\mathrm{fm}\), with 1700 common configurations and the original
3000 configuration-level paired bootstrap replicas, seed 20261015. Each
replica draws 1700 configurations with replacement. Gauge measurements
start from unsmeared links and use four-dimensional Wilson RK3 flow with
\(\epsilon=0.01\). Wilson rectangles average their starting points and the
XT/YT/ZT planes; fundamental and adjoint traces are normalized separately.
These are pure-gauge measurements: no Blending A/B weight enters a loop.

The three relevant analysis roots below share the prefix
\begin{center}
\path{/public/group/lqcd/donghx/Gradient_Flow_GluonPDF/}.
\end{center}
\begin{itemize}
 \item \path{line_mass_d_boot3000_20261018/}: original paired loop averages
       and bootstrap indices.
 \item \path{wilson_relative_Tw_static_20261020/}: direct fixed-flow-pair
       Wilson-time fits and static-potential figures.
 \item \path{wilson_relative_d_from_Tw_20261020/}: subsequent flow fits,
       spatial plateaus, joint fits and new-reference scans.
\end{itemize}
The earlier \path{wilson_relative_radius030_20261020/} contains separate
raw-loop fits reused for the full static-energy plots. It is a different
finite-window estimator from the direct log-ratio slope.

\paragraph{Status, not a final subtraction prescription.}
Numerical completion and independent implementation checks have passed.
They do not establish a unique universal \(d_F\) or \(d_A\). The current
machine-readable status is \texttt{final\_d\_approved=false}, with both
\texttt{final\_d\_F} and \texttt{final\_d\_A} equal to \texttt{null}.
The finite-window estimates below must not be promoted to final coefficients
by selecting a small error or an attractive spatial plateau. No new
Polyakov fit or Wilson/Polyakov joint determination was performed in this
update, and no resulting factor was applied to a PDF.

\subsection{Conventions, flow window and paired statistics}

Let \(s=F,A\) denote the fundamental or adjoint representation, distinct
from the spatial distance \(R\). Preserve the positive-line-self-energy
convention
\begin{equation}
 h_s^{\rm GF}\sim e^{-m_{{\rm line},s}|z|},\qquad
 am_{{\rm line},s}(\tau)=d_s q(\tau)+\cdots,\qquad
 q(\tau)=\frac{1}{\sqrt{8\tau}},\qquad
 \delta m_s=-m_{{\rm line},s}.
 \label{eq:line-mass-oct20-convention}
\end{equation}
Here \(d_s\) is dimensionless and already includes representation
dependence. Do not multiply by another Casimir, average the two
representations, set \(m_A=2m_F\), or substitute Casimir scaling for an
unavailable adjoint measurement. Fundamental lines are relevant to quark
bilocals and adjoint lines to gluon bilocals; this does not by itself fix
their finite scheme conversion.

The user-selected strict radius window is
\begin{equation}
 a<r_f=\sqrt{8\tf}<0.3\,\mathrm{fm}
 \quad\Longleftrightarrow\quad
 0.125<\tau<1.3981946511.
 \label{eq:line-mass-oct20-radius}
\end{equation}
The actually measured checkpoints retained in this interval are
\begin{equation}
 \begin{split}
 \tau\in\{&0.15,0.20,0.25,0.30,0.35,0.40,0.45,\,0.50,\\
           &0.55,0.60,0.65,0.80,1.00\}.
 \end{split}
\end{equation}
The point 1.4 is outside the strict bound. The main reference is
\(\tau_r=0.15\); \(R/a=1,\ldots,12\); the two inclusive Wilson-time
windows are \(T_w/a=10,\ldots,16\) and \(8,\ldots,12\). Old additional
scale masks are not silently reapplied. Short-distance flow overlap is
recorded as a risk, not removed by the absence of a cut.

All flow times, distances, Wilson times and representations retain the
same configuration/bootstrap identities. Each stage constructs a fixed
full covariance from its 1700 delete-one jackknife estimators. Correlation
eigenmodes below \(10^{-8}\) of the largest eigenvalue are deleted; the
unregulated information-rank threshold is \(10^{-12}\). Rank and finite
input checks remain mandatory. There are no priors or
\(\chi^2/Q/\mathrm{AIC}\) quality cuts, covariance diagonalization,
negative-loop absolute values, or failed-replica replacements. Quality
diagnostics are nevertheless reported and interpreted.

A full statistical error is reported only for 3000/3000 usable bootstrap
replicas. It is the population standard deviation, \texttt{ddof=0}, not
that spread divided by \(\sqrt{3000}\). Incomplete cases retain their
center and usable count where defined, but no full bootstrap error. The
configuration-level iid resampling has not acquired a new autocorrelation
blocking study; the errors below are not total systematic uncertainties.

\subsection{Fixed-flow-pair finite Wilson-time fit}

For the central estimator, every JK sample, and every bootstrap sample,
first average each loop and then form
\begin{equation}
 Y_s(R,\ell;\tau,\tau_r)
 =-\log\frac{\overline W_s(R,\ell;\tau)}
                 {\overline W_s(R,\ell;\tau_r)},
 \qquad \ell=T_w/a.
 \label{eq:line-mass-oct20-logratio}
\end{equation}
This is not the average of per-configuration logarithms or ratios. Both
mean loops must individually be positive and finite. The positive ratio
of two negative means is not accepted as this real-log estimator.

If a finite time window admits
\(\overline W_s\simeq A_s\exp[-\ell aV_s^{\rm eff}]\), then
\begin{align}
 Y_s&=b_s+2\ell\,a\Delta m_s^{\rm eff},
 &b_s&=-\log\frac{A_s(R,\tau)}{A_s(R,\tau_r)},
 \label{eq:line-mass-oct20-time-fit}\\
 a\Delta m_s^{\rm eff}
 &=\frac{aV_s^{\rm eff}(R;\tau)-aV_s^{\rm eff}(R;\tau_r)}{2}.
 \label{eq:line-mass-oct20-factor-two}
\end{align}
The design matrix is \([1,2\ell]\) and parameter order is
\([b_s,a\Delta m_s^{\rm eff}]\). The free intercept absorbs
time-independent spatial-edge, endpoint and overlap-amplitude differences;
forcing it to vanish changes the estimator. The factor two accounts for
the two temporal long edges. It has already been removed from the stored
relative mass and must not be removed a second time in a flow fit.
In particular,
\begin{equation}
 \frac{Y_s}{2\ell}=a\Delta m_s^{\rm eff}+\frac{b_s}{2\ell}
 \label{eq:line-mass-oct20-finite-time}
\end{equation}
need not be a constant plateau at finite \(T_w\).

For the fundamental representation, \(R/a=4\), \(\tau=0.5\) relative to
0.15, both windows have 3000/3000 usable replicas:
\begin{center}
\small
\begin{tabular}{@{}cccc@{}}
 \toprule
 \(T_w/a\) & \(a\Delta m_F^{\rm eff}\) & \(b_F\) & \(\chi^2/\mathrm{dof}\) \\
 \midrule
 \(10..16\) & \(-0.145148\pm0.001636\) & \(-0.725736\pm0.032274\) & 0.246464 \\
 \(8..12\) & \(-0.143753\pm0.000486\) & \(-0.754096\pm0.007604\) & 0.492713 \\
 \bottomrule
\end{tabular}
\end{center}
The corresponding \(Q\) values are 0.941739 and 0.687325. Negative
relative masses here mean a lower effective energy at the larger flow
time, not an automatic sign error. Figure~\ref{fig:line-mass-oct20-time}
shows this example using the saved samples and fit parameters, without
repeating the fit. Of the 576 main time-fit candidates, 228 have finite
centers and 168 have complete bootstrap errors. Reference-loop signal
loss at large \(R\) propagates to multiple target flow times.

\begin{figure}[tbp]
 \centering
 \includegraphics[width=0.96\linewidth]{figures/wilson_direct_Tw_note_20261020.pdf}
 \caption{Direct paired log-ratio at fixed \(R/a=4\), \(\tau=0.5\),
 reference 0.15, fundamental representation. Blue and red denote the
 \(10..16\) and \(8..12\) Wilson-time windows, respectively, not the
 two flow models. Gray early-time points are context only. The fitted
 lines have a free intercept; the data are not joined.}
 \label{fig:line-mass-oct20-time}
\end{figure}

\subsection{Full static-energy plots and the Pion-DA presentation}

For full static energies, reuse the completed separate fits
\(\overline W=Ae^{-\ell aV_{\rm eff}}\) from
\path{wilson_relative_radius030_20261020/}; do not reinterpret the
relative-mass slope as the full potential. Convert units explicitly:
\begin{equation}
 z=R=na,\qquad
 V_{\rm eff}[\mathrm{GeV}]
 =(aV_{\rm eff})\frac{0.1973269804\,\mathrm{GeV\,fm}}
                          {0.0897\,\mathrm{fm}}.
 \label{eq:line-mass-oct20-static-units}
\end{equation}
These are unrenormalized finite-time effective static energies, with no
additive shift and no factor of one half. The half belongs only to the
relative line-mass difference in Eq.~\eqref{eq:line-mass-oct20-factor-two}.

The local presentation \path{Pion_DA.pdf}, page 15, uses \(z\,[\mathrm{fm}]\)
and \(V\,[\mathrm{GeV}]\), grouped by smoothing radius. Its two-state
model average on F32P30 is not the single-exponential finite-window
estimator on our L32x64 ensemble. Figures~\ref{fig:line-mass-oct20-static-F}
and \ref{fig:line-mass-oct20-static-A} adopt its axis/grouping style only,
not its spectral analysis. Filled points carry full bootstrap errors;
hollow points have a center but no full error; bottom crosses denote
undefined estimates, not zero potential. No absolute-value repair or
interpolation has filled failed distances.

\begin{figure}[tbp]
 \centering
 \includegraphics[page=1,width=0.96\linewidth]{figures/wilson_static_note_20261020.pdf}
 \caption{Fundamental finite-time static energy versus physical separation,
 \(T_w/a=10..16\), for five representative flows. Larger flow lowers
 much of the self-energy contribution, but shape changes can also include
 a flow-dependent interaction potential. The result is neither shifted
 nor divided by two. Both time windows and all 13 flows are available in
 the original figures listed below.}
 \label{fig:line-mass-oct20-static-F}
\end{figure}
\begin{figure}[tbp]
 \centering
 \includegraphics[page=2,width=0.96\linewidth]{figures/wilson_static_note_20261020.pdf}
 \caption{The corresponding adjoint effective static energy. Large-distance
 failures and incomplete replica coverage are explicit. Adjoint screening
 and finite flow/time effects preclude assuming that every distance obeys
 the same unscreened Cornell form or that a short-distance coefficient is
 already a pure line self-energy.}
 \label{fig:line-mass-oct20-static-A}
\end{figure}

\subsection{Flow dependence, spatial plateaus and the fundamental candidate}

At fixed \(R\), compare
\begin{align}
 \text{leading:}\quad a\Delta m_s^{\rm eff}
   &=d_s(q-q_r),\\
 \text{flow correction:}\quad a\Delta m_s^{\rm eff}
   &=d_s(q-q_r)+c_s(\tau-\tau_r).
 \label{eq:line-mass-oct20-flow-models}
\end{align}
There is no extra intercept: the relative definition cancels a
flow-independent constant. The reference's identically zero difference
is not an additional data point. An empirical
\(e_s(q^3-q_r^3)\) term, with or without \(c_s\), is also retained as a
single-spacing discretization-sensitivity diagnostic, not a validated
complete cutoff correction or a way to certify a fit by its \(Q\).

With fixed reference 0.15, upper bounds are 0.30, 0.50, 0.65, 0.80 and
1.00. The resulting \(d_s(R)\) values are fit to a constant using their
paired JK covariance across \(R\). Predeclared spatial windows are
\(1..2,2..3,3..4,4..5,4..6,5..7,6..7,6..8,6..12,7..12,8..12\).
An unavailable required distance is not silently omitted to create a
plateau. A separate full-covariance joint fit uses one common \(d_s\)
and, in the corrected model, a separate nuisance \(c_R\) for every
distance. The two routes and the two temporal windows are correlated
cross-checks, not independent measurements to average.

For \(T_w/a=8..12\), \(R/a=6,7\), reference 0.15 and the model with
\(c_s\), the spatial plateau gives:
\begin{center}
\small
\begin{tabular}{@{}cccc@{}}
 \toprule
 \(\tau_{\max}\) & Candidate \(d_F\) & Plateau \(\chi^2/\mathrm{dof}\) & Plateau \(Q\) \\
 \midrule
 0.50 & \(0.404186\pm0.002872\) & 1.94527 & 0.16310 \\
 0.65 & \(0.390556\pm0.001717\) & 1.27326 & 0.25916 \\
 0.80 & \(0.382635\pm0.001445\) & 0.000146 & 0.99038 \\
 1.00 & \(0.374752\pm0.001211\) & 1.78842 & 0.18112 \\
 \bottomrule
\end{tabular}
\end{center}
All these errors use 3000/3000 replicas and are conditional statistical
errors, not total uncertainties. A paired difference, retaining the common
bootstrap rather than adding two errors in quadrature, is
\begin{equation}
 d_F(\tau_{\max}=0.5)-d_F(\tau_{\max}=1.0)
 =0.02943358\pm0.00276048.
 \label{eq:line-mass-oct20-window-drift}
\end{equation}
Figure~\ref{fig:line-mass-oct20-d} shows this drift and model sensitivity.
The range is not automatically a completed systematic-error budget.

\paragraph{Why a reasonable spatial plateau is insufficient.}
At \(\tau_{\max}=1\), the input fixed-distance flow fits give
\begin{center}
\small
\begin{tabular}{@{}cccc@{}}
 \toprule
 \(R/a\) & Candidate \(d_F\) & Flow \(\chi^2/\mathrm{dof}\) & Flow \(Q\) \\
 \midrule
 6 & \(0.374137\pm0.001316\) & 62.06299 & \(1.55\times10^{-52}\) \\
 7 & \(0.371697\pm0.002685\) & 20.42724 & \(7.57\times10^{-17}\) \\
 \bottomrule
\end{tabular}
\end{center}
Thus a spatial-plateau \(Q\simeq0.18\) does not repair the poor upstream
flow model. The joint fit gives \(d_F=0.375269\pm0.001200\), but with
\(\chi^2/\mathrm{dof}=29.0264\) and \(Q=4.22\times10^{-51}\).
Neither is approved as the final coefficient. Rank identifiability,
model goodness of fit and physical validity are distinct questions.

\begin{figure}[tbp]
 \centering
 \includegraphics[page=1,width=0.96\linewidth]{figures/wilson_d_stability_note_20261020.pdf}
 \caption{Fundamental spatial-plateau candidates from \(R/a=6,7\) and
 \(T_w/a=8..12\), versus the flow-window upper bound. Blue: leading
 model; red: an added \(c(\tau-\tau_r)\); green: an empirical
 \(q^3-q_r^3\) sensitivity term. Small horizontal offsets aid comparison;
 no data are joined. The error bars are conditional bootstrap errors and
 do not cover the model/window drift.}
 \label{fig:line-mass-oct20-d}
\end{figure}

\subsection{New references and the adjoint candidate}

To test the poor small-reference signal, declare
\begin{equation}
 (\tau_r,\tau_{\max})\in
 \{(0.2,0.5),(0.2,0.65),(0.2,1.0),
    (0.3,0.65),(0.3,1.0),(0.5,1.0)\}.
\end{equation}
For each pair, rebuild Eq.~\eqref{eq:line-mass-oct20-logratio} from the
original mean-loop center/JK/bootstrap arrays and refit the time slope.
This is not subtraction of old slopes that used different correlated
weights. Resampling identities are unchanged; the new estimator is not
a replacement value for failed replicas of the old estimator.

For the adjoint representation with reference 0.5, maximum 1.0,
\(T_w/a=8..12\), and the model including \(c_s\), the nonreference
checkpoints are \(0.55,0.60,0.65,0.80,1.00\):
\begin{center}
\small
\begin{tabular}{@{}ccccc@{}}
 \toprule
 \(R/a\) & Candidate \(d_A\) & Flow \(\chi^2/\mathrm{dof}\) & \(Q\) & Usable/3000 \\
 \midrule
 4 & \(0.869908\pm0.002214\) & 1.46156 & 0.2228 & 3000 \\
 5 & \(0.833204\pm0.008895\) & 0.201074 & 0.8957 & 3000 \\
 6 & 0.840247; no full error & 3.93717 & 0.008058 & 2964 \\
 \bottomrule
\end{tabular}
\end{center}
The new reference recovers some signals, but at the largest flow
\(r_f/a=\sqrt8\); \(r_f/R\) is about 0.707 and 0.566 at \(R/a=4,5\).
Flow-dependent interaction contamination has not been eliminated, and
a controlled distance plateau is not established. A good \(Q\) at one
short distance is insufficient. Fundamental new-reference results are
also retained; Fig.~\ref{fig:line-mass-oct20-ref} illustrates their
sensitivity rather than selecting the smallest error.

\begin{figure}[tbp]
 \centering
 \includegraphics[page=2,width=0.96\linewidth]{figures/wilson_d_stability_note_20261020.pdf}
 \caption{Fundamental \(R/a=6,7\) plateau candidates with new direct
 references at fixed \(\tau_{\max}=1\), using the corrected flow model.
 Blue/red label the two Wilson-time windows; hollow markers have no
 complete bootstrap error. Agreement between selected points does not
 override their upstream flow-fit quality. No old slope has been
 subtracted to simulate a newly fitted direct ratio.}
 \label{fig:line-mass-oct20-ref}
\end{figure}

\subsection{What remains before a final coefficient can be used}

The measured relative slope still has the schematic content
\begin{equation}
 \Delta m_s^{\rm eff}(R;\tau,\tau_r)
 =\Delta m_{{\rm line},s}
  +\frac12\Delta V_{{\rm int},s}(R)
  +\text{finite-Wilson-time effects}.
 \label{eq:line-mass-oct20-contamination}
\end{equation}
Current fits have not completely separated these terms. All retained
measured flows satisfy \(\tau\le1\), hence \(a^2/\tf=1/\tau\ge1\),
not a parametrically small cutoff correction. The inequality \(r_f>a\)
is not proof of \(a^2/\tf\ll1\). Flow/model dependence, missing running
or higher-order terms, excited states, finite-spacing effects and finite
mass-scheme conversion remain explicit limitations. Canceling an additive
constant in relative differences does not fix an absolute residual-mass
prescription, endpoint conversion, \(Z_\chi\), operator mixing or LaMET.
An independently fitted nonperturbative factor must also not double count
the line term already present in an existing perturbative conversion.

\textbf{Pending, not executed:} authorize a separate analysis with a relaxed
0.3 fm upper-radius bound on already measured larger flows, or specify an
independently validated finite-\(a\) flow-self-energy model; then check
Wilson-time, distance, reference, covariance and flow-window stability,
autocorrelation blocking, and finite-scheme conversion. Relaxing a window
is not guaranteed to resolve competing scales. This documentation update
does not authorize that change, choose a final coefficient, or apply a
subtraction to a PDF.

\subsection{Code, samples, original figures and completed checks}

Under \path{wilson_relative_Tw_static_20261020/}, the primary products are
\path{analysis/direct_ratio_time_fits.csv},
\path{analysis/ratio_samples.npz} and
\path{analysis/static_potential_points.csv}. The code is
\path{code/tw_static.py}; the parameter order of time fits is
\([b,a\Delta m_{\rm eff}]\). The original figures are
\path{figures/01_fixed_flow_pair_Tw_fit_examples.pdf},
\path{figures/02_finite_Tw_relative_mass_examples.pdf},
\path{figures/03_all_fixed_flow_pairs_Tw_fits.pdf} (48 pages),
\path{figures/05_static_potential_PionDA_style.pdf}, and
\path{figures/06_static_potential_all_flow_times.pdf}.

Under \path{wilson_relative_d_from_Tw_20261020/}, use:
\begin{longtable}{@{}p{0.49\textwidth}p{0.46\textwidth}@{}}
 \toprule
 Relative path & Content \\
 \midrule
 \endhead
 \path{code/fit_d.py}, \path{code/audit.py} & Fixed-R flow fits, correlated
 spatial plateaus, joint checks and independent SVD audit.\\
 \path{code/reference_scan.py}, \path{code/final_review.py} & New direct
 reference fits, independent flow reconstruction and paired window drift.\\
 \path{analysis/flow_d_results.csv},
 \path{analysis/R_plateau_results.csv},
 \path{analysis/joint_R_flow_results.csv} & Windows, parameters, covariance
 ranks, fit quality, usable counts and failure states.\\
 \path{analysis/fit_samples.npz} & Suffixes \texttt{\_\_center},
 \texttt{\_\_jk}, \texttt{\_\_boot}; JK/bootstrap leading axes have
 lengths 1700/3000, with the reported parameter order on the last axis.\\
 \path{analysis/input_mass_ledger.csv},
 \path{analysis/paired_sensitivity.csv} & Traceable inputs and paired
 time-window/model changes.\\
 \path{analysis/final_d_status.json} & Final coefficients are null;
 candidates and paired drift are retained, with no automatic selection.\\
 \path{reference_scan/analysis/} & New-reference time, flow and spatial
 fits and their parameter samples.\\
 \path{figures/01_flow_fits_full_window.pdf},
 \path{figures/03_d_vs_R_window_stability.pdf} & All-R flow fits and
 distance/window sensitivity, including unavailable cases.\\
 \path{figures/04_R_plateau_fits.pdf},
 \path{figures/05_joint_d_window_scan.pdf} & Spatial plateaus and joint-fit
 window comparison.\\
 \path{figures/06_final_d_diagnostics.pdf} & Compact flow-window, model
 and new-reference comparison.\\
 \path{reference_scan/figures/01_reference_scan_flow_fits.pdf},
 \path{reference_scan/figures/02_reference_scan_d_vs_R.pdf} & Full
 new-reference flow fits and distance scans.\\
 \path{analysis/report.json}, \path{analysis/complete.json},
 \path{analysis/validation.json}, \path{analysis/independent_audit.json}
 & Source hashes, completion and implementation checks; analogous files
 are in the reference-scan analysis. PASS is not physical approval.\\
 \bottomrule
\end{longtable}

All numerical fitting and plotting used CPU partitions
\texttt{a192c1t}, \texttt{cpu6248R}, \texttt{cpueicc} or
\texttt{i72c512g}, not the login node. Main jobs 3860805, 3860811,
3860819 and 3860829 completed with exit 0:0 on cpu001. The initial review
job 3860824 failed on a same-name Python-module import and was corrected
without changing scientific fits. The preceding corrected time/static
job 3860685 and its audit 3860692 also completed. Independent SVD checks
cover 825 usable main fits and 316 new-reference flow fits; complete-error
checks cover 585 main fits and 895 new-reference stage fits. The main
audit retains 240 partial fits without full errors and 1455 unavailable
candidates. These checks establish implementation consistency, not a
controlled final line mass.

The main CPU entry is \path{slurm/run_cpu.slurm}, requiring
\texttt{TW\_D\_ROOT}; independent and reference jobs are
\path{slurm/audit_cpu.slurm}, \path{slurm/reference_cpu.slurm} and
\path{slurm/review_cpu.slurm}. The typical submission is
\texttt{sbatch --parsable --export=ALL,TW\_D\_ROOT=<analysis-root>}
followed by the relevant script, with explicit log paths. Completed
versions refuse in-place reruns: create a separate version for changed
windows or schemes. The Chinese handoff with full absolute links is
\path{/public/home/donghx/.codex/Markdown/divergence.md}, Section 26.
The document date is 20 October 2026; this section records only verified
existing results and pending decisions, not a new physics approval.

\clearpage
\section{Validation and systematic-error program}

A production calculation should pass:

\begin{enumerate}
  \item \textbf{Basis:} check \(V^\dagger V=1\), axes, color order, and
        \(\omega\) after compression.
  \item \textbf{Gauge flow:} monitor action, plaquette, unitarity, and the
        saved temporal phase; verify that the initial field is
        \texttt{Nosmear}.
  \item \textbf{Kernel:} reproduce \(e^{-\tf\widehat p^2}\) on unit links,
        including boundary conditions, and halve \(\epsilon\).
  \item \textbf{Gauge covariance:} transform a small field and verify \(K\),
        flowed links, and the final vertex.
  \item \textbf{Zero flow:} verify Eq.~\eqref{eq:zero-flow-identity} and compare
        all six Wick terms, \(C_2\), and \(C_3\) elementwise with the reference.
  \item \textbf{Temporal support:} monitor Eq.~\eqref{eq:discarded-norm} and
        compare correlators at successive \(w\).
  \item \textbf{Numerics:} vary precision, \(N_v\), batches, noise, and
        dilution.
  \item \textbf{Symmetries:} test \(\pm z\), \(\pm P_z\), and spatial
        rotations.
  \item \textbf{Spectrum:} test dispersion and compare plateau, summation,
        and multi-state fits while excluding contact regions.
  \item \textbf{Limits:} vary \(a\), \(\tf\), and \(P_z\); quantify volume,
        winding, \(z_{\max}\), perturbative, and power corrections.
\end{enumerate}

\section{Recommended analysis chain}

A complete chain is:
\begin{enumerate}
  \item fix the thin-link flowed operator, normalization, and matching scheme;
  \item flow links and all-time basis vectors at several \(\tf\), with
        step-size and temporal-support tests;
  \item build \(J^{(\tf)}\) and contract connected and required disconnected
        diagrams;
  \item apply \(\Gamma\) and fit \(C_2,C_3\) to extract \(\hflow(z,P_z)\);
  \item determine flowed-field normalization in the same scheme;
  \item take the continuum limit at fixed physical kinematics while testing
        volume and distillation effects;
  \item convert schemes and apply flowed LaMET matching
        \cite{MonahanOrginos2017,Monahan2018};
  \item apply mass and higher-twist corrections and reconstruct \(x\);
  \item vary \(P_z,\tf,a\), fits, and reconstruction priors for systematics.
\end{enumerate}
Matching coefficients must correspond to the implemented operator and
normalization \cite{Xiong2014,Izubuchi2018,Constantinou2021}.

\section{Conclusions}

Four-dimensional flow mixes time slices on both sides of the bilocal current.
Flowing all-time LapH vectors moves this mixing into a two-source-time vertex,
so existing all-time light perambulators need no new operator-only inversion.
Time-diagonal strange data support 3D flow, but complete 4D strange flow needs
off-diagonal perambulators.

Three-dimensional flow exponentially suppresses spatial momentum and, at one
loop, makes the straight spatial Wilson-line self-energy finite at fixed flow
time, replacing the cutoff linear divergence by a
\(|z|/\sqrt{\tf}\) term.  It does not damp arbitrarily large Euclidean energy,
so this result is not a proof that the complete bilocal and its normalization
are UV finite to all orders.  A 3D calculation therefore needs its own
renormalization theorem, ringed-field prescription, operator-mixing analysis,
and matching coefficient.  Existing gradient-flow quasi-PDF matching assumes
the standard 4D kernel and cannot be transferred by identifying the numerical
value of \(\tf\).

The \(32^3\times96\) quark outputs are connected finite-flow-time correlators,
while the gluon products are vacuum-subtracted disconnected finite-flow-time
ratios.  Neither is a PDF.  The gluon calculation uses a frozen common set of
406 configurations, three positive \(z\)-momenta, 14 flow times, and saves
the combined, \(M_{ti;ti}\), and \(M_{ij;ij}\) components for both
unpolarized and helicity channels.  The unpolarized combined selector is
consistent with the Euclidean continuation.  The helicity raw components are
also correct, but its legacy plus-sign combined ratio is not: the common
Euclidean convention requires \(M_{ti;ti}-M_{ij;ij}\) in both channels, with
an odd-separation projection for helicity.  Existing raw OPE files permit this
correction without repeating the flow.  The new C3/C2 product reconstructs, in
order, \(T\), \(T+S\), and \(T-S\) for each channel and resamples each of the
six choices to Nboot=1000.  Physical results still require state isolation,
normalization and mixing, continuum and flow-time control, matching, and
finite-\(z\) reconstruction.  No HYP pre-smearing is used.
Temporal truncation, distillation, and the limited scale hierarchy remain
explicit systematics.

\appendix

\section{Tensor dictionary for the implementation}

With spin indices shown explicitly, the light perambulator is represented as
\begin{equation}
  \tau_f(t,\tau)_{\alpha\beta;mn}
  =
  \left[V^\dagger(t)D_f^{-1}(t,\tau)V(\tau)\right]_{\alpha\beta;mn}.
\end{equation}
The sequential orientation follows from \(\gamma_5\)-hermiticity with an
axis-aware transpose, not unchecked elementwise conjugation.

For the production defaults, the principal arrays are
\begin{center}
\begin{tabular}{p{0.28\textwidth}p{0.25\textwidth}p{0.39\textwidth}}
  \toprule
  Object & Shape & Axes \\
  \midrule
  Flowed basis slice
  & \((400,32,32,32,3)\)
  & \((n,z,y,x,c)\) \\
  Generalized vertex
  & \((2,21,400,400)\)
  & \((s_z,|z|,n_{\rm left},n_{\rm right})\) \\
  \texttt{corr\_2pt}
  & \((4,3,96,96)\)
  & \((\Pi,P,t_{\rm src},t_{\rm sink})\) \\
  \texttt{corr\_3pt\_u/d}
  & \((2,4,3,21,4,4,96,96)\)
  & \((s_z,\Pi,P,|z|,m,o,t_{\rm src},t_J)\) \\
  \bottomrule
\end{tabular}
\end{center}
The three momentum entries are \(P_z=+3,+4,+5\).  Files retain full
\(96\times96\) time axes, with only allowed entries populated.

\section{Code sources and assumptions}

The principal files inspected through September 2026 are:
\begin{itemize}
  \raggedright
  \item \nolinkurl{/public/group/lqcd/donghx/Gradient_Flow/Blending_L32x96/flow_common_l32x96.py};
  \item \nolinkurl{/public/group/lqcd/donghx/Gradient_Flow/Blending_L32x96/generate_flowed_vertex_l32x96.py};
  \item \nolinkurl{/public/group/lqcd/donghx/Gradient_Flow/Blending_L32x96/calc_flowed_c2c3_l32x96.py};
  \item \nolinkurl{/public/group/lqcd/donghx/Gradient_Flow/Blending_L32x96/calc_flowed_c2c3_factorized_l32x96.py};
  \item \nolinkurl{/public/group/lqcd/donghx/Gradient_Flow/Blending_L32x96/validate_flow_l32x96.py};
  \item \nolinkurl{/public/group/lqcd/donghx/Strange_PDF/Conneted_Diagram_final/L32x96/Calc_L32x96.py};
  \item \nolinkurl{/public/group/lqcd/donghx/Strange_PDF/Strange_PDF_sink_distillation/Calc_ga_gpu_L24x72_VDzU0_hyp_new.py};
  \item \nolinkurl{/public/home/sush/EasyDistillation/work/eigenvectors.one_bin.py};
  \item \nolinkurl{/public/home/sush/EasyDistillation/work/perambulator.beta6.41_mu-0.2295_ms-0.2050_L32x96.py};
  \item \nolinkurl{/public/home/donghx/PyQuda_a100_2608/PyQUDA-gflow-src/pyquda/src/quda.in.pyx};
  \item \nolinkurl{/public/home/donghx/PyQuda_a100_2608/quda/lib/interface_quda.cpp};
  \item \nolinkurl{/public/home/donghx/PyQuda_a100_2608/quda/lib/gauge_wilson_flow.cu};
  \item \nolinkurl{/public/home/donghx/Gradient_Flow/Blending_L32x96/generate_flowed_vertex_l32x96.py}.
  \item \nolinkurl{/public/group/lqcd/donghx/Gradient_Flow_GluonPDF/quark_l32x64_flow_20261007/reference/flow_common_l32x64.py};
  \item \nolinkurl{/public/group/lqcd/donghx/Gradient_Flow_GluonPDF/quark_l32x64_flow_20261007/reference/generate_flowed_vertex_l32x64.py};
  \item \nolinkurl{/public/group/lqcd/donghx/Gradient_Flow_GluonPDF/quark_l32x64_flow_20261007/full_chain.py};
  \item \nolinkurl{/public/group/lqcd/donghx/Gradient_Flow_GluonPDF/src/flowed_gluon_ope.cpp};
  \item \nolinkurl{/public/home/donghx/.codex/Markdown/gradient_flow_gluon.md};
  \item \nolinkurl{/public/group/lqcd/donghx/Gradient_Flow_GluonPDF/handoff/Gradient_flow_gluon/calc_ratio_gradient_flow.py};
  \item \nolinkurl{/public/group/lqcd/donghx/Gradient_Flow_GluonPDF/ratio_bootstrap_all_operators_v1/construct_bootstrap.py};
  \item \nolinkurl{/public/group/lqcd/donghx/Gradient_Flow_GluonPDF/ratio_bootstrap_all_operators_v1/validate_product.py}.
\end{itemize}
The analysis assumes that light perambulators cover all required times and
distillation indices. Each run must verify orthonormality, metadata, gauge
layout, boundary phases, and axis order. The strange conclusion assumes the
active \texttt{tXXX\_tsinkXXX\_gen} dataset; legacy full-time files require
the audit in Section~\ref{sec:3d-vs-4d}.

\begingroup
\small
\bibliographystyle{unsrt}
\bibliography{reference}
\endgroup

\end{document}
